
% BEGIN REVISED SOURCE main.tex SHA256=9e9afa86d52a621f8b2616b126e6b62c374116fb82a3953b8ade31769d577408
\documentclass[11pt,oneside,a4paper]{book}

% BEGIN REVISED SOURCE preamble.tex SHA256=875c7f480b0458122d3ee7ea1b603458d21d5d30104027b8acbf5d326370e741
\usepackage[T1]{fontenc}
\usepackage{lmodern}
\usepackage{amsmath,amssymb,amsthm,mathtools}
\usepackage[a4paper,margin=25mm,headheight=15pt]{geometry}
\usepackage{microtype,enumitem,xcolor,longtable,booktabs,array,needspace}
\usepackage[colorlinks=true,linkcolor=blue!40!black,urlcolor=blue!40!black]{hyperref}
\hypersetup{pdftitle={Single-Variable Analysis -- Unified Undergraduate Notes},pdfauthor={MathGenome mathematical preparation series}}
\usepackage{fancyhdr}
\pagestyle{fancy}\fancyhf{}
\fancyhead[L]{\small MathGenome: Single-Variable Analysis}
\fancyhead[R]{\small\thepage}
\renewcommand{\headrulewidth}{0.3pt}
\setlength{\parindent}{0pt}
\setlength{\parskip}{5pt}
\setlength{\emergencystretch}{3em}
\allowdisplaybreaks[2]\raggedbottom
\setcounter{tocdepth}{1}
\newtheorem{theorem}{Theorem}[chapter]
\newtheorem{lemma}[theorem]{Lemma}
\newtheorem{proposition}[theorem]{Proposition}
\newtheorem{corollary}[theorem]{Corollary}
\theoremstyle{definition}
\newtheorem{definition}[theorem]{Definition}
\newtheorem{example}[theorem]{Example}
\newtheorem{construction}[theorem]{Construction}
\theoremstyle{remark}
\newtheorem{remark}[theorem]{Remark}
\newcommand{\N}{\mathbb N}
\newcommand{\Z}{\mathbb Z}
\newcommand{\Q}{\mathbb Q}
\newcommand{\R}{\mathbb R}
\newcommand{\C}{\mathbb C}
\newcommand{\K}{\mathbb K}
\newcommand{\eps}{\varepsilon}
\newcommand{\abs}[1]{\lvert#1\rvert}
\newcommand{\norm}[1]{\lVert#1\rVert}
\newcommand{\set}[1]{\{#1\}}
\newcommand{\ip}[2]{\langle#1,#2\rangle}
\newcommand{\dd}{\,\mathrm d}
\DeclareMathOperator{\Var}{Var}
\DeclareMathOperator{\osc}{osc}
\DeclareMathOperator{\supp}{supp}
\DeclareMathOperator{\dist}{dist}
\DeclareMathOperator{\diam}{diam}
\DeclareMathOperator{\RePart}{Re}
\DeclareMathOperator{\ImPart}{Im}
\newcommand{\KEntry}[2]{\Needspace{7\baselineskip}\section{#2}\label{#1}{\small\sffamily\color{black!60}\detokenize{#1}}\par}
\newcommand{\KRef}[1]{\hyperref[#1]{\S\,\ref*{#1}}}
\newcommand{\SourceRef}[1]{{\small\textsf{Source comparison: #1.}}}
\begin{document}
\frontmatter
\title{Single-Variable Analysis\\Unified Undergraduate Notes}
\author{MathGenome mathematical preparation series}
\date{31 October 2026}
\hypersetup{pageanchor=false}
\maketitle
\clearpage
\pagenumbering{roman}
\hypersetup{pageanchor=true}

% BEGIN REVISED SOURCE frontmatter.tex SHA256=5bf2ac5272778e10a6e6c3f58b503bde5c5d562600314f970b3529ff44845312
\chapter*{Purpose, scope and reading conventions}
\addcontentsline{toc}{chapter}{Purpose, scope and reading conventions}
These notes organize a single-variable undergraduate analysis course by mathematical
content rather than by copying the chapter order of one textbook. The primary
comparison sources are Zorich I and II, Rudin, Apostol and Pugh. Tao's first
analysis volume supplies a second foundational viewpoint; Olver is used only
for bounded real-variable asymptotic supplements. The source register explains
the selected scope and the original-source locations.

The stable identifiers \texttt{K03.xx.name} refer to the accompanying knowledge
registry. They identify mathematical content, not Lean declarations or a prescribed
number of formalization tasks. A definition, a counterexample, a representation
comparison and a theorem are different kinds of content. The main text retains
the fourteen original core modules and five explicit supplements. Five records
only route material to another direction: Riemann--Lebesgue and Stieltjes--measure
interfaces, advanced asymptotics, general special functions, and general Fourier
transform theory. Those routing records assert no unproved theorem and are not
prerequisites of the proved D03 content.

We use the standard complete ordered real field and its complex extension, along
with elementary set theory and finite arithmetic. These are declared shared
foundations, not newly proposed axioms for a Lean project. Consequences used by
analysis, including completeness arguments and convergence criteria, are proved
in the text. A separate construction of the reals by cuts or rational Cauchy
sequences is not required to use this interface. We do not assume Lebesgue
integration, general multidimensional change of variables, complex contour
integration or a general ODE existence theorem in order to prove one-variable
results.

Unless stated otherwise, an interval is an interval in $\R$, a scalar belongs to
$\R$ or $\C$, and derivatives use real increments. An exact mean-value equality
with one intermediate point is a real-valued statement; complex-valued functions
have a norm estimate instead. The domains of logarithms, real powers and inverse
functions, endpoint continuity, integrability and directions of limits are part
of each assertion. A total integral expression in a future formal library never
replaces an integrability hypothesis.

The elementary functions are built from general series and power-series results;
the logarithm is first the inverse of the exponential. Its integral formula is
then a comparison theorem. For this reason the exposition places general series
before the elementary-function chapter. A few applications refer forward to
later established functions or integration rules; these applications are not
used to prove the foundational series or completeness results. The companion
audit records the prerequisite choices and actual review limits.

The mathematical prose and proofs are newly written syntheses. Source references
identify comparison material; they do not mean that every theorem or exercise in
every book has been included. Distinct useful hypotheses are kept distinct, and
the generality of a comparison bridge is not imposed on the underlying theory.
In particular, refinement and unrestricted mesh versions of a Stieltjes integral
are genuinely different for some jump integrators.
\tableofcontents
\mainmatter

% BEGIN REVISED SOURCE foundations.tex SHA256=2d0bb628eec467164ea16232a77ee6c60e1b8488b423a462c96d79441ebad68c
\chapter*{Declared shared foundations}
\addcontentsline{toc}{chapter}{Declared shared foundations}
\label{shared-foundations}
The set-theoretic language consists of functions, subsets, finite and countable
families, the usual natural numbers with induction, and elementary finite sums
and products. We use countable choices when choosing a sequence of witnesses
from nonempty sets. Finite binomial identities and factorial arithmetic are part
of elementary algebra and can also be established by induction. For example,
$(1+t)^n=\sum_{k=0}^n\binom nk t^k$ follows by multiplying the $n$th identity by
$(1+t)$ and using Pascal's relation. No convergence assertion is encoded in this
finite identity.

The scalar interface is a complete ordered field $\R$: every nonempty real set
bounded above has a supremum. Its complex extension is $\C=\R+i\R$ with
$\abs{x+iy}=\sqrt{x^2+y^2}$. Real inequalities concern real quantities. The
triangle inequality for complex scalars follows by squaring and using the finite
real Cauchy--Schwarz inequality. That inequality itself follows from
$\sum_j(u_j-tv_j)^2\geq0$: minimizing the quadratic in $t$ gives
$(\sum_j u_jv_j)^2\leq(\sum_j u_j^2)(\sum_jv_j^2)$, with the zero-denominator
case immediate. Complex finite Cauchy--Schwarz follows by choosing the phase of
the complex inner product and applying the same real estimate, or by expanding
$\sum_j\abs{u_j-tv_j}^2\geq0$ with complex $t$.

A metric has nonnegative symmetric distances, distance zero exactly at equal
points, and the triangle inequality. Open and closed sets, metric completeness
and compactness by finite subcovers have their usual definitions; the compactness
and completeness results actually needed below are proved. Compactness does not
mean mere boundedness in an arbitrary metric space. For the rare finite-dimensional
algebraic application we may reuse the accepted field arithmetic and polynomial
factorization interfaces of the fixed mathematical library. This does not import
a theorem about infinite sums or analytic integrals.

All other results invoked by a stable D03 identifier are proved or constructed
within the notes. In particular a Gaussian integral is not justified by an
unprepared polar-coordinate theorem: the text uses Wallis integrals and a
one-variable Gamma limit. Likewise uniform convergence of proper truncations
and explicit tail bounds replace an implicit appeal to a Lebesgue dominated
convergence or Fubini theorem.

% BEGIN REVISED SOURCE chapters/M03.01.tex SHA256=a5792520eebd1bf866ce3b7a1b53fb08ca4afc9fa3491eb12bace5f2c900d3ac
\chapter{Sequences and completeness}
Throughout this chapter indices start at $1$, and $\K$ denotes $\R$ or $\C$. Statements involving order apply to real numbers. The proofs use completeness at identifiable points; no logarithm or exponential is needed for the sequence tests.
\SourceRef{Zorich I, 3.1; Rudin, Chapter 3; Tao, Chapter 6}

\KEntry{K03.01.real_interface}{The scalar foundation}
\begin{definition}
We use one complete ordered field $\R$: every nonempty subset bounded above has a least upper bound. We use $\C=\R^2$ with its usual field operations, conjugation, and modulus $|u+iv|=(u^2+v^2)^{1/2}$. The field laws, order laws, finite sums and products, natural-number induction, elementary set theory, and metric axioms are shared foundations. In particular the triangle inequality and $|zw|=|z||w|$ are available scalar algebra. Completeness of $\C$ as a metric space will be proved below, rather than separately postulated.
\end{definition}
The nonnegative square root in this interface agrees with the completeness construction in \KRef{K03.01.positive_roots}; uniqueness there makes the notations compatible. Constructing the real field itself belongs to the preceding direction of study.

\KEntry{K03.01.archimedean}{Archimedean bounds and rational density}
\begin{theorem}
For every $x\in\R$ some $n\in\N$ satisfies $n>x$. If $x<y$, some $q\in\Q$ satisfies $x<q<y$.
\end{theorem}
\begin{proof}
If $\N$ had a supremum $s$, then $s-1$ would not be an upper bound. Choose $n>s-1$; then $n+1>s$, a contradiction. To obtain density choose $n$ with $n(y-x)>1$. There is an integer $m$ with $nx<m\le nx+1$: indeed choose an integer $N>|nx|+1$ and take the least member of the nonempty finite set of integers in $[-N,N]$ greater than $nx$. Its predecessor is at most $nx$. Consequently $x<m/n\le x+1/n<y$.
\end{proof}
In particular $1/n\to0$ directly from its defining inequality, and rational approximations with arbitrarily small prescribed errors exist.

\KEntry{K03.01.sup_inf}{Suprema and infima}
\begin{definition}
An upper bound of $A\subset\R$ is a number $b$ with $a\le b$ for every $a\in A$. Its least upper bound is $\sup A$. The complete-field foundation asserts its existence when $A$ is nonempty and bounded above. Define $\inf A=-\sup(-A)$ for nonempty sets bounded below. In $\overline\R=\R\cup\{-\infty,+\infty\}$ we set $\sup\varnothing=-\infty$, $\inf\varnothing=+\infty$, and assign infinite bounds to unbounded sets. These conventions do not assert that an empty set has a real supremum.
\end{definition}
\begin{proposition}
If $s=\sup A\in\R$ and $\eps>0$, some $a\in A$ satisfies $s-\eps<a\le s$. If $t=\inf A\in\R$, some $a\in A$ satisfies $t\le a<t+\eps$.
\end{proposition}
\begin{proof}
Otherwise $s-\eps$ would be a smaller upper bound. Apply this conclusion to $-A$ for the infimum assertion.
\end{proof}

\KEntry{K03.01.positive_roots}{Positive integer roots}
\begin{construction}
For $a\ge0$ and integer $m\ge1$ there is exactly one $r\ge0$ with $r^m=a$; write $r=a^{1/m}$.
\end{construction}
\begin{proof}
For $a=0$ take $r=0$. For $a>0$, let $S=\{t\ge0:t^m\le a\}$ and $r=\sup S$. The set contains a positive number, for $0<t<\min(1,a)$ gives $t^m\le t<a$, and is bounded above by $a+1$. Thus $r>0$. For $u,v\in[0,r+1]$ the factorization of $u^m-v^m$ gives
\[ |u^m-v^m|\le m(r+1)^{m-1}|u-v|. \]
If $r^m<a$, choose $0<h<1$ with $m(r+1)^{m-1}h<a-r^m$. Then $(r+h)^m<a$, contradicting the upper-bound property of $r$. If $r^m>a$, choose $0<h<\min(1,r)$ with $m(r+1)^{m-1}h<r^m-a$. Then $(r-h)^m>a$. Since $t\mapsto t^m$ is strictly increasing on $[0,\infty)$ by its factorization, $r-h$ is an upper bound of $S$, again a contradiction. The same strict increase proves uniqueness. For a positive base $a$, rational powers $a^{p/q}=(a^{1/q})^p$, with integer $p$ and positive integer $q$, use these roots, and negative $p$ uses the reciprocal. Their exponent laws and independence of the fraction representation follow by raising to common positive integer powers and using uniqueness. At base zero, positive rational exponents extend with value zero; no negative-power law is asserted there.
\end{proof}

\KEntry{K03.01.sequence_limit}{Convergence of a sequence}
\begin{definition}
A sequence $a:\N\to\K$ converges to $L\in\K$ if
\[ (\forall\eps>0)(\exists N\in\N)(\forall n\ge N)\quad |a_n-L|<\eps. \]
Write $a_n\to L$. For real sequences $a_n\to+\infty$ means that for every real $B$ eventually $a_n>B$; $-\infty$ is defined with the reverse inequality. These are not finite scalar limits. An assertion holding for all $n\ge N$ for some $N$ holds \emph{eventually}.
\end{definition}

\KEntry{K03.01.unique_limit}{Uniqueness, boundedness, and finite prefixes}
\begin{proposition}
A convergent scalar sequence has exactly one limit and is bounded. Changing or deleting finitely many terms does not change its limit.
\end{proposition}
\begin{proof}
If $L\ne M$ were limits, take $\eps=|L-M|/3$ and a common tail. There $|L-M|\le|L-a_n|+|a_n-M|<2|L-M|/3$, impossible. On a tail with $|a_n-L|<1$, $|a_n|\le |L|+1$; taking the maximum with the finitely many earlier moduli proves boundedness. Each limit condition concerns only a tail, so finitely many changes are irrelevant.
\end{proof}

\KEntry{K03.01.limit_algebra}{Algebra and order of limits}
\begin{theorem}
If $a_n\to a$, $b_n\to b$, and $c\in\K$, then $a_n+b_n\to a+b$, $ca_n\to ca$, $a_nb_n\to ab$, and $|a_n|\to|a|$. If $b\ne0$, eventually $b_n\ne0$ and $a_n/b_n\to a/b$. For real sequences, eventual $a_n\le b_n$ implies $a\le b$.
\end{theorem}
\begin{proof}
The estimates for sums and scaling are $|(a_n+b_n)-(a+b)|\le|a_n-a|+|b_n-b|$ and $|ca_n-ca|=|c||a_n-a|$. If $|a_n|\le C$, then $|a_nb_n-ab|\le C|b_n-b|+|b||a_n-a|$. The reverse triangle inequality gives $||a_n|-|a||\le|a_n-a|$. If $b\ne0$, eventually $|b_n|\ge |b|/2$, and
\[ \left|\frac1{b_n}-\frac1b\right|\le\frac{2|b_n-b|}{|b|^2}. \]
These estimates tend to zero and give the quotient by multiplication. Finally if $a>b$, tails with $|a_n-a|,|b_n-b|<(a-b)/3$ contradict $a_n\le b_n$. Strict inequalities need not survive: $0<1/n$ but both limits are zero.
\end{proof}

\KEntry{K03.01.squeeze}{Squeeze and absolute estimates}
\begin{theorem}
If real $a_n\le b_n\le c_n$ eventually and $a_n,c_n\to L$, then $b_n\to L$. If $|z_n-L|\le d_n$ eventually and $d_n\to0$, then $z_n\to L$.
\end{theorem}
\begin{proof}
On a common sufficiently late tail $L-\eps<a_n\le b_n\le c_n<L+\eps$. The second assertion follows directly from $d_n<\eps$ on a tail.
\end{proof}

\KEntry{K03.01.monotone}{Monotone convergence}
\begin{theorem}
A nondecreasing real sequence bounded above converges to the supremum of its range. If unbounded above it tends to $+\infty$. A nonincreasing sequence has the dual conclusions with the infimum and $-\infty$.
\end{theorem}
\begin{proof}
If $s$ is the supremum, choose $N$ with $s-\eps<a_N\le s$. For $n\ge N$, $s-\eps<a_n\le s$. If unbounded above, choose $a_N>B$ and use monotonicity for every later term. Negation of the sequence proves the decreasing assertions.
\end{proof}

\KEntry{K03.01.nested_intervals}{Nested closed intervals}
\begin{theorem}
For nonempty closed bounded intervals $I_n=[a_n,b_n]$ with $I_{n+1}\subset I_n$, the intersection is nonempty. If $b_n-a_n\to0$, it contains exactly one point.
\end{theorem}
\begin{proof}
Let $s=\sup_n a_n$, which exists since $a_n\le b_1$. For each fixed $j$, all $a_n\le b_j$: this follows from nesting for $n\ge j$, and from $a_n\le a_j$ for $n<j$. Thus $a_j\le s\le b_j$. If $s,t$ both lie in every interval, $|s-t|\le b_n-a_n\to0$, so $s=t$. Without shrinking lengths, $I_n=[0,1]$ has a whole interval as intersection.
\end{proof}

\KEntry{K03.01.subsequence}{Subsequences}
\begin{definition}
A subsequence has the form $a_{n_j}$ where $n_1<n_2<\cdots$ are positive integers. A limit of a convergent subsequence is a subsequential limit of $a_n$.
\end{definition}
\begin{proposition}
If $a_n\to L$, then every subsequence tends to $L$.
\end{proposition}
\begin{proof}
Induction gives $n_j\ge j$. A tail estimate for $a_n$ therefore applies whenever $j$ is at least its threshold.
\end{proof}

\KEntry{K03.01.bolzano}{Bolzano--Weierstrass}
\begin{theorem}
Every bounded real or complex sequence has a convergent subsequence.
\end{theorem}
\begin{proof}
Enclose a real sequence in a closed bounded interval. Bisect the interval and retain a closed half containing infinitely many terms; at least one half does, since two finite collections cannot contain infinitely many indices. Repeat. The retained intervals are nested with lengths tending to zero, so have a unique common point $L$. Choose recursively $n_j>n_{j-1}$ with $a_{n_j}$ in the $j$th interval. Its distance from $L$ is at most that interval's length. For a complex sequence first choose a subsequence whose real parts converge, and then a subsequence of that subsequence whose imaginary parts converge. The real-part limit persists, and $|u+iv|\le |u|+|v|$ proves convergence to the pair of limits.
\end{proof}

\KEntry{K03.01.cauchy}{The Cauchy criterion}
\begin{theorem}
A scalar sequence converges in $\R$ or $\C$ if and only if it is Cauchy, meaning that for every $\eps>0$ eventually $|a_n-a_m|<\eps$ for all $m,n$ on the same tail.
\end{theorem}
\begin{proof}
Convergence implies the assertion by the triangle inequality with two errors $<\eps/2$. A Cauchy sequence has a bounded tail by fixing one term after the threshold for $\eps=1$, and hence is bounded. Choose a subsequence tending to $L$ by \KRef{K03.01.bolzano}. Given $\eps>0$, choose the Cauchy threshold for $\eps/2$, and a subsequence term beyond it within $\eps/2$ of $L$. Every term after the threshold is then within $\eps$ of $L$. This also proves metric completeness of $\C$.
\end{proof}
The rational version fails: rational numbers within $1/n$ of $\sqrt2$ form a Cauchy sequence, but their unique real limit is irrational. Irrationality follows because $p^2=2q^2$ for relatively prime integers would force both $p$ and $q$ even.

\KEntry{K03.01.limsup}{Upper and lower limits}
\begin{definition}
For a real sequence put $U_N=\sup_{n\ge N}a_n$ and $L_N=\inf_{n\ge N}a_n$ in $\overline\R$. Define
\[ \limsup a_n=\inf_N U_N,\qquad \liminf a_n=\sup_N L_N. \]
For bounded sequences these are real numbers, $U_N$ decreases to the former, and $L_N$ increases to the latter.
\end{definition}
\begin{proposition}
For a bounded real sequence, $\limsup a_n$ and $\liminf a_n$ are subsequential limits, and every subsequential limit lies between them.
\end{proposition}
\begin{proof}
Write $u=\lim U_N$. Recursively choose $n_j\ge\max(j,n_{j-1}+1)$ with $a_{n_j}>U_{\max(j,n_{j-1}+1)}-1/j\ge u-1/j$, using supremum approximation. Also $a_{n_j}\le U_j\to u$, so squeeze gives the desired subsequence. Apply this to $-a_n$ for the infimum. For any subsequential limit $v$, fixing $N$ and passing to the limit in $L_N\le a_{n_j}\le U_N$ gives $L_N\le v\le U_N$; then let $N\to\infty$.
\end{proof}

\KEntry{K03.01.limsup_rules}{Limsup rules and convergence}
\begin{theorem}
For a bounded real sequence, $a_n\to L$ if and only if $\liminf a_n=\limsup a_n=L$. For bounded real sequences $a,b$,
\[ \liminf a+\liminf b\le\liminf(a+b)\le\limsup(a+b)\le\limsup a+\limsup b. \]
Moreover $\limsup(-a)=-\liminf a$.
\end{theorem}
\begin{proof}
Convergence places all sufficiently late tail bounds in $[L-\eps,L+\eps]$. Conversely equality of the two limits traps $a_n$ between tails arbitrarily close to $L$. On every tail the infimum of $a+b$ is at least the sum of the two infima, and the supremum is at most the sum of the two suprema; passage to limits proves the displayed bounds. Negation interchanges upper and lower bounds. Equalities in the sum bounds need not hold: $a_n=(-1)^n$, $b_n=-a_n$ give $\limsup(a+b)=0<2=\limsup a+\limsup b$.
\end{proof}

\KEntry{K03.01.cesaro}{Arithmetic averages}
\begin{theorem}
If $a_n\to L$ in $\K$, then $n^{-1}\sum_{j=1}^n a_j\to L$.
\end{theorem}
\begin{proof}
For $j\ge N$ arrange $|a_j-L|<\eps/2$. Then for $n\ge N$ the modulus of the averaged error is at most
\[ \frac1n\sum_{j<N}|a_j-L|+\frac{n-N+1}{n}\frac\eps2. \]
The first term tends to zero. The converse fails for $a_n=(-1)^n$: its partial sums have modulus at most $1$, so its averages tend to zero, although its even and odd subsequences have distinct limits.
\end{proof}

\KEntry{K03.01.stolz}{Stolz--Cesaro with unbounded denominator}
\begin{theorem}
Let $b_n$ be real, strictly increasing, and tend to $+\infty$. If $a_n\in\K$ and
\[ \frac{a_{n+1}-a_n}{b_{n+1}-b_n}\longrightarrow L\in\K, \]
then $a_n/b_n\to L$.
\end{theorem}
\begin{proof}
Set $c_n=a_n-Lb_n$. Given $\eps>0$, choose $N$ so that $|c_{j+1}-c_j|\le\eps(b_{j+1}-b_j)$ for $j\ge N$. Telescoping yields $|c_n|\le|c_N|+\eps(b_n-b_N)$ for $n>N$. Eventually $b_n>0$, and division shows $\limsup |c_n/b_n|\le\eps$. Since $\eps$ is arbitrary, the quotient tends to zero. Strict increase supplied the positive weights; no decreasing-denominator variant has been asserted.
\end{proof}

\KEntry{K03.01.standard_limits}{Reference sequences without logarithms}
\begin{example}
If $|q|<1$, then $q^n\to0$. Also $n^{1/n}\to1$, and $n^k/r^n\to0$ for fixed integer $k\ge0$ and $r>1$.
\end{example}
\begin{proof}
For $0<|q|<1$ write $1/|q|=1+d$, $d>0$. Induction gives $(1+d)^n\ge1+nd$, so $|q|^n\le(1+nd)^{-1}\to0$; $q=0$ is immediate. For each $\eps>0$ the binomial formula gives $(1+\eps)^n\ge\binom n2\eps^2>n$ for sufficiently large $n$. Root monotonicity gives $1\le n^{1/n}<1+\eps$. Finally write $r=1+d$ and take $m=k+1$. For $n\ge2m$,
\[ r^n\ge\binom nm d^m\ge \frac{(n/2)^m d^m}{m!},\qquad 0\le\frac{n^k}{r^n}\le\frac{2^m m!}{d^m n}. \]
Every bound used here is algebraic.
\end{proof}

\KEntry{K03.01.e_sequence}{The classical sequence for the exponential constant}
\begin{construction}
The sequence $u_n=(1+1/n)^n$ converges to a finite positive number. After the power-series construction of $\exp$, this limit equals $\exp(1)$.
\end{construction}
\begin{proof}
The finite binomial formula writes
\[ u_n=\sum_{j=0}^n\frac1{j!}\prod_{\ell=0}^{j-1}(1-\ell/n). \]
For a fixed $j\le n$ every factor is nonnegative and nondecreasing as $n$ increases; the new term at $n+1$ is positive. Thus $u_n$ increases. It is bounded by $\sum_{j=0}^n1/j!\le 1+\sum_{j=1}^n2^{1-j}<3$, since $j!\ge2^{j-1}$ for $j\ge1$. Monotone convergence gives a limit $u\ge2$.
For the comparison use the later definition $\exp(1)=\sum_{j=0}^\infty1/j!$ in \KRef{K03.07.exp_construction}. For fixed $J$, the first $J+1$ terms of the displayed binomial sum tend to $\sum_{j=0}^J1/j!$. All remaining terms are nonnegative and their sum is at most $\sum_{j>J}1/j!\le\sum_{j>J}2^{1-j}$, uniformly in $n$. Letting first $n$ and then $J$ tend to infinity identifies $u$. This last paragraph is a comparison after the series construction, not a prerequisite for it.
\end{proof}

\KEntry{K03.01.counterexamples}{Boundary cases for sequences}
\begin{example}
The bounded sequence $(-1)^n$ diverges because its even and odd subsequences tend to $1$ and $-1$. The sequence $n$ has no finite limit because it is unbounded, but tends to $+\infty$ since $n>B$ eventually for every $B$. Rational approximations to $\sqrt2$ within $1/n$ converge in $\R$ and are Cauchy in $\Q$, yet have no rational limit, as proved in \KRef{K03.01.cauchy}. Thus boundedness alone, existence of one convergent subsequence, and the Cauchy property without a complete target are different conditions.
\end{example}

% BEGIN REVISED SOURCE chapters/M03.02.tex SHA256=fc338c9dc76be5aadf6f98123f669abcd2b287e81d93d8db4c61df9d31f0ee57
\chapter{Function limits and asymptotic comparison}
All limits in this chapter have an explicitly chosen approach. Unless indicated otherwise $E\subset\R$ and $f:E\to\K$, where $\K=\R$ or $\C$. The same finite-point metric arguments also work for $E\subset\C$.
\SourceRef{Zorich I, 3.2; Rudin, Chapter 4; Olver, Chapter 1}

\KEntry{K03.02.relative_limit}{Limits relative to a domain}
\begin{definition}
A point $a$ is an accumulation point of $E$ if every $\delta>0$ admits $x\in E$ with $0<|x-a|<\delta$. We write $\lim_{E\ni x\to a}f(x)=L\in\K$ if
\[ (\forall\eps>0)(\exists\delta>0)(\forall x\in E)\quad 0<|x-a|<\delta\ \Longrightarrow\ |f(x)-L|<\eps. \]
The value at $a$, if defined, is absent from this condition. We require accumulation whenever discussing uniqueness of a punctured limit.
\end{definition}
\begin{proposition}
At an accumulation point the limit is unique.
\end{proposition}
\begin{proof}
If $L\ne M$ were limits, take their common neighborhood for errors $|L-M|/3$ and choose an approach point in it. The triangle inequality gives $|L-M|<2|L-M|/3$. Without accumulation the implication can be vacuous for every $L$, so the hypothesis cannot be removed.
\end{proof}

\KEntry{K03.02.one_sided}{One-sided and infinite limits}
\begin{definition}
The right-hand limit uses $x\in E\cap(a,a+\delta)$, and the left-hand limit uses $x\in E\cap(a-\delta,a)$. Each stated one-sided limit presupposes approach points arbitrarily near $a$ on that side. At $+\infty$, replace $0<|x-a|<\delta$ by $x>A$; at $-\infty$, use $x<-A$, where $A$ may be taken positive. Such approaches require $E$ to be unbounded in the indicated direction. For a real-valued function, $f\to+\infty$ means that for every $B\in\R$, $f(x)>B$ eventually along the stated approach; $f\to-\infty$ uses $f(x)<B$. For complex values only $|f|\to\infty$ has this order-free meaning.
\end{definition}
If both sides accumulate, the two-sided finite limit exists with value $L$ exactly when both one-sided limits are $L$: the forward implication restricts a neighborhood, and the reverse takes the smaller of its two radii. On $E=(0,\infty)$ an approach to zero is a right-hand relative limit; adding a nonexistent left side conveys no extra information.

\KEntry{K03.02.sequential}{The sequential criterion}
\begin{theorem}
At an accumulation point $a$, $f(x)\to L$ within $E$ if and only if $f(x_n)\to L$ for every sequence $x_n\in E\setminus\{a\}$ with $x_n\to a$. The analogous criterion holds for a chosen side and for approaches to $\pm\infty$.
\end{theorem}
\begin{proof}
A limit estimate applies to every sufficiently late term of such a sequence. Conversely if the limit condition fails, some $\eps_0>0$ admits, for each $n$, a point $x_n\in E$ with $0<|x_n-a|<1/n$ and $|f(x_n)-L|\ge\eps_0$. This sequence contradicts the asserted sequential property. For a side choose each point on that side; for $+\infty$ choose $x_n>n$ when constructing a failed test sequence, and similarly for $-\infty$.
\end{proof}

\KEntry{K03.02.algebra}{Algebra and composition of limits}
\begin{theorem}
Along a common nontrivial approach, if $f\to A$ and $g\to B$ are finite, then $cf+dg\to cA+dB$, $fg\to AB$, and $|f|\to|A|$. If $B\ne0$, eventually $g\ne0$ and $f/g\to A/B$.
If $f:E\to D$, $f(x)\to b$, and $g:D\to\K$ is continuous at $b\in D$, then $g(f(x))\to g(b)$. Alternatively, if $g(y)\to C$ as $D\ni y\to b$ in the punctured sense and $f(x)\ne b$ eventually, then $g(f(x))\to C$.
\end{theorem}
\begin{proof}
The sequential criterion reduces every algebra assertion to \KRef{K03.01.limit_algebra}; alternatively its inequalities apply verbatim on small neighborhoods. For continuous composition, choose $\eta$ so that $y\in D$, $|y-b|<\eta$ implies $|g(y)-g(b)|<\eps$, then choose a sufficiently small input neighborhood with $|f(x)-b|<\eta$. For punctured composition the same argument requires $0<|f(x)-b|$, supplied by eventual avoidance. The latter condition matters: a punctured limit of $g$ says nothing about its actual value at $b$, while $f$ may equal $b$ repeatedly.
\end{proof}

\KEntry{K03.02.squeeze}{Comparison along an approach}
\begin{theorem}
Along the same domain and side, eventual real inequalities $f\le g\le h$ and $f,h\to L$ imply $g\to L$. An eventual estimate $|u-L|\le v$ with $v\to0$ implies $u\to L$. Eventual real $f\le g$ and finite limits imply $\lim f\le\lim g$.
\end{theorem}
\begin{proof}
Intersect the neighborhoods in which the inequalities hold with neighborhoods where both outer functions lie in $(L-\eps,L+\eps)$. The middle function then lies there too. The absolute estimate is the definition of a limit. The order assertion follows by contradiction, choosing disjoint neighborhoods of two incorrectly ordered limits. These arguments also apply on tails at infinity.
\end{proof}

\KEntry{K03.02.cauchy}{A Cauchy criterion for finite limits}
\begin{theorem}
At an accumulation point $a$, a function $f:E\to\K$ has a finite limit if and only if for every $\eps>0$ some $\delta>0$ satisfies
\[ |f(x)-f(y)|<\eps\quad( x,y\in E,\ 0<|x-a|,|y-a|<\delta). \]
There are identical one-sided and infinite-endpoint versions, with both points in the selected approach neighborhoods.
\end{theorem}
\begin{proof}
A finite limit gives this estimate by two errors $<\eps/2$. Conversely choose $x_n\in E$ with $0<|x_n-a|<1/n$. The condition makes $f(x_n)$ Cauchy, so by scalar completeness it tends to some $L$. Given $\eps>0$, choose $\delta$ for pairwise error $<\eps/2$ and a sufficiently late $x_n$ in this neighborhood with $|f(x_n)-L|<\eps/2$. Every approach point $x$ in the same neighborhood satisfies $|f(x)-L|<\eps$. For a side or infinity choose a sequence of points in its nested approach neighborhoods and use the same proof.
\end{proof}

\KEntry{K03.02.oscillation}{Punctured and unpunctured oscillation}
\begin{definition}
For a nonempty set $A\subset E$ set $\osc(f,A)=\sup_{x,y\in A}|f(x)-f(y)|\in[0,+\infty]$, and set its value to zero for the empty set. Define punctured oscillation at an accumulation point by
\[ \omega^*(f,a)=\inf_{\delta>0}\osc\bigl(f,E\cap\{0<|x-a|<\delta\}\bigr). \]
If $a\in E$, unpunctured oscillation $\omega(f,a)$ instead uses $E\cap\{|x-a|<\delta\}$.
\end{definition}
\begin{proposition}
$\omega^*(f,a)=0$ is equivalent to existence of a finite punctured limit. At $a\in E$, $\omega(f,a)=0$ is equivalent to continuity at $a$.
\end{proposition}
\begin{proof}
An infimum of these nonnegative nested oscillations is zero exactly when, for every $\eps>0$, some neighborhood has oscillation $<\eps$. The punctured statement is \KRef{K03.02.cauchy}. In an unpunctured neighborhood one may set $y=a$, obtaining $|f(x)-f(a)|<\eps$. Conversely continuity gives oscillation $\le2\eps$ on a sufficiently small neighborhood. For example $f(x)=0$ when $x\ne0$, $f(0)=1$ has $\omega^*(f,0)=0$ and $\omega(f,0)=1$.
\end{proof}

\KEntry{K03.02.big_o}{Big-O and little-o}
\begin{definition}
Fix an approach within a common domain $E$. For scalar functions $f,g$, write $f=O(g)$ if eventually $|f|\le C|g|$ for some finite $C\ge0$. Write $f=o(g)$ if for every $\eps>0$ eventually $|f|\le\eps|g|$. The neighborhood may depend on $\eps$ in the latter definition. These are inequalities, so they remain meaningful at zeros of $g$ and force $f=0$ there eventually. If $g$ is eventually nonzero, they are equivalent respectively to eventual boundedness of $f/g$ and to $f/g\to0$.
\end{definition}
The ratio characterizations follow by division by $|g|>0$ in both directions. In particular $o(g)$ implies $O(g)$ by taking $\eps=1$.

\KEntry{K03.02.equivalence}{Asymptotic equivalence}
\begin{definition}
Write $f\sim g$ along a fixed approach if $g$ is eventually nonzero and $f/g\to1$.
\end{definition}
\begin{proposition}
Under eventual nonvanishing of $g$, $f\sim g$ if and only if $f-g=o(g)$. Equivalence implies eventual nonvanishing of $f$, $f=O(g)$, $g=O(f)$, symmetry, and transitivity.
\end{proposition}
\begin{proof}
The identity $(f-g)/g=f/g-1$ gives the first assertion. Eventually $|f/g-1|<1/2$, so $1/2<|f/g|<3/2$. This proves nonvanishing and the two bounds. The reciprocal of a limit $1$ is $1$, and products of two such ratios tend to $1$, proving symmetry and transitivity on a common tail. Without eventual nonzero denominators this ratio definition is not being asserted.
\end{proof}

\KEntry{K03.02.o_algebra}{Calculus of error bounds}
\begin{proposition}
All the following statements concern one common approach. Bounds $f=O(g)$, $g=O(h)$ imply $f=O(h)$; if either of the two bounds is little-o, their composite is little-o. Sums satisfy $O(g)+O(g)\subset O(g)$ and $o(g)+o(g)\subset o(g)$. More generally $f=O(g)$, $u=O(v)$ imply $fu=O(gv)$, with little-o for the product if either factor bound is little-o. Fixed scalar multiples preserve both bounds. If $v$ is eventually nonzero, dividing an estimate by $v$ gives the corresponding estimate with comparison function divided by $v$.
\end{proposition}
\begin{proof}
Intersect all required neighborhoods. Transitivity of big-O follows from $|f|\le C|g|\le CD|h|$. If $f=o(g)$ choose its bound with coefficient $\eps/\max(1,D)$; if $g=o(h)$ use $\eps/\max(1,C)$. For sums add the inequalities, taking coefficients $\eps/2$ for two little-o terms. For products multiply the bounds and use the same coefficient choice as for transitivity. Scaling multiplies constants; division divides both sides of the pointwise inequality by $|v|$. These arguments permit zeros in comparison functions and introduce nonvanishing only when division actually occurs.
\end{proof}

\KEntry{K03.02.equivalence_replacement}{Replacement of factors and failure under subtraction}
\begin{proposition}
If $f\sim F$ and $g\sim G$, then $fg\sim FG$ and $f/g\sim F/G$, on their eventual domains. Subtraction admits no analogous unrestricted replacement rule.
\end{proposition}
\begin{proof}
The first quotient of products is $(f/F)(g/G)\to1$, and the quotient of ratios is $(f/F)/(g/G)\to1$; eventual nonvanishing follows from \KRef{K03.02.equivalence}. For the failure, as $x\to0$ through nonzero reals, $x+x^2\sim x$ and $x\sim x$, but $(x+x^2)-x=x^2$ whereas replacing both terms gives $x-x=0$. The residual leading term was destroyed by cancellation, and $x^2$ is not asymptotically equivalent to zero under our definition.
\end{proof}

\KEntry{K03.02.growth_scale}{Logarithms, powers, and exponentials}
\begin{theorem}
After the constructions in \KRef{K03.07.exp_construction}, \KRef{K03.07.log_construction}, and \KRef{K03.07.real_powers}, for fixed real $m,k\ge0$ and $\eps,c>0$,
\[ (\log x)^m=o(x^\eps),\qquad x^k=o(\exp(cx))\quad(x\to+\infty). \]
\end{theorem}
\begin{proof}
For $y\ge0$ and any integer $N$, the positive-term exponential series gives $\exp(\eps y)\ge (\eps y)^N/N!$. Choose $N>m$. For $y\ge1$, $y^m/\exp(\eps y)\le N!\eps^{-N}y^{m-N}\to0$, since positive powers tend to infinity by the real-power properties in module 07. Put $y=\log x\to+\infty$ and $\exp(\eps\log x)=x^\eps$. Likewise choose an integer $N>k$ and use $\exp(cx)\ge(cx)^N/N!$, obtaining $x^k/\exp(cx)\le N!c^{-N}x^{k-N}\to0$. None of these facts is used in constructing the earlier sequence limits or the general series theory.
\end{proof}

\KEntry{K03.02.path_dependence}{Different approaches can give different answers}
\begin{example}
On $\R\setminus\{0\}$ the function $x/|x|$ is $1$ on the right and $-1$ on the left. Hence its two one-sided limits at zero differ and it has no two-sided limit. Its values along $x_n=1/n$ all equal $1$, showing that one convergent sampling sequence cannot prove the full limit. The sequential criterion requires every approach sequence; $-1/n$ already defeats this claim.
\end{example}

\KEntry{K03.02.domain_zeros}{A bounded ratio need not converge}
\begin{example}
For $x\ne0$ put $f(x)=x\sin(1/x)$ and $g(x)=x$. The trigonometric construction and identities in \KRef{K03.07.trig_properties} give $|f(x)|\le|g(x)|$, so $f=O(g)$ near zero and $f\to0$. But $f/g=\sin(1/x)$ has no limit: at
\[ x_n=(\pi/2+2\pi n)^{-1},\qquad y_n=(3\pi/2+2\pi n)^{-1} \]
it equals $1$ and $-1$, respectively. Both sequences tend to zero. Thus $f$ is neither $o(g)$ nor equivalent to $g$. The repeated zeros of $f$ cause no difficulty for the inequality definition of big-O, while ratio manipulations always retain the nonzero denominator $g(x)=x$ on the punctured domain.
\end{example}

% BEGIN REVISED SOURCE chapters/M03.03.tex SHA256=6dbeff7f2afcffaa77143b7ad6dc8009cfa5f0c212339ef287f803415247737f
\chapter{Continuity, compactness, and intervals}
Metric terminology in this chapter is relative to the stated domain. An open set in $E$ is an intersection of $E$ with an ambient open set; closed sets are their complements. Compactness means that every relative open cover has a finite subcover.
\SourceRef{Zorich I, Chapter 4; Rudin, Chapter 4; Tao, Chapter 9; Pugh, 2.2}

\KEntry{K03.03.continuity}{Continuity at a point}
\begin{definition}
For metric $E$ and $f:E\to\K$, continuity at $a\in E$ means that for every $\eps>0$ some $\delta>0$ satisfies $|f(x)-f(a)|<\eps$ whenever $x\in E$ and $d(x,a)<\delta$. A function is continuous if it is continuous at every point. If $a$ accumulates $E$, this is equivalent to the relative punctured limit $f(x)\to f(a)$. At an isolated point the condition holds automatically.
\end{definition}
\begin{proposition}
Continuity at $a$ is equivalent to $f(x_n)\to f(a)$ for every sequence $x_n\to a$ in $E$.
\end{proposition}
\begin{proof}
The forward implication inserts sufficiently late sequence terms into a continuity neighborhood. If continuity fails, choose $\eps_0>0$ and $x_n$ with $d(x_n,a)<1/n$ and $|f(x_n)-f(a)|\ge\eps_0$, contradicting sequential continuity. This argument includes sequences equal to $a$ and isolated points.
\end{proof}

\KEntry{K03.03.continuity_algebra}{Operations on continuous functions}
\begin{proposition}
Finite scalar linear combinations and products of continuous scalar functions are continuous. Quotients are continuous where the denominator is nonzero. Composites of continuous functions with compatible domains are continuous, including metric intermediate spaces.
\end{proposition}
\begin{proof}
Apply the sequential continuity criterion to any converging input sequence and then \KRef{K03.01.limit_algebra}. For composition, continuity of the inner map first gives convergence to its actual value, and continuity of the outer map then gives convergence of the composite. The reciprocal estimate in that algebra theorem requires precisely a nonzero value at the point under consideration.
\end{proof}

\KEntry{K03.03.local_bounds}{Local boundedness, signs, and denominators}
\begin{proposition}
If $f$ is continuous at $a$, then on a relative neighborhood $|f(x)|\le|f(a)|+1$. If $f(a)\ne0$, on a possibly smaller neighborhood $|f(x)|\ge|f(a)|/2$. For real $f$ with $f(a)>0$, $f(x)>f(a)/2>0$ locally.
\end{proposition}
\begin{proof}
Use continuity with $\eps=1$ for the first bound, and with $\eps=|f(a)|/2$ and the reverse triangle inequality for the second. For the positive real case $f(x)>f(a)-f(a)/2$. In particular a denominator continuous and nonzero at $a$ stays nonzero nearby.
\end{proof}

\KEntry{K03.03.closed_preimage}{Open and closed preimages}
\begin{theorem}
A map between metric subspaces is continuous if and only if every relative open set has relative open preimage, and if and only if every relative closed set has relative closed preimage.
\end{theorem}
\begin{proof}
For a continuous map and $x\in f^{-1}(U)$ with $U$ open, choose a target ball about $f(x)$ lying in $U$, then a source ball mapped into it. Thus every point of $f^{-1}(U)$ has a neighborhood there. Conversely the preimage of the target ball of radius $\eps$ around $f(a)$ is an open set containing $a$, so contains a source ball witnessing continuity. The closed formulation is equivalent because $f^{-1}(Y\setminus F)=E\setminus f^{-1}(F)$.
\end{proof}

\KEntry{K03.03.compact_interval}{Compactness and sequential compactness}
\begin{theorem}
In a metric space, compactness is equivalent to sequential compactness: every sequence has a subsequence converging to a point of the space. A subset of $\R$ (or $\C$) is compact exactly when it is closed and bounded. In particular every closed bounded interval is compact.
\end{theorem}
\begin{proof}
First suppose a metric space $K$ is compact. If a sequence had no convergent subsequence, then for every $x\in K$ some ball about $x$ would contain only finitely many sequence indices. Indeed infinitely many indices in every ball would let us choose increasing indices in balls of radii $1/j$, yielding a subsequence tending to $x$. These balls cover $K$; a finite subcover would imply that the sequence has only finitely many indices, a contradiction.
Conversely suppose $K$ is sequentially compact. For each $\eta>0$ it is covered by finitely many $\eta$-balls: otherwise choose recursively points at pairwise distance at least $\eta$, a sequence without any convergent subsequence. Given an open cover, there is some $\delta>0$ such that every relative $\delta$-ball is contained in one cover member. If not, choose $x_n$ whose relative $1/n$-ball is contained in no member. A subsequence tends to $x$; a cover member containing $x$ contains a relative ball of radius $r>0$ about $x$. For sufficiently late terms, $d(x_n,x)<r/2$ and $1/n<r/2$, a contradiction. Choose finitely many $\delta/2$-balls covering $K$, and for each center a cover member containing its $\delta$-ball. These members give a finite subcover.
Now a closed bounded scalar set is sequentially compact by \KRef{K03.01.bolzano}, since a sequence limit remains in a closed set: if it were outside, the open complement would eventually contain the sequence. Conversely a compact scalar set is bounded by a finite subcover of the balls centered at zero with positive integer radii. It is closed because any point in its closure is a limit of a sequence in the set (choose distances $<1/n$); sequential compactness gives a subsequence with limit in the set, and uniqueness identifies this point with the original limit. The empty set satisfies all assertions directly.
\end{proof}

\KEntry{K03.03.extreme_values}{Continuous images and extreme values}
\begin{theorem}
The continuous image of a compact set is compact. Hence a continuous real function on a nonempty compact set attains its maximum and minimum, and a continuous scalar function on a compact set is bounded.
\end{theorem}
\begin{proof}
Pull an open cover of the image back to the domain and use a finite subcover there. For a scalar target compactness gives boundedness and closedness by \KRef{K03.03.compact_interval}. For real $f$ the nonempty bounded image has a supremum $s$. Select image points in $(s-1/n,s]$; they tend to $s$, which belongs to the image by closedness. Negate $f$ to obtain a minimum. Nonemptiness is required for attainment. On $(0,1)$, $f(x)=x$ attains neither bound; on $\R$, $f(x)=x$ is continuous and unbounded.
\end{proof}

\KEntry{K03.03.intermediate_values}{The intermediate value theorem}
\begin{theorem}
If $f:[a,b]\to\R$ is continuous and $a\le b$, every value between $f(a)$ and $f(b)$ occurs.
\end{theorem}
\begin{proof}
The case $a=b$ and the endpoint values are immediate. For a strictly intermediate $v$ suppose first $f(a)<v<f(b)$. Let $S=\{x\in[a,b]:f(x)\le v\}$ and $c=\sup S$. Continuity near $b$ gives $c<b$, and $a\le c$. Choose points of $S$ tending to $c$ by supremum approximation; continuity gives $f(c)\le v$. If $f(c)<v$, continuity and $c<b$ give a point of $S$ strictly beyond $c$, impossible. Thus $f(c)=v$. When the endpoint order is reversed apply the result to $-f$. No differentiability has been used.
\end{proof}

\KEntry{K03.03.interval_images}{Intervals and connectedness}
\begin{definition}
A subset $E$ of a metric space is connected if it cannot be written as the disjoint union of two nonempty relatively open subsets. A real set is an interval if it contains every point between any two of its points; this convention includes the empty set and singletons.
\end{definition}
\begin{theorem}
A real set is connected if and only if it is an interval. Continuous images of connected spaces are connected; consequently continuous real functions map intervals to intervals.
\end{theorem}
\begin{proof}
If a real set contains $x<y$ but omits $z\in(x,y)$, its intersections with $(-\infty,z)$ and $(z,\infty)$ disconnect it. Conversely suppose an interval $I$ were disconnected as $U\cup V$ and take $u\in U$, $v\in V$, relabeling so $u<v$. The function equal to $0$ on $U$ and $1$ on $V$ is continuous, since each point has a relative neighborhood in one piece. Its restriction to $[u,v]\subset I$ contradicts \KRef{K03.03.intermediate_values}. Finally a separation of a continuous image pulls back to a separation of its domain by open preimages. This proves the image assertion, also obtainable by applying the intermediate value theorem separately to every segment in the domain.
\end{proof}

\KEntry{K03.03.monotone_limits}{One-sided limits and jumps of monotone functions}
\begin{theorem}
For a nondecreasing real function on an interval and an interior point $x$,
\[ f(x-)=\sup_{t<x}f(t),\qquad f(x+)=\inf_{t>x}f(t) \]
exist and are finite, with $f(x-)\le f(x)\le f(x+)$. At an interval endpoint the available one-sided limit exists in $\overline\R$. Interior discontinuities are jumps and there are at most countably many. The decreasing version follows by negation.
\end{theorem}
\begin{proof}
For the left formula, values to the left are bounded above by $f(x)$ and bounded below on a sufficiently short left neighborhood by a fixed interior value. If $L$ is their supremum choose $t_0<x$ with $f(t_0)>L-\eps$; then $L-\eps<f(t)\le L$ whenever $t_0<t<x$. The right formula is the infimum argument. At endpoints the same argument applies, using arbitrary finite bounds when a supremum or infimum is infinite. At an interior point continuity is equivalent to equality of both one-sided limits with $f(x)$, so a discontinuity has $f(x-)<f(x+)$. Assign to each such point a rational in the nonempty interval $(f(x-),f(x+))$, choosing the first one in a fixed enumeration. For $x<y$, monotonicity gives $f(x+)\le f(y-)$: choose $t$ between $x$ and $y$ to bound both limits. The assigned rational intervals are therefore disjoint, making the assignment injective. At most two included endpoints can add discontinuities, so the total remains countable.
\end{proof}

\KEntry{K03.03.inverse_continuity}{Continuous monotone inverses}
\begin{theorem}
A continuous strictly monotone real function $f$ on an interval $I$ is a bijection onto an interval $J=f(I)$ and has a continuous inverse $J\to I$.
\end{theorem}
\begin{proof}
Strict monotonicity gives injectivity, and \KRef{K03.03.interval_images} gives the interval image. Suppose $f$ increases, $y_0=f(x_0)$, and fix $\eps>0$. If $I$ contains points at least $\eps$ below $x_0$, its interval property provides $u\in I$ with $x_0-\eps<u<x_0$; choose such a $u$ and require $|y-y_0|<y_0-f(u)$. This forces $f^{-1}(y)>u>x_0-\eps$. If no such points exist the lower bound is automatic. Treat the upper side with $v\in(x_0,x_0+\eps)$ whenever needed and require $|y-y_0|<f(v)-y_0$. The minimum of the available positive bounds gives inverse continuity, also at endpoints of $J$. Negation treats decreasing $f$.
For the domain $I=[0,1)\cup\{2\}$, define $f(x)=x$ for $x<1$ and $f(2)=1$. This is continuous and strictly increasing on that domain, but its inverse on $[0,1]$ is discontinuous at $1$, since $f^{-1}(1)=2$ while $f^{-1}(1-1/n)\to1$. The interval hypothesis has content.
\end{proof}

\KEntry{K03.03.uniform_continuity}{Uniform continuity}
\begin{definition}
A scalar function on a metric set $E$ is uniformly continuous if
\[ (\forall\eps>0)(\exists\delta>0)(\forall x,y\in E)\quad d(x,y)<\delta\Longrightarrow |f(x)-f(y)|<\eps. \]
Thus $\delta$ depends on $\eps$, but not on a chosen base point. Uniform continuity implies continuity by fixing $y=a$. Its failure is equivalent to existence of $\eps_0>0$ and pairs $x_n,y_n$ with $d(x_n,y_n)\to0$ and $|f(x_n)-f(y_n)|\ge\eps_0$: choose distances $<1/n$ from the negated definition, and use the definition to prove the reverse implication.
\end{definition}

\KEntry{K03.03.heine_cantor}{Heine--Cantor}
\begin{theorem}
A continuous scalar function on a compact metric set is uniformly continuous.
\end{theorem}
\begin{proof}
Otherwise take the pairs in \KRef{K03.03.uniform_continuity}. A subsequence $x_{n_j}$ tends to some $x$ by compactness. Since $d(y_{n_j},x)\le d(y_{n_j},x_{n_j})+d(x_{n_j},x)$, the paired subsequence tends to $x$ too. Continuity makes both image subsequences tend to $f(x)$, contradicting their fixed separation. A Lipschitz bound does not follow, as the square-root example below shows.
\end{proof}

\KEntry{K03.03.lipschitz}{Lipschitz and Holder estimates}
\begin{proposition}
If $|f(x)-f(y)|\le C d(x,y)^\alpha$ for all $x,y\in E$, with fixed $C\ge0$ and $0<\alpha\le1$, then $f$ is uniformly continuous. The case $\alpha=1$ is Lipschitz continuity; $0<\alpha<1$ is Holder continuity.
\end{proposition}
\begin{proof}
If $C=0$ the function is constant. Otherwise choose $\delta=(\eps/C)^{1/\alpha}$ using the positive real power of module 07 (or any smaller positive number whose $\alpha$ power is small enough). Then the displayed bound is $<\eps$ whenever $d(x,y)<\delta$. In the common algebraic case $\alpha=1/m$ this already uses only integer roots. Such an estimate does not imply differentiability: $|x|$ is $1$-Lipschitz by the reverse triangle inequality but its two difference quotients at zero are $1$ and $-1$.
\end{proof}

\KEntry{K03.03.extension}{Extension to the closure}
\begin{theorem}
A uniformly continuous function $f:E\to\K$ on a subset of a metric space extends uniquely to a continuous function $\bar f:\overline E\to\K$. This extension is uniformly continuous.
\end{theorem}
\begin{proof}
The empty case is immediate. For $x\in\overline E$ choose $x_n\in E$ tending to $x$. Such a sequence is Cauchy; uniform continuity makes $f(x_n)$ Cauchy, hence convergent in the complete scalar target. Define $\bar f(x)$ as its limit. If $y_n\to x$ is another choice, $d(x_n,y_n)\to0$, so uniform continuity gives $|f(x_n)-f(y_n)|\to0$ and the definition is independent of choice. Constant sequences show agreement on $E$.
Given $\eps>0$, choose $\delta>0$ for $f$ with output error $<\eps/2$. If $x,y\in\overline E$ and $d(x,y)<\delta/3$, approximating sequences $x_n,y_n\in E$ eventually satisfy $d(x_n,y_n)<\delta$. Passing to limits gives $|\bar f(x)-\bar f(y)|\le\eps/2<\eps$. Hence the extension is uniformly continuous. Every continuous extension must take at $x$ the limit along $x_n\to x$, proving uniqueness.
\end{proof}

\KEntry{K03.03.fixed_point}{The contraction principle}
\begin{theorem}
Let $(X,d)$ be a nonempty complete metric space, and $T:X\to X$ satisfy $d(Tx,Ty)\le qd(x,y)$ with $0\le q<1$. Starting at any $x_0\in X$, put $x_{n+1}=Tx_n$. There is a unique fixed point $p$, and
\[ d(x_n,p)\le \frac{q^n}{1-q}d(x_1,x_0),\qquad d(x_n,p)\le\frac{d(x_{n+1},x_n)}{1-q}. \]
The first bound uses $q^0=1$, also when $q=0$.
\end{theorem}
\begin{proof}
Induction gives $d(x_{j+1},x_j)\le q^jd(x_1,x_0)$. For $m>n$, summing yields $d(x_m,x_n)\le q^n d(x_1,x_0)/(1-q)$, so the sequence is Cauchy and converges to $p\in X$. Lipschitz continuity gives $Tx_n\to Tp$, while $Tx_n=x_{n+1}\to p$, hence $Tp=p$. Taking $m\to\infty$ proves the first bound. Applying the same geometric estimate starting at $x_n$ proves the second. If $r$ is another fixed point, $d(p,r)\le qd(p,r)$ forces equality to zero. The condition $T:X\to X$ is essential for all iterates and their complete limit to lie in the specified space.
\end{proof}

\KEntry{K03.03.counterexamples}{Separating the continuity conditions}
\begin{example}
The function $1/x$ is continuous on $(0,1)$ but unbounded and not uniformly continuous: for $x_n=1/(n+1)$, $y_n=1/(n+2)$ the input distances tend to zero and the image distances equal $1$. The continuous function $x^2$ on $\R$ is not uniformly continuous, since $x_n=n$, $y_n=n+1/n$ have input distance $1/n$ and image distance $2+1/n^2$. On $[0,1]$, $\sqrt{x}$ is uniformly continuous because
\[ |\sqrt x-\sqrt y|^2\le |x-y|: \quad x\ge y\Longrightarrow (\sqrt x-\sqrt y)^2\le(\sqrt x-\sqrt y)(\sqrt x+\sqrt y). \]
It is not Lipschitz, since $|\sqrt h-0|/|h-0|=1/\sqrt h$ is unbounded as $h\downarrow0$. Finally the step function $f(x)=0$ for $x<0$, $f(x)=1$ for $x\ge0$ is bounded but discontinuous at zero, as the sequence $-1/n$ shows. These computations separate boundedness, continuity, uniform continuity, and Lipschitz continuity.
\end{example}

% BEGIN REVISED SOURCE chapters/M03.04.tex SHA256=9557a855c9dd37fac213e66f36372ca5bedfe2ead4807b604f87d46474ccea5e
\chapter{Derivatives and differential rules}
The independent variable is real, even when the value lies in $\C$. Unless an endpoint is explicitly mentioned, differentiation takes place at an interior point of an interval.
\SourceRef{Zorich I, 5.1--5.2; Apostol, Chapter 5; Pugh, 3.1}

\KEntry{K03.04.derivative}{The real-variable derivative}
\begin{definition}
For $f:I\to\K$ and interior $a\in I$, the derivative is the finite limit
\[ f'(a)=\lim_{h\to0,\ h\in\R\setminus\{0\}}\frac{f(a+h)-f(a)}h. \]
For $f=u+iv$, existence is equivalent to existence of $u'(a)$ and $v'(a)$, and $f'=u'+iv'$. This follows by taking real and imaginary parts of the difference quotients and using $|u+iv|\le|u|+|v|$ for the converse. It is not the definition of a derivative with complex increments; that separate notion belongs to complex analysis.
\end{definition}

\KEntry{K03.04.one_sided}{One-sided derivatives}
\begin{definition}
The right derivative $f'_+(a)$ uses $h\downarrow0$ and the left derivative $f'_-(a)$ uses $h\uparrow0$, wherever the corresponding side of the domain contains an interval. A derivative at an included endpoint means the available one-sided derivative only when that convention is stated.
\end{definition}
\begin{proposition}
At an interior point a finite ordinary derivative exists if and only if both finite one-sided derivatives exist and agree.
\end{proposition}
\begin{proof}
A two-sided limit restricts to each side. If the one-sided limits equal $L$, choose for each $\eps$ a radius on either side making the quotient within $\eps$ of $L$; their minimum works for every nonzero $h$ of sufficiently small modulus.
\end{proof}

\KEntry{K03.04.peano}{Differentiability as a linear approximation}
\begin{theorem}
$f$ is differentiable at $a$ with derivative $L$ if and only if
\[ f(a+h)=f(a)+Lh+r(h),\qquad r(h)=o(h)\quad(h\to0). \]
Differentiability implies continuity at $a$.
\end{theorem}
\begin{proof}
For $h\ne0$, $r(h)/h$ is the difference between the difference quotient and $L$, so the two conditions are equivalent by \KRef{K03.02.big_o}. Since $r(h)=o(h)$ implies $r(h)\to0$, the increment $Lh+r(h)$ tends to zero. This is continuity. For example $|h|$ tends to zero but has unequal left and right quotients, so the converse is false.
\end{proof}

\KEntry{K03.04.linear_rules}{Linear rules}
\begin{proposition}
If $f_1,\ldots,f_m:I\to\K$ are differentiable at $a$ and $c_1,\ldots,c_m\in\K$ are constant, then
\[ \left(\sum_{j=1}^m c_jf_j\right)'(a)=\sum_{j=1}^m c_j f_j'(a). \]
In particular constants have derivative zero and the identity has derivative one.
\end{proposition}
\begin{proof}
The difference quotient of the sum is the finite sum of the difference quotients with the same constants. Limit algebra gives the formula. Constant and identity difference quotients are identically zero and one, respectively.
\end{proof}

\KEntry{K03.04.product}{Products, reciprocals, and quotients}
\begin{theorem}
For scalar functions differentiable at $a$,
\[ (fg)'(a)=f'(a)g(a)+f(a)g'(a). \]
If $g(a)\ne0$, it stays nonzero near $a$ and
\[ (1/g)'(a)=-\frac{g'(a)}{g(a)^2},\qquad (f/g)'(a)=\frac{f'(a)g(a)-f(a)g'(a)}{g(a)^2}. \]
\end{theorem}
\begin{proof}
Subtract $f(a)g(a)$ by first subtracting $f(a)g(a+h)$ to get the quotient
\[ \frac{f(a+h)-f(a)}h g(a+h)+f(a)\frac{g(a+h)-g(a)}h. \]
Differentiability gives continuity, hence the stated limit. If $g(a)\ne0$, continuity bounds $|g(a+h)|$ below by $|g(a)|/2$, and
\[ \frac{1/g(a+h)-1/g(a)}h=-\frac{g(a+h)-g(a)}{h\,g(a+h)g(a)}\to-\frac{g'(a)}{g(a)^2}. \]
The product rule applied to $f(1/g)$ proves the last formula.
\end{proof}

\KEntry{K03.04.chain}{The chain rule}
\begin{theorem}
Let $f:I\to\R$ be differentiable at $a$, and let $g:J\to\K$ be differentiable at $b=f(a)$, where $J$ contains a neighborhood of $b$ and $f$ maps a neighborhood of $a$ into $J$. Then
\[ (g\circ f)'(a)=g'(f(a))f'(a). \]
\end{theorem}
\begin{proof}
Set $k=f(a+h)-b=f'(a)h+o(h)$, so $k=O(h)$ and $k\to0$. Write $g(b+k)-g(b)=g'(b)k+\eta(k)k$, with $\eta(k)\to0$ for $k\ne0$ and $\eta(0)=0$. Such a remainder exists by differentiability. Substitution gives
\[ g(f(a+h))-g(f(a))=g'(b)f'(a)h+o(h)+\eta(k)O(h). \]
The last term is $o(h)$, including when $k=0$, and \KRef{K03.04.peano} finishes the proof. A complex-valued inner function would require an outer function on a complex domain, for which this theorem supplies no assertion.
\end{proof}

\KEntry{K03.04.inverse}{The inverse derivative}
\begin{theorem}
Let $f:I\to\R$ be continuous and strictly monotone on an interval, let $a$ be interior, and suppose $f'(a)\ne0$. Then its inverse $g:f(I)\to I$ is differentiable at $b=f(a)$, with
\[ g'(b)=\frac1{f'(a)}. \]
The corresponding available one-sided statement holds at an included endpoint if its one-sided derivative is finite and nonzero.
\end{theorem}
\begin{proof}
The inverse is continuous by \KRef{K03.03.inverse_continuity}. For $y\ne b$ put $x=g(y)\ne a$; as $y\to b$, $x\to a$, and
\[ \frac{g(y)-g(b)}{y-b}=\left(\frac{f(x)-f(a)}{x-a}\right)^{-1}\to\frac1{f'(a)}. \]
The image contains points on both sides of $b$ when $a$ is interior, by strict monotonicity on an interval. For endpoints the same ratio is taken on the available side.
\end{proof}

\KEntry{K03.04.higher}{Higher derivatives and regularity classes}
\begin{definition}
On an open interval set $f^{(0)}=f$ and recursively $f^{(j+1)}=(f^{(j)})'$ wherever defined. A function is $k$ times differentiable on the interval when all derivatives through $k$ exist there. It belongs to $C^k$ when these derivatives are continuous, and to $C^\infty$ when it belongs to $C^k$ for every finite $k$. On a compact interval, $C^k$ means the interior derivatives have continuous extensions to the endpoints agreeing with their available one-sided derivatives.
\end{definition}
Existence of a derivative gives continuity of the preceding derivative, but gives no continuity of the new derivative itself. The examples \KRef{K03.04.noncontinuous_derivative} and \KRef{K03.04.unbounded_derivative} quantify this distinction already at order one.

\KEntry{K03.04.leibniz}{The higher product rule}
\begin{theorem}
If $f,g$ are $k$ times differentiable on a common open interval, then
\[ (fg)^{(k)}=\sum_{j=0}^k\binom kj f^{(j)}g^{(k-j)}. \]
\end{theorem}
\begin{proof}
The case $k=0$ is the product itself. Differentiate the finite sum at order $k$ using the product rule. The coefficient of $f^{(j)}g^{(k+1-j)}$ for $1\le j\le k$ becomes $\binom{k}{j-1}+\binom{k}{j}=\binom{k+1}{j}$ by the finite factorial identity, and the two boundary coefficients remain one. Induction gives every order up to the available differentiability; continuity of the highest derivative is not required.
\end{proof}

\KEntry{K03.04.elementary_rules}{Elementary derivative formulas and their domains}
\begin{proposition}
The following real-variable formulas apply after the indicated functions have been constructed in module 07:
\[
\begin{array}{c|c|c}
 f(x)&f'(x)&\text{domain of this formula}\\\hline
 x^n&nx^{n-1}&x\in\R,\ n\in\N\\
 x^\alpha&\alpha x^{\alpha-1}&x>0,\ \alpha\in\R\\
 \exp x&\exp x&x\in\R\\
 \log x&1/x&x>0\\
 \sin x&\cos x&x\in\R\\
 \cos x&-\sin x&x\in\R\\
 \tan x&1/\cos^2x&\cos x\ne0\\
 \arcsin x&(1-x^2)^{-1/2}&-1<x<1\\
 \arccos x&-(1-x^2)^{-1/2}&-1<x<1\\
 \arctan x&(1+x^2)^{-1}&x\in\R
\end{array}
\]
The derivative of $x^0=1$ is zero everywhere; its formula does not require interpreting $0\cdot0^{-1}$ at zero. Negative integer powers obey the power rule on $\R\setminus\{0\}$.
\end{proposition}
\begin{proof}
Induction with the product rule proves positive integer powers; the reciprocal rule gives negative integers on their domain. The series differentiation arguments in \KRef{K03.07.exp_construction} and \KRef{K03.07.sin_cos} give the exponential and trigonometric formulas. Applying the inverse rule to the positive exponential gives the logarithm formula, as established in \KRef{K03.07.log_derivative}. The identity $x^\alpha=\exp(\alpha\log x)$ then gives $\alpha\exp(\alpha\log x)/x=\alpha x^{\alpha-1}$. The quotient rule and $\cos^2x+\sin^2x=1$ give the tangent formula. For $x=\sin y$, $y\in(-\pi/2,\pi/2)$, $\cos y=\sqrt{1-x^2}>0$, so the inverse rule gives the arcsine formula. For $x=\cos y$, $y\in(0,\pi)$, $\sin y=\sqrt{1-x^2}>0$, giving the arccosine sign. For $x=\tan y$, $y\in(-\pi/2,\pi/2)$, $1/\cos^2y=1+x^2$ gives the arctangent formula. The branch and range facts used here are proved in \KRef{K03.07.trig_inverses}; no formula is asserted at endpoints where its displayed denominator vanishes.
\end{proof}

\KEntry{K03.04.modulus}{Differentiating an absolute value}
\begin{proposition}
The real function $|x|$ has derivative $1$ for $x>0$, $-1$ for $x<0$, and no derivative at zero. If $f:I\to\C$ is differentiable at $a$ and $f(a)\ne0$, then
\[ (|f|)'(a)=\frac{\RePart(\overline{f(a)}f'(a))}{|f(a)|}. \]
If $f(a)=0$, then $|f|$ is differentiable at $a$ if and only if $f'(a)=0$, in which case its derivative is zero.
\end{proposition}
\begin{proof}
The first claim follows from the formulas $|x|=x$ and $|x|=-x$ on the two open half-lines; the one-sided quotients at zero are $1$ and $-1$. For the nonzero case rationalize the norm increment:
\[ \frac{|f(a+h)|-|f(a)|}{h}=\frac{|f(a+h)|^2-|f(a)|^2}{h\,(|f(a+h)|+|f(a)|)}. \]
The derivative of $|f|^2=u^2+v^2$ is $2(uu'+vv')=2\RePart(\bar f f')$, and the denominator tends to $2|f(a)|>0$. If $f(a)=0$, differentiability gives $f(a+h)=h(f'(a)+o(1))$, hence $|f(a+h)|/h$ tends to $|f'(a)|$ from the right and $-|f'(a)|$ from the left. Equality of these limits is precisely $f'(a)=0$.
\end{proof}

\KEntry{K03.04.noncontinuous_derivative}{A discontinuous derivative}
\begin{example}
Set $f(0)=0$ and $f(x)=x^2\sin(1/x)$ for $x\ne0$. Then $f'(0)=0$, since $|f(h)/h|\le|h|\to0$. For $x\ne0$ the product and chain rules give
\[ f'(x)=2x\sin(1/x)-\cos(1/x). \]
At $x_n=1/(2\pi n)$ the derivative equals $-1$, so it does not tend to $f'(0)=0$. Thus $f$ is differentiable on $\R$ but is not $C^1$ at zero. The trigonometric identities used in this and the next example are supplied by module 07.
\end{example}

\KEntry{K03.04.unbounded_derivative}{An unbounded derivative}
\begin{example}
Set $g(0)=0$ and $g(x)=x^2\sin(1/x^2)$ for $x\ne0$. The bound $|g(h)/h|\le|h|$ proves $g'(0)=0$. Away from zero,
\[ g'(x)=2x\sin(1/x^2)-\frac2x\cos(1/x^2). \]
For $x_n=(2\pi n)^{-1/2}$, this equals $-2\sqrt{2\pi n}$ and is unbounded in every neighborhood of zero. Proper Riemann integrability on a compact interval requires boundedness, by \KRef{K03.08.integral}; consequently this derivative is not properly Riemann integrable on any compact interval containing zero in its interior. Differentiability alone therefore supplies no such integrability condition.
\end{example}

\KEntry{K03.04.geometric}{Tangents and velocities}
\begin{proposition}
At a differentiability point $a$, the affine function $\ell(x)=f(a)+f'(a)(x-a)$ satisfies $f(a+h)-\ell(a+h)=o(h)$. It is the unique affine function through $f(a)$ with this property. For a scalar path $f(t)$, the same derivative is its instantaneous velocity.
\end{proposition}
\begin{proof}
The approximation is \KRef{K03.04.peano}. If another slope $v$ gave the same error, subtraction would give $(f'(a)-v)h=o(h)$; division for $h\ne0$ forces $v=f'(a)$. The velocity statement is the definition as the limit of displacement divided by elapsed real time. At a point where this limit does not exist no tangent or velocity conclusion is supplied by this result.
\end{proof}

% BEGIN REVISED SOURCE chapters/M03.05.tex SHA256=80db8feab6cae7d1ebaeaac4e46dc8dc8a130ef637f8606391befed016efaa67
\chapter{Mean values, extrema, and convexity}
Exact mean-value equalities in this chapter concern real-valued functions. For complex-valued functions the conclusion is instead a norm estimate. Every use of a closed segment includes continuity at both endpoints.
\SourceRef{Zorich I, 5.3--5.4; Apostol, 5.9--5.12}

\KEntry{K03.05.fermat}{Fermat's necessary condition}
\begin{theorem}
If a real function has a local maximum or minimum at an interior point $a$ and has a finite derivative there, then $f'(a)=0$.
\end{theorem}
\begin{proof}
At a maximum, for sufficiently small $h>0$ the difference quotient is nonpositive, while for $h<0$ it is nonnegative. The common finite limit is therefore both nonpositive and nonnegative. Apply this argument to $-f$ at a minimum. Interior location and existence of the derivative matter: $f(x)=x$ has an endpoint minimum on $[0,1]$ with right derivative $1$, and $|x|$ has an interior minimum with no derivative.
\end{proof}

\KEntry{K03.05.rolle}{Rolle's theorem}
\begin{theorem}
For $a<b$, if $f:[a,b]\to\R$ is continuous, differentiable on $(a,b)$, and $f(a)=f(b)$, some $c\in(a,b)$ satisfies $f'(c)=0$.
\end{theorem}
\begin{proof}
The function attains both extreme values by \KRef{K03.03.extreme_values}. If constant, every interior point works. Otherwise some value differs from the common endpoint value, so either a maximum or a minimum is attained at an interior point. Fermat's theorem applies there.
\end{proof}

\KEntry{K03.05.lagrange}{The real mean value theorem}
\begin{theorem}
If real $f$ is continuous on $[a,b]$, $a<b$, and differentiable inside, then
\[ f(b)-f(a)=f'(c)(b-a)\quad\text{for some }c\in(a,b). \]
\end{theorem}
\begin{proof}
Apply Rolle's theorem to $F(t)=f(t)-f(a)-(f(b)-f(a))(t-a)/(b-a)$, whose endpoint values are zero. Its zero derivative at an interior point is exactly the required formula. For a complex counterexample take $f(t)=\cos t+i\sin t$ on $[0,2\pi]$ after module 07: the endpoint difference is zero but $f'(t)=-\sin t+i\cos t$ always has modulus one. Thus the scalar equality cannot be asserted for this target.
\end{proof}

\KEntry{K03.05.cauchy}{Cauchy's mean value theorem}
\begin{theorem}
For $a<b$, if real $f,g$ are continuous on $[a,b]$ and differentiable on $(a,b)$, then some $c\in(a,b)$ satisfies
\[ (f(b)-f(a))g'(c)=(g(b)-g(a))f'(c). \]
If $g'$ is nowhere zero inside, then $g(b)\ne g(a)$ and this gives
\[ \frac{f(b)-f(a)}{g(b)-g(a)}=\frac{f'(c)}{g'(c)}. \]
\end{theorem}
\begin{proof}
The function $H(t)=(f(b)-f(a))g(t)-(g(b)-g(a))f(t)$ has equal endpoint values; Rolle gives $H'(c)=0$. If $g(b)=g(a)$, Rolle applied to $g$ would give a zero derivative, contradicting the additional hypothesis. Both denominators in the ratio formula are then nonzero.
\end{proof}

\KEntry{K03.05.mean_inequality}{The scalar and complex mean inequality}
\begin{theorem}
Let $f:[x,y]\to\K$ be continuous, differentiable inside, and satisfy $|f'(t)|\le M$ for all interior $t$, where $x<y$ and $M\ge0$. Then
\[ |f(y)-f(x)|\le M(y-x). \]
\end{theorem}
\begin{proof}
If $v=f(y)-f(x)=0$ the result is immediate. Otherwise set $u=v/|v|$ and $F(t)=\RePart(\bar u f(t))$. For real scalars this is simply a choice of sign. Then $F$ is real, $F(y)-F(x)=|v|$, and $|F'(t)|\le|f'(t)|\le M$. The real mean value theorem gives $|v|=F'(c)(y-x)\le M(y-x)$. A fixed projection was used throughout, without requiring a complex derivative vector to equal the secant vector at any point.
\end{proof}

\KEntry{K03.05.monotonicity}{Derivatives and monotonicity}
\begin{theorem}
Let real $f$ be continuous on an interval $I$ and differentiable at every interior point. If $f'\ge0$ inside then $f$ is nondecreasing; if $f'>0$ inside it is strictly increasing. Reversing the derivative signs gives decreasing versions. A nondecreasing differentiable function necessarily has $f'\ge0$, but strict increase does not force $f'>0$ everywhere.
\end{theorem}
\begin{proof}
For any $x<y$ in $I$, apply \KRef{K03.05.lagrange} on $[x,y]$ to obtain $f(y)-f(x)=f'(c)(y-x)$ and the required sign. For necessity, all difference quotients of a nondecreasing function at an interior point are nonnegative, including negative increments, so their limit is nonnegative. The strictly increasing function $x^3$ has derivative zero at zero. Endpoint continuity cannot be discarded: on $[0,1]$ set $f(0)=2$, $f(x)=x$ for $x>0$. Its interior derivative is positive, but it is not nondecreasing on the stated closed interval.
\end{proof}

\KEntry{K03.05.zero_derivative}{Zero derivatives and constants}
\begin{corollary}
A scalar function continuous on an interval and differentiable inside, with derivative zero at every interior point, is constant. Two such functions with the same derivative differ by a constant. On a domain consisting of several interval components the constants may differ between components.
\end{corollary}
\begin{proof}
Apply \KRef{K03.05.mean_inequality} with $M=0$ to any segment in the interval. Apply this to $f-g$ for equal derivatives. On $(-2,-1)\cup(1,2)$ the function equal to zero on the first component and one on the second has zero derivative everywhere on its domain but is not globally constant.
\end{proof}

\KEntry{K03.05.root_uniqueness}{Existence and uniqueness of roots}
\begin{proposition}
If real $f$ is continuous on $[a,b]$ and $f(a)f(b)<0$, it has a root in $(a,b)$. If it is also strictly monotone, that root is unique. In particular interior differentiability with $f'>0$ or $f'<0$ suffices for uniqueness.
\end{proposition}
\begin{proof}
Zero is strictly between the endpoint values, so \KRef{K03.03.intermediate_values} gives an interior root. A strictly monotone function cannot take the value zero at two different points. The derivative condition implies strict monotonicity by the preceding theorem. Existence alone is insufficient: $f(x)=x(x^2-1)$ has roots $-1,0,1$ while its values at $-2$ and $2$ have opposite signs.
\end{proof}

\KEntry{K03.05.derivative_sign_test}{The first derivative sign test}
\begin{proposition}
Suppose real $f$ is continuous at interior $a$, differentiable on a punctured neighborhood, and $f'>0$ just to the left of $a$, $f'<0$ just to the right. Then $a$ is a strict local maximum. Reversed signs yield a strict local minimum.
\end{proposition}
\begin{proof}
On a short segment $[x,a]$ with $x<a$, endpoint continuity and interior differentiability give $f(a)-f(x)=f'(c)(a-x)>0$. On $[a,x]$ to the right the corresponding difference $f(x)-f(a)$ is negative. Negation gives the minimum test. A stationary point alone is insufficient: $f(x)=x^3$ has $f'(0)=0$ but increases through zero.
\end{proof}

\KEntry{K03.05.second_test}{The second derivative test}
\begin{theorem}
Suppose $f$ is differentiable in a neighborhood of $a$, $f'$ is differentiable at $a$, and $f'(a)=0$. If $f''(a)>0$ then $a$ is a strict local minimum; if $f''(a)<0$ it is a strict local maximum.
\end{theorem}
\begin{proof}
When $f''(a)>0$, the quotient $(f'(a+h)-f'(a))/h$ is positive for all sufficiently small $h\ne0$. Thus $f'$ is negative to the left and positive to the right; use \KRef{K03.05.derivative_sign_test}. Negation gives the other case. If $f''(a)=0$ there is no conclusion: $x^4$, $-x^4$, and $x^3$ at zero give a minimum, a maximum, and neither, respectively.
\end{proof}

\KEntry{K03.05.darboux}{The intermediate value property of derivatives}
\begin{theorem}
If real $f$ is differentiable on an open interval and $u<v$ are in it, then $f'$ takes every value between $f'(u)$ and $f'(v)$ at a point of $[u,v]$, and every strictly intermediate value at a point of $(u,v)$.
\end{theorem}
\begin{proof}
Suppose $f'(u)<\lambda<f'(v)$ and put $F(t)=f(t)-\lambda t$. Its derivative at $u$ is negative, so a point just to the right has value less than $F(u)$. Its derivative at $v$ is positive, so a point just to the left has value less than $F(v)$. The continuous function $F$ attains a minimum on $[u,v]$, and neither endpoint can be a minimizer. Fermat gives $F'(c)=0$, hence $f'(c)=\lambda$, at an interior point. For reversed endpoint derivative order apply this argument to $-f$; endpoint values themselves are already attained. No continuity of $f'$ was assumed. In particular a nowhere-zero derivative on an interval has a constant sign, since opposite signs would force a zero.
\end{proof}

\KEntry{K03.05.lhopital_zero}{L'Hopital's rule with zero limits}
\begin{theorem}
Fix a one-sided endpoint approach along an open interval $J$, including an approach to $+\infty$ or $-\infty$. Let real $f,g$ be differentiable on $J$, let $g'$ be nowhere zero there, and suppose $f,g\to0$. If $f'/g'\to L\in\overline\R$ along that approach, then $g$ is eventually nonzero and $f/g\to L$.
\end{theorem}
\begin{proof}
By Darboux, $g'$ has a constant sign, so $g$ is strictly monotone. A strictly monotone function tending to zero at the selected endpoint cannot vanish at an interior point: its subsequent values in the direction of approach would all lie strictly on one side of zero and be bounded away from zero. Thus $g\ne0$ on $J$.
Suppose $L$ is finite. Restrict to an approach neighborhood where $|f'/g'-L|<\eps$. For two distinct points $x,t$ there, Cauchy's theorem gives
\[ \left|\frac{f(x)-f(t)}{g(x)-g(t)}-L\right|<\eps. \]
The denominator is nonzero by strict monotonicity. Keep $x$ fixed and let $t$ tend to the endpoint within that neighborhood. Since $f(t),g(t)\to0$ and $g(x)\ne0$, it follows that $|f(x)/g(x)-L|\le\eps$. Choosing $\eps$ smaller than any prescribed output error proves the finite limit. If $L=+\infty$, choose instead a neighborhood with $f'/g'>A$ for any fixed real $A$. The same mean value argument and limit give $f(x)/g(x)\ge A$. Since $A$ is arbitrary the quotient tends to $+\infty$. Reverse inequalities treat $-\infty$.
\end{proof}

\KEntry{K03.05.lhopital_infinite}{L'Hopital's rule with an unbounded denominator}
\begin{theorem}
Under the same specified endpoint and real differentiability conditions, suppose $g'$ is nowhere zero, $|g|\to\infty$, and $f'/g'\to L\in\overline\R$. Then $f/g\to L$. No limit assumption on $f$ is needed.
\end{theorem}
\begin{proof}
Eventually $g\ne0$. For finite $L$, restrict to a neighborhood where the derivative ratio differs from $L$ by $<\eps$, and fix one point $t$ in that neighborhood. Cauchy's theorem on the segment joining $t$ to a point $x$ closer to the endpoint gives
\[ \left|\frac{f(x)-f(t)}{g(x)-g(t)}-L\right|<\eps. \]
Here $g(x)\ne g(t)$ by Rolle's theorem and $g'\ne0$. Hence
\[ \left|\frac{f(x)}{g(x)}-L\right|\le\eps\left|1-\frac{g(t)}{g(x)}\right|+\frac{|f(t)-Lg(t)|}{|g(x)|}. \]
The second term tends to zero and the factor of $\eps$ tends to one. Arbitrarily small $\eps$ proves the result. If $L=+\infty$, the corresponding increment ratio is greater than any chosen $A$ once $t,x$ lie sufficiently far into the approach. The identity
\[ \frac{f(x)}{g(x)}=\frac{f(x)-f(t)}{g(x)-g(t)}\left(1-\frac{g(t)}{g(x)}\right)+\frac{f(t)}{g(x)} \]
has a positive factor tending to one. Thus the lower limit of $f/g$ is at least $A$ for every real $A$ (or use a larger positive $A$ for any desired threshold), giving $+\infty$. Negating $f$ treats $-\infty$.
\end{proof}

\KEntry{K03.05.convexity}{Convexity and secant slopes}
\begin{definition}
A real function on an interval is convex if
\[ f((1-t)x+ty)\le(1-t)f(x)+tf(y)\quad(x,y\in I,\ 0\le t\le1). \]
It is strictly convex if the inequality is strict whenever $x\ne y$ and $0<t<1$. Concavity means convexity of $-f$.
\end{definition}
\begin{proposition}
Convexity is equivalent to the ordered secant inequalities, for every $x<y<z$ in the interval,
\[ \frac{f(y)-f(x)}{y-x}\le\frac{f(z)-f(x)}{z-x}\le\frac{f(z)-f(y)}{z-y}. \]
\end{proposition}
\begin{proof}
Write $y=(1-t)x+tz$ with $t=(y-x)/(z-x)$. Subtracting $f(x)$ from the convexity inequality and dividing by $y-x>0$ yields the first secant inequality; subtracting instead from $f(z)$ yields the second. Conversely the first secant inequality rearranges to the convexity inequality for each strict interior combination. The cases $t=0,1$ and $x=y$ are equalities.
\end{proof}

\KEntry{K03.05.convex_derivative}{Convexity and derivatives}
\begin{theorem}
For differentiable real $f$ on an open interval, convexity is equivalent to $f'$ being nondecreasing, and equivalent to every tangent supporting the graph:
\[ f(y)\ge f(x)+f'(x)(y-x)\quad(x,y\in I). \]
If $f$ is twice differentiable, convexity is equivalent to $f''\ge0$.
\end{theorem}
\begin{proof}
For convex $f$, let the near endpoint of a secant tend to $x$ or $y$ in the inequalities above. For $x<y$ this gives $f'(x)\le(f(y)-f(x))/(y-x)\le f'(y)$. Thus $f'$ is nondecreasing and, considering both orders of $x,y$, the tangent inequality holds. Conversely if $f'$ is nondecreasing, the mean value theorem on $[x,y]$ gives the tangent inequality: when $y>x$ its intermediate derivative is at least $f'(x)$, and when $y<x$ it is at most $f'(x)$ and the negative factor reverses the inequality. If every tangent supports, apply the two inequalities based at $z=(1-t)x+ty$, multiply by $1-t$ and $t$, and add; the linear terms cancel, giving convexity. Finally $f''\ge0$ makes $f'$ nondecreasing by \KRef{K03.05.monotonicity}. Conversely a differentiable nondecreasing $f'$ has derivative $f''\ge0$ by the necessity part of that result.
\end{proof}

\KEntry{K03.05.convex_onesided}{One-sided derivatives of finite convex functions}
\begin{theorem}
A finite convex function on an open interval is continuous and has finite one-sided derivatives at every interior point, with $f'_-(x)\le f'_+(x)$. Every $p\in[f'_-(x),f'_+(x)]$ is a supporting slope:
\[ f(y)\ge f(x)+p(y-x)\quad(y\in I). \]
\end{theorem}
\begin{proof}
Fix $u<u_1<v_1<v$ in the interval. For $u_1\le x<y\le v_1$, the secant inequalities repeatedly give
\[ \frac{f(u_1)-f(u)}{u_1-u}\le\frac{f(y)-f(x)}{y-x}\le\frac{f(v)-f(v_1)}{v-v_1}. \]
Thus $f$ is Lipschitz on $[u_1,v_1]$, with constant the maximum of the moduli of these two fixed finite slopes. Such segments cover the interval, proving continuity. At fixed $x$, right secants $(f(x+h)-f(x))/h$ are nondecreasing as $h>0$ increases, so their limit as $h\downarrow0$ is their infimum; nearby external secants give finite upper and lower bounds. Left secants have a finite supremum giving $f'_-(x)$. Every left secant is at most every right secant by convexity, giving $f'_-(x)\le f'_+(x)$. For $y>x$, its secant is at least $f'_+(x)\ge p$. For $y<x$, its secant toward $x$ is at most $f'_-(x)\le p$; multiplication by $y-x<0$ gives the same supporting inequality. These arguments also identify all supporting slopes, since their one-sided limits force $f'_-(x)\le p\le f'_+(x)$.
\end{proof}

\KEntry{K03.05.jensen}{Finite Jensen inequality}
\begin{theorem}
If $f$ is convex, $w_j\ge0$, $\sum_{j=1}^m w_j=1$, and $x_j\in I$, then
\[ f\left(\sum_{j=1}^m w_jx_j\right)\le\sum_{j=1}^m w_jf(x_j). \]
If $f$ is strictly convex, equality holds precisely when all $x_j$ with positive weights are equal.
\end{theorem}
\begin{proof}
Discard zero weights. The one-weight case is equality. For $m>1$ write the average as $(1-w_m)y+w_mx_m$, where $y=(\sum_{j<m}w_jx_j)/(1-w_m)$. Convexity followed by the induction hypothesis for $y$ proves the inequality. If all positively weighted inputs agree, equality is immediate. If two such inputs are distinct, combine just those two into their weighted average. Strict convexity makes their contribution strictly larger than the function at that combined average; Jensen applied to the remaining list then gives a strict total inequality. This proves the converse equality statement without any unstated condition on zero weights.
\end{proof}

\KEntry{K03.05.inequalities}{Exponential, logarithmic, and arithmetic--geometric inequalities}
\begin{proposition}
After module 07, $\exp x\ge1+x$ for real $x$, $\log x\le x-1$ for $x>0$, and, for positive $a_j$ and weights $w_j\ge0$ summing to one,
\[ \prod_{j=1}^m a_j^{w_j}\le\sum_{j=1}^m w_ja_j. \]
Equality in the first two inequalities occurs at $x=0$ and $x=1$, respectively; equality in the weighted inequality occurs exactly when all positively weighted $a_j$ coincide.
\end{proposition}
\begin{proof}
The elementary formulas give $(\exp x)''=\exp x>0$, so its tangent at zero supports the graph and gives $\exp x\ge1+x$. Strictness for $x\ne0$ follows because the derivative $\exp x$ is strictly increasing: the mean value secant from zero has slope strictly above $1$ to the right and strictly below $1$ to the left. Substitute $x=\log u$ to obtain $u\ge1+\log u$, with equality only at $u=1$. Also $(\log x)''=-1/x^2<0$, so $-\log$ is strictly convex, using the same strict secant argument. Jensen for $-\log$ gives $\log(\sum w_ja_j)\ge\sum w_j\log a_j$. Apply the increasing exponential and its addition and real-power identities to get the displayed inequality. Strict Jensen gives the equality condition.
\end{proof}

% BEGIN REVISED SOURCE chapters/M03.11.tex SHA256=6ecb6ab761587512173051ff8653fdb882943fe1a79e879d5900a19a9fcf4539
\chapter{Numerical series, products, and rearrangements}
\SourceRef{Rudin, Chapter 3 (including Theorems 3.33--3.55); Apostol, Chapter 8; Tao, Analysis I, Chapter 7. The proofs below use completeness and integer roots, not elementary transcendental functions. The explicitly identified logarithmic applications occur only after Chapter 7's constructions.}

\KEntry{K03.11.partial_sums}{Series and finite changes}
\begin{definition}
For scalars $a_n\in\R$ or $\C$, indexed by $n\geq0$, put $s_N=\sum_{n=0}^N a_n$. The equation $\sum_{n=0}^\infty a_n=S$ means $s_N\to S$ with $S$ finite. A divergent series has no such scalar sum.
\end{definition}
\begin{proof}
If $b_n=a_n$ for $n>m$, then for $N\geq m$, $\sum_{n=0}^N b_n=s_N+\sum_{n=0}^m(b_n-a_n)$. Thus either both series converge or neither does; their sums differ by the displayed finite correction. The limit laws similarly prove linearity of convergent series.
\end{proof}

\KEntry{K03.11.necessary}{Term test and Cauchy criterion}
\begin{theorem}
A series converges if and only if for every $\eps>0$ there is $N$ such that $\abs{\sum_{n=p}^q a_n}<\eps$ whenever $q\geq p\geq N$. Convergence implies $a_n\to0$, whose converse is false.
\end{theorem}
\begin{proof}
The finite tail is $s_q-s_{p-1}$, so the assertion is precisely the Cauchy criterion in the complete scalar field (\KRef{K03.01.cauchy}). Also $a_n=s_n-s_{n-1}\to0$. For the converse counterexample, $a_n=1/n$ tends to zero, whereas the terms from $2^k$ through $2^{k+1}-1$ sum to at least $1/2$; its partial sums are unbounded.
\end{proof}

\KEntry{K03.11.nonnegative}{Nonnegative sums and regrouping}
\begin{theorem}
For $a_n\geq0$, the series converges exactly when its partial sums are bounded. Otherwise write $\sum a_n=+\infty$. Its extended sum is $\sup\{\sum_{n\in F}a_n:F\subset\N_0\text{ finite}\}$, and is unchanged by any ordering or by a partition of the indices into finite or countable blocks, each summed nonnegatively.
\end{theorem}
\begin{proof}
The partial sums increase, so completeness gives their limit, or unboundedness gives divergence to $+\infty$. Every finite set lies in an initial segment and each initial segment is a finite set, proving the supremum formula. Every finite selection within finitely many blocks is a finite selection of original terms, giving one inequality for iterated sums. Conversely a finite set of original indices meets only finitely many blocks, whose sums dominate its sum. Taking suprema proves the reverse inequality, also when the supremum is infinite.
\end{proof}

\KEntry{K03.11.comparison}{Comparison tests}
\begin{theorem}
For eventually nonnegative $a_n,b_n$, $a_n\leq Cb_n$ eventually implies convergence of $\sum a_n$ if $\sum b_n$ converges. The equivalent divergence implication follows by contraposition. If $b_n>0$ eventually and $a_n/b_n\to L\in(0,\infty)$, the series converge together. If $L=0$, convergence of $\sum b_n$ implies that of $\sum a_n$; if $L=+\infty$, convergence of $\sum a_n$ implies that of $\sum b_n$.
\end{theorem}
\begin{proof}
Remove a finite prefix and compare partial sums. For $0<L<\infty$ use $L/2\leq a_n/b_n\leq2L$ eventually. For $L=0$ use $a_n\leq b_n$ eventually; for $L=\infty$ use $b_n\leq a_n$. No reverse implication in these last cases follows: $a_n=1/n^2,b_n=1/n$ gives a zero ratio and unequal behavior; $a_n=1/n,b_n=1/n^2$ gives an infinite ratio and unequal behavior.
\end{proof}

\KEntry{K03.11.geometric_p}{Geometric, harmonic, and p-series}
\begin{example}
The geometric series $\sum_{n\geq0}q^n$ has sum $(1-q)^{-1}$ for $\abs q<1$ and diverges for $\abs q\geq1$. The harmonic series diverges. For real $p$, $\sum_{n\geq1}n^{-p}$ converges exactly when $p>1$.
\end{example}
\begin{proof}
The finite identity $(1-q)\sum_{n=0}^Nq^n=1-q^{N+1}$ proves the geometric assertion: $\abs q^n\to0$ for $\abs q<1$ by monotone convergence (a positive limit would equal $\abs q$ times itself); otherwise the term test applies. Harmonic divergence was proved in \KRef{K03.11.necessary}. For $p\leq0$, terms do not tend to zero. For $0<p\leq1$, compare with $1/n$. For $p>1$, condensation, proved below without using this example, gives the geometric series $\sum_k2^{k(1-p)}$, with ratio $2^{1-p}<1$. Here the meaning and order laws of arbitrary real powers are the downstream interface \KRef{K03.07.real_powers}; for positive integral $p$ this proof needs only arithmetic. Nothing in the root or general power-series test uses this real-power application.
\end{proof}

\KEntry{K03.11.ratio}{Ratio test}
\begin{theorem}
For eventually nonzero terms, if $\limsup\abs{a_{n+1}/a_n}<1$, then $\sum a_n$ converges absolutely. If $\abs{a_{n+1}/a_n}\geq1$ eventually, the series diverges; in particular $\liminf\abs{a_{n+1}/a_n}>1$ suffices. A ratio limit $1$ is inconclusive.
\end{theorem}
\begin{proof}
Choose $q<1$ strictly above the limsup; for $n\geq N$, $\abs{a_n}\leq\abs{a_N}q^{n-N}$, whose sum is geometric. Absolute convergence implies convergence by \KRef{K03.11.absolute}, whose proof is independent of this test. In the second case the eventually nonzero magnitudes never decrease, so cannot tend to zero. The series $\sum1/n$ and $\sum1/n^2$ have ratio limit $1$ and different behavior (for the latter use condensation below).
\end{proof}

\KEntry{K03.11.root}{Root test without logarithms}
\begin{theorem}
Let $L=\limsup_{n\to\infty}\abs{a_n}^{1/n}\in[0,\infty]$, using positive integer roots from \KRef{K03.01.positive_roots}. If $L<1$, the series converges absolutely. If $L>1$, it diverges. When $L=1$, no conclusion follows.
\end{theorem}
\begin{proof}
For $L<q<1$, eventually $\abs{a_n}^{1/n}\leq q$, hence $\abs{a_n}\leq q^n$. The finite identity $\sum_{n=p}^m t^n=t^p(1-t^{m-p+1})/(1-t)$ for $0\leq t<1$, with $t^p\to0$ by monotone convergence, makes the comparison tails arbitrarily small. Absolute convergence follows directly from the numerical Cauchy criterion. For $1<q<L$, infinitely many $n$ satisfy $\abs{a_n}>q^n\geq1$, so the term test fails. For inconclusiveness compare $1/n$ and $1/n^2$: their roots tend to $1$. Indeed, for any $h>0$, the binomial theorem gives $(1+h)^n\geq\binom n3 h^3>n^2$ eventually, while $1\leq n^{2/n}$; squeeze and reciprocation prove the claim. No exponential or logarithm has been used.
\end{proof}

\KEntry{K03.11.condensation}{Cauchy condensation}
\begin{theorem}
If $a_n\geq0$ is nonincreasing for $n\geq1$, then $\sum a_n$ converges if and only if $\sum_{k\geq0}2^k a_{2^k}$ converges.
\end{theorem}
\begin{proof}
On a dyadic block,
\[2^k a_{2^{k+1}}\leq\sum_{n=2^k}^{2^{k+1}-1}a_n\leq2^k a_{2^k}.\]
Summing these inequalities compares block partial sums with the condensed partial sums, the left bound being one half of the condensed sum with its first term removed. Boundedness of one sequence of partial sums is therefore equivalent to boundedness of the other. Nonnegativity controls all intermediate partial sums.
\end{proof}

\KEntry{K03.11.alternating}{Alternating series and remainders}
\begin{theorem}
If $b_n\geq0$, $b_{n+1}\leq b_n$, and $b_n\to0$, then $\sum_{n\geq0}(-1)^nb_n$ converges. The remainder after $n=N$ has sign $(-1)^{N+1}$ or is zero, and magnitude at most $b_{N+1}$.
\end{theorem}
\begin{proof}
The partial sums ending at odd indices increase because $b_{2j}-b_{2j+1}\geq0$, while those ending at even indices decrease because $-b_{2j+1}+b_{2j+2}\leq0$. Each odd partial sum is below each later even partial sum, and the gap between adjacent such sums is a term $b_n\to0$. The two bounded monotone sequences thus have the same limit. Applying the same bracketing to the tail places its signed value between $0$ and $b_{N+1}$.
\end{proof}

\KEntry{K03.11.summation_parts}{Finite summation by parts}
\begin{theorem}
For $A_n=\sum_{j=0}^n a_j$, $A_{-1}=0$, and $0\leq p\leq q$,
\[\sum_{n=p}^q a_nb_n=A_qb_q-A_{p-1}b_p+\sum_{n=p}^{q-1}A_n(b_n-b_{n+1}).\]
Equivalently, with $T_n=\sum_{j=p}^n a_j$, the right side is $T_qb_q+\sum_{n=p}^{q-1}T_n(b_n-b_{n+1})$.
\end{theorem}
\begin{proof}
Insert $a_n=A_n-A_{n-1}$ and shift the index in the second finite sum. The uncancelled endpoints are exactly $A_qb_q$ and $-A_{p-1}b_p$. Replacing $A_n$ by $T_n$ gives the second formula. Limits may be taken only after the resulting boundary and series terms have been controlled.
\end{proof}

\KEntry{K03.11.dirichlet}{Dirichlet's test}
\begin{theorem}
If partial sums of scalar $a_n$ are bounded and real $b_n$ is monotone with limit zero, then $\sum a_nb_n$ converges.
\end{theorem}
\begin{proof}
Suppose first $b_n\downarrow0$ and $\abs{A_n}\leq M$. Then $\abs{T_n}\leq2M$, and summation by parts bounds every tail by
\[2M\left(b_q+\sum_{n=p}^{q-1}(b_n-b_{n+1})\right)=2Mb_p\longrightarrow0.\]
Use the Cauchy criterion. An increasing sequence tending to zero is nonpositive; replace it by its negative.
\end{proof}

\KEntry{K03.11.abel}{Abel's test}
\begin{theorem}
If $\sum a_n$ converges and real $b_n$ is monotone and bounded, then $\sum a_nb_n$ converges.
\end{theorem}
\begin{proof}
Monotone convergence gives $b_n\to b$. Split $a_nb_n=ba_n+a_n(b_n-b)$. The second factor in the latter product is monotone and tends to zero; the partial sums of $a_n$ are bounded, so Dirichlet applies.
\end{proof}

\KEntry{K03.11.absolute}{Absolute and conditional convergence}
\begin{theorem}
If $\sum\abs{a_n}<\infty$, then $\sum a_n$ converges, and every convergent tail satisfies $\abs{\sum_{n=p}^\infty a_n}\leq\sum_{n=p}^\infty\abs{a_n}$. Convergence without absolute convergence is called conditional.
\end{theorem}
\begin{proof}
For finite tails the inequality is the triangle inequality. Its right side tends uniformly to zero as $p\to\infty$, so the scalar Cauchy criterion gives convergence. Passing $q\to\infty$ in the finite inequality gives the displayed bound. The alternating harmonic series is conditional by the alternating test and harmonic divergence.
\end{proof}

\KEntry{K03.11.rearrange_absolute}{Grouping and absolute rearrangements}
\begin{theorem}
Consecutive finite grouping of a convergent series preserves its sum. Every bijective rearrangement of an absolutely convergent scalar series is absolutely convergent with the same sum.
\end{theorem}
\begin{proof}
Grouped partial sums are a subsequence of the original partial sums. For a permutation $\pi$, choose $m$ so $\sum_{n>m}\abs{a_n}<\eps$. Once a rearranged partial sum contains all indices $0,\ldots,m$, its difference from the original sum has magnitude at most $\sum_{n>m}\abs{a_n}<\eps$: subtract the finite included terms and take limits in the remaining absolutely convergent tail. The nonnegative regrouping theorem applied to $\abs{a_n}$ proves absolute convergence of the rearrangement.
\end{proof}

\KEntry{K03.11.rearrange_conditional}{Riemann rearrangement theorem}
\begin{theorem}
For a conditionally convergent real series, the positive terms sum to $+\infty$ and the negative terms sum to $-\infty$. A permutation can have any prescribed finite sum, can tend to either infinity, or can have any prescribed extended lower and upper limits $\alpha\leq\beta$.
\end{theorem}
\begin{proof}
Write $a_n=p_n-q_n$, $p_n=\max(a_n,0)$, $q_n=\max(-a_n,0)$. If one of $\sum p_n,\sum q_n$ were finite, convergence of $\sum a_n$ would make the other finite, contradicting $\sum\abs{a_n}=\infty$. Enumerate positive and negative terms in their original orders; both enumerations tend to zero.

For a finite target $s$, take unused positive terms until the partial sum first exceeds $s$, then unused negative terms until it first falls below $s$, and repeat. Each crossing exists by the infinite one-sided sums, and each overshoot is at most the last term used, tending to zero. All intermediate sums lie between the neighboring crossing sums up to these vanishing overshoots; hence they tend to $s$. Both lists are exhausted because infinitely many blocks of each sign are used. Insert zero terms successively, without affecting any sum.

For $+\infty$, at stage $j$ use the next negative term and then at least the next positive term, continuing positively until the sum exceeds $j$. The negative terms tend to zero, so even the minima of successive stages tend to $+\infty$; each list is exhausted. Reverse signs for $-\infty$.

For $\alpha<\beta$, choose finite $u_j<v_j$ with $u_j\to\alpha$, $v_j\to\beta$ in the extended sense. Alternately take at least one unused positive term and continue until above $v_j$, then at least one unused negative term and continue until below $u_j$. The crossing errors tend to zero; all intermediate values lie between neighboring endpoints, so the lower and upper limits are $\alpha,\beta$. The same proof covers infinite endpoints. When $\alpha=\beta$ use the finite or infinite constructions already given. This constructs the permutation, not merely a grouping.
\end{proof}

\KEntry{K03.11.cauchy_product}{Cauchy products and Mertens' theorem}
\begin{theorem}
Put $c_n=\sum_{k=0}^n a_kb_{n-k}$. If $\sum a_n=A$ absolutely and $\sum b_n=B$ converges, then $\sum c_n=AB$. If both factors are absolutely convergent, so is their Cauchy product.
\end{theorem}
\begin{proof}
Writing $B_j=\sum_{k=0}^j b_k$ gives $C_N=\sum_{n=0}^Nc_n=\sum_{k=0}^N a_kB_{N-k}$. Thus
\[C_N-AB=\sum_{k=0}^Na_k(B_{N-k}-B)-B\sum_{k>N}a_k.\]
The second term tends to zero. For the first, $\abs{B_j-B}\leq M$; choose $m$ with $M\sum_{k>m}\abs{a_k}<\eps$, and then make each of the finitely many factors $B_{N-k}-B$, $k\leq m$, small. This proves convergence. In the absolute case,
\[\sum_{n=0}^N\abs{c_n}\leq\sum_{k+j\leq N}\abs{a_k}\abs{b_j}\leq\Big(\sum_k\abs{a_k}\Big)\Big(\sum_j\abs{b_j}\Big),\]
so the product is absolutely convergent.
\end{proof}

\KEntry{K03.11.product_failure}{Two conditional factors can have a divergent product}
\begin{example}
Let $a_n=b_n=(-1)^n/\sqrt{n+1}$, $n\geq0$. Each series converges by the alternating test. Its Cauchy product has
\[\abs{c_n}=\sum_{k=0}^n\frac1{\sqrt{(k+1)(n-k+1)}}\geq\frac{2(n+1)}{n+2}\geq1,\]
because $(k+1)(n-k+1)\leq(n+2)^2/4$. Thus $c_n\not\to0$ and the product diverges. The factors are conditional, since $1/\sqrt{n+1}\geq1/(n+1)$.
\end{example}

\KEntry{K03.11.double_series}{Double sums and summation order}
\begin{theorem}
For nonnegative $a_{ij}$, all iterated sums and every enumeration of $\N_0^2$ have the same extended sum, the supremum over finite subsets. If $\sum_{i,j}\abs{a_{ij}}<\infty$, all such scalar sums exist and agree.
\end{theorem}
\begin{proof}
The nonnegative proof is exactly the finite-subset argument in \KRef{K03.11.nonnegative}, with rows or columns as blocks. In the absolute case choose a finite rectangle outside which the sum of magnitudes is $<\eps$ (a finite subset achieving this is contained in a rectangle). Any two partial sums containing that rectangle differ by $<2\eps$, proving existence and agreement; each iterated tail obeys the same bound.
For failure without the hypothesis, set $a_{ij}=1$ if $j=i$, $a_{ij}=-1$ if $j=i+1$, and zero otherwise. Every row sums to zero, whereas the first column sums to $1$ and every later column to zero. Thus the two iterated sums are $0$ and $1$.
\end{proof}

\KEntry{K03.11.products}{Infinite products and logarithms}
\begin{definition}
For nonzero scalar factors $v_n$, convergence of $\prod_{n\geq0}v_n$ means that $P_N=\prod_{n=0}^Nv_n$ tends to a nonzero scalar. A zero factor, or a nonzero-factor product whose limit is zero, is recorded separately, not called a convergent nonzero product. A convergent product necessarily has $v_n=P_n/P_{n-1}\to1$.
\end{definition}
\begin{theorem}
For positive factors, convergence to a positive nonzero product is equivalent to convergence of $\sum\log v_n$.
\end{theorem}
\begin{proof}
This is a downstream logarithmic characterization using \KRef{K03.07.log_construction} and \KRef{K03.07.exp_addition}: $\log P_N=\sum_{n=0}^N\log v_n$. Continuity of log proves one implication and continuity and positivity of exp prove the other. Finite nonzero prefixes simply multiply the limit. The example $v_n=1/2$ has limit zero and shows why the nonzero convention matters.
\end{proof}

\KEntry{K03.11.product_absolute}{Absolutely convergent products without logarithms}
\begin{theorem}
If $u_n\in\C$, $\sum\abs{u_n}<\infty$, and $1+u_n\ne0$ for all $n$, then $\prod(1+u_n)$ converges to a nonzero value.
\end{theorem}
\begin{proof}
For finitely many $t_j\geq0$ with $\sum t_j=T<1$, expansion gives $\prod(1+t_j)\leq\sum_{k\geq0}T^k=(1-T)^{-1}$: each elementary symmetric sum is at most $T^k$. Also
\[\abs{\prod(1+u_j)-1}\leq\prod(1+\abs{u_j})-1\leq T/(1-T).\]
Choose a tail with total magnitude $<1/2$. Its partial products, and hence the full partial products, are uniformly bounded; the displayed inequality for later tails makes $P_N$ Cauchy. Finally the elementary inequality $\prod(1-t_j)\geq1-\sum t_j$ (proved inductively) gives the modulus of every such tail product at least $1-T>1/2$. Multiplication by the fixed nonzero prefix bounds $\abs{P_N}$ away from zero. Its limit is therefore nonzero.
\end{proof}

\KEntry{K03.11.cesaro}{Cesaro summation}
\begin{definition}
With $s_N=\sum_{n=0}^Na_n$, the Cesaro means are $\sigma_N=(N+1)^{-1}\sum_{j=0}^Ns_j$. A finite limit of these means is the Cesaro sum.
\end{definition}
\begin{proof}
If $s_j\to S$, split the average of $s_j-S$ at a fixed $m$: the finite prefix divided by $N+1$ tends to zero, and the remaining average is bounded by $\sup_{j>m}\abs{s_j-S}$. Thus $\sigma_N\to S$. For $a_n=(-1)^n$, $s_N$ alternates between $1$ and $0$, so $\sigma_N\to1/2$ although the series diverges. Summability and ordinary convergence are distinct.
\end{proof}

\KEntry{K03.11.acceleration}{Remainders and a proved telescoping correction}
\begin{proposition}
An alternating remainder is bounded by its first omitted term. For a nonnegative decreasing $f$ on $[1,\infty)$ with convergent improper integral,
\[\int_{N+1}^\infty f(x)\dd x\leq\sum_{n>N}f(n)\leq\int_N^\infty f(x)\dd x.\]
For the reciprocal-square sum, a concrete correction is
\[\sum_{n=1}^\infty\frac1{n^2}=2+\sum_{n=1}^\infty\left(\frac1{n^2}-\frac1{n-1/2}+\frac1{n+1/2}\right),\]
and after $N\ge 1$ corrected terms the remaining error has magnitude at most $1/(9N^3)$.
\end{proposition}
\begin{proof}
The first assertion is \KRef{K03.11.alternating}. The integral inequalities follow from $f(n+1)\leq f(x)\leq f(n)$ on $[n,n+1]$, summed on finite truncations before passing to the improper limit (\KRef{K03.10.integral_test}). For the correction,
\[\frac1{n-1/2}-\frac1{n+1/2}=\frac1{n^2-1/4},\qquad
\frac1{n^2}-\frac1{n^2-1/4}=-\frac1{4n^2(n^2-1/4)}.\]
The telescoping sum of the left difference through $N$ is $2-1/(N+1/2)\to2$. Since $n^2-1/4\geq3n^2/4$, the correction magnitude is at most $1/(3n^4)$. Its tail is at most $\frac13\int_N^\infty x^{-4}\dd x=1/(9N^3)$, proving both the identity and its useful error bound.
\end{proof}

% BEGIN REVISED SOURCE chapters/M03.12.tex SHA256=3ac52d57ce442c026eec52c1d9d5e1e4ac90b8ad30e4a9020ea9f7cf4fc3a25a
\chapter{Sequences and series of functions}
\SourceRef{Rudin, Chapter 7, especially Theorems 7.8--7.17; Apostol, Chapter 9; Zorich II, Sections 16.1--16.3. Scalars are real or complex; all derivatives below use real increments.}

\KEntry{K03.12.pointwise}{Pointwise convergence}
\begin{definition}
For functions $f_n:E\to\K$, $\K=\R$ or $\C$, pointwise convergence to $f:E\to\K$ means
\[\forall x\in E\ \forall\eps>0\ \exists N\ \forall n\geq N:\quad\abs{f_n(x)-f(x)}<\eps.\]
The index may depend on $x$. The pointwise limit is unique by uniqueness of scalar limits at each point. On the empty set the assertion is vacuous and the unique empty function is the limit.
\end{definition}

\KEntry{K03.12.uniform}{Uniform and locally uniform convergence}
\begin{definition}
Uniform convergence on $E$, written $f_n\to f$ uniformly, means
\[\forall\eps>0\ \exists N\ \forall n\geq N\ \forall x\in E:\quad\abs{f_n(x)-f(x)}<\eps.\]
It implies pointwise convergence by retaining any fixed $x$. On an open set $U\subset\R$, locally uniform convergence means uniform convergence on each compact subset of $U$. Uniform convergence on $U$ implies locally uniform convergence by restriction.
\end{definition}

\KEntry{K03.12.sup_norm}{The supremum criterion}
\begin{theorem}
For nonempty $E$, define $e_n=\sup_{x\in E}\abs{f_n(x)-f(x)}\in[0,\infty]$. Then uniform convergence is equivalent to $e_n\to0$, which in particular makes $e_n$ finite eventually. For $E=\varnothing$ set $e_n=0$ by convention. No boundedness of $f$ or of the individual $f_n$ is needed, only eventual boundedness of their differences.
\end{theorem}
\begin{proof}
If $e_n<\eps$, every pointwise error is $<\eps$. Conversely uniform errors $<\eps/2$ imply $e_n\leq\eps/2<\eps$. These implications prove the equivalence, including the stated convention for the empty set.
\end{proof}

\KEntry{K03.12.cauchy}{Uniform Cauchy criterion}
\begin{theorem}
A scalar function sequence on any set $E$ converges uniformly if and only if for each $\eps>0$ there is $N$ such that $\abs{f_n(x)-f_m(x)}<\eps$ for every $m,n\geq N$ and $x\in E$.
\end{theorem}
\begin{proof}
Uniform convergence gives the assertion by the triangle inequality with error $\eps/2$. Conversely, for each fixed $x$, scalar completeness produces $f(x)=\lim_m f_m(x)$. Apply the uniform Cauchy bound with $\eps/2$ and let $m\to\infty$ to obtain $\abs{f_n(x)-f(x)}\leq\eps/2<\eps$ for all $x$ and $n\geq N$. The same $N$ works everywhere.
\end{proof}

\KEntry{K03.12.continuity}{Continuity of uniform limits}
\begin{theorem}
A uniform limit on a metric domain of continuous scalar functions is continuous. On an open real domain locally uniform convergence also suffices.
\end{theorem}
\begin{proof}
Fix $x$ and $\eps>0$. Select $n$ with $\sup_E\abs{f-f_n}<\eps/3$, then a neighborhood of $x$ on which $\abs{f_n(y)-f_n(x)}<\eps/3$. The triangle inequality bounds $\abs{f(y)-f(x)}$ by the sum of these three errors. In the locally uniform case put a small compact interval around $x$ inside the open domain and apply this argument there. Uniform convergence is sufficient, not necessary: on $[0,1]$ the continuous triangles of \KRef{K03.12.integral_failure} converge pointwise to the continuous zero function, but their heights are unbounded.
\end{proof}

\KEntry{K03.12.dini}{Dini's theorem}
\begin{theorem}
Let $K$ be compact and $f_n,f:K\to\R$ continuous. If $f_n(x)\uparrow f(x)$ for every $x$, or $f_n(x)\downarrow f(x)$ for every $x$, then convergence is uniform.
\end{theorem}
\begin{proof}
For increasing convergence put $g_n=f-f_n\geq0$. For fixed $\eps>0$, the open sets $U_n=\{x:g_n(x)<\eps\}$ increase and cover $K$. A finite subcover, and hence its largest index $N$, covers $K$. Thus $0\leq g_n(x)<\eps$ for $n\geq N$ and every $x$. Replace $f_n,f$ by their negatives in the decreasing case. Continuity of the limit is essential: $x^n$ on $[0,1]$ decreases pointwise to the discontinuous function in \KRef{K03.12.pointwise_discontinuous} and is not uniformly convergent.
\end{proof}

\KEntry{K03.12.interchange_limits}{Interchanging a sequence limit and a domain limit}
\begin{theorem}
Let $E\subset\R$, let $c$ be an accumulation point approached through $E$, and let $f_n:E\to\K$ converge uniformly to $f$. Suppose $b_n=\lim_{x\to c,\,x\in E}f_n(x)$ exists for every $n$. Then $b_n$ converges to a scalar $b$, the domain limit of $f$ exists, and
\[\lim_{x\to c,\,x\in E}\lim_n f_n(x)=b=\lim_n\lim_{x\to c,\,x\in E}f_n(x).\]
The same proof applies to a specified one-sided or infinite approach; uniformity on a fixed approach neighborhood suffices.
\end{theorem}
\begin{proof}
Uniform Cauchy convergence gives $\abs{f_n(x)-f_m(x)}<\eps$ on $E$ for large $m,n$. Taking the domain limit gives $\abs{b_n-b_m}\leq\eps$, so $b_n\to b$ by completeness. Choose $n$ with $\sup_E\abs{f-f_n}<\eps/3$ and $\abs{b_n-b}<\eps/3$. Near $c$ in the specified approach, $\abs{f_n(x)-b_n}<\eps/3$. Their sum bounds $\abs{f(x)-b}$ and proves the required domain limit.
\end{proof}

\KEntry{K03.12.integration}{Uniform convergence and the Riemann integral}
\begin{theorem}
If $f_n$ are Riemann integrable on $[a,b]$ and converge uniformly to $f$, then $f$ is Riemann integrable, and
\[\abs{\int_a^b f-\int_a^b f_n}\leq(b-a)\sup_{[a,b]}\abs{f-f_n}\longrightarrow0.\]
\end{theorem}
\begin{proof}
For real functions, let $e_n=\sup\abs{f-f_n}$. For a partition $P$,
\[U(f,P)-L(f,P)\leq U(f_n,P)-L(f_n,P)+2e_n(b-a).\]
Choose $n$ making the second term $<\eps/2$, then a partition making the first $<\eps/2$. The Darboux criterion proves integrability of $f$. Apply the integral norm bound to $f-f_n$ to obtain the estimate. For complex functions apply the real argument to both components; the scalar integral norm bound of \KRef{K03.08.order} gives the same displayed constant. For $a=b$ all integrals are zero.
\end{proof}

\KEntry{K03.12.differentiation}{Uniform convergence of derivatives}
\begin{theorem}
Suppose $a<b$, $f_n:[a,b]\to\K$ are continuous and differentiable on $(a,b)$, $f_n(x_0)$ converges at one $x_0\in[a,b]$, and $f_n'\to g$ uniformly on $(a,b)$. Then $f_n$ converge uniformly to a continuous $f$, differentiable in $(a,b)$, with $f'=g$. Continuity of the individual derivatives is not assumed.
\end{theorem}
\begin{proof}
The norm mean value inequality \KRef{K03.05.mean_inequality} gives
\[\abs{(f_n-f_m)(x)}\leq\abs{(f_n-f_m)(x_0)}+(b-a)\sup_{(a,b)}\abs{f_n'-f_m'}.\]
Uniform Cauchy convergence follows, and its limit $f$ is continuous. Fix an interior $x$ and $\eta>0$. Choose $N$ with $\sup\abs{f_n'-g}<\eta$ for $n\geq N$. For $m,n\geq N$, the same mean value inequality on the interval from $x$ to $x+h$ yields
\[\abs{\frac{(f_m-f_n)(x+h)-(f_m-f_n)(x)}h}\leq2\eta.\]
Pass $m\to\infty$ at fixed nonzero $h$ to replace $f_m$ by $f$. Fix one $n\geq N$; differentiability of $f_n$ gives, for small $h$, a difference between its quotient and $f_n'(x)$ less than $\eta$. Since $\abs{f_n'(x)-g(x)}<\eta$, the quotient of $f$ differs from $g(x)$ by at most $4\eta$. Let $\eta\downarrow0$. This proves the derivative assertion. Endpoint derivatives require their own one-sided assumptions and are not included here.
\end{proof}

\KEntry{K03.12.series}{Three convergence conditions for function series}
\begin{definition}
The pointwise or uniform convergence of $\sum_{n\geq0}u_n$ means the corresponding convergence of $S_N=\sum_{n=0}^Nu_n$. Uniform convergence of $\sum\abs{u_n}$ is a stronger assertion; a summable bound $\sum\sup_E\abs{u_n}<\infty$ is stronger still.
\end{definition}
\begin{example}
The scalar alternating harmonic series, viewed as constant functions, converges uniformly but not absolutely. To separate the latter two conditions, let $E=\N$, and let $u_n(k)=1/n$ for $k=n\geq1$, zero otherwise. Then the uniform tail of $\sum\abs{u_n}$ is $1/(N+1)\to0$, but $\sum_n\sup_E\abs{u_n}=\sum1/n$ diverges. Thus the predicates cannot be identified.
\end{example}

\KEntry{K03.12.m_test}{Weierstrass' majorant test}
\begin{theorem}
If $\abs{u_n(x)}\leq M_n$ on $E$, $M_n\geq0$, and $\sum M_n<\infty$, then $\sum u_n$ and $\sum\abs{u_n}$ converge uniformly on $E$; the original series is absolutely convergent at every point.
\end{theorem}
\begin{proof}
Every finite tail of either series has magnitude at most $\sum_{n=p}^qM_n$, uniformly in $x$. The numerical Cauchy criterion makes this arbitrarily small. Apply the uniform Cauchy criterion, and the numerical comparison theorem at each point.
\end{proof}

\KEntry{K03.12.uniform_dirichlet}{Uniform Dirichlet test}
\begin{theorem}
Suppose $A_N(x)=\sum_{n=0}^Na_n(x)$ satisfies $\abs{A_N(x)}\leq M$ for all $x,N$. If real $b_n(x)$ is monotone in $n$ for each $x$ and tends uniformly to zero, then $\sum a_n(x)b_n(x)$ converges uniformly.
\end{theorem}
\begin{proof}
For each $x$ a real monotone sequence tending to zero has constant sign and variation on $[p,q]$ equal to $\abs{b_p(x)}-\abs{b_q(x)}$. In the tail form of summation by parts, $\abs{T_n(x)}\leq2M$. Therefore
\[\abs{\sum_{n=p}^q a_n(x)b_n(x)}\leq2M\abs{b_p(x)}.\]
The right side tends uniformly to zero. Apply the uniform Cauchy criterion. The direction of monotonicity may depend on $x$.
\end{proof}

\KEntry{K03.12.uniform_abel}{Uniform Abel test}
\begin{theorem}
If $\sum a_n(x)$ converges uniformly and real $b_n(x)$ is monotone in $n$ at each point, with $\abs{b_n(x)}\leq M$ for all $x,n$, then $\sum a_n(x)b_n(x)$ converges uniformly.
\end{theorem}
\begin{proof}
For large $p$, every $T_n(x)=\sum_{j=p}^na_j(x)$ has magnitude $<\eta$, uniformly for $n\geq p$ and $x$. Summation by parts bounds a product tail by
\[\eta\left(\abs{b_q(x)}+\sum_{n=p}^{q-1}\abs{b_n(x)-b_{n+1}(x)}\right)\leq3M\eta.\]
Use uniform Cauchy convergence; if $M=0$ the assertion is immediate. This proof does not require uniform convergence of $b_n$ to its pointwise limit.
\end{proof}

\KEntry{K03.12.termwise_integration}{Termwise integration}
\begin{theorem}
If each $u_n$ is Riemann integrable on $[a,b]$ and $\sum u_n$ converges uniformly there to $f$, then
\[\int_a^b f(x)\dd x=\sum_{n=0}^\infty\int_a^bu_n(x)\dd x.\]
\end{theorem}
\begin{proof}
Finite sums $S_N$ are integrable and their integrals are the corresponding finite sums of integrals by linearity. Apply \KRef{K03.12.integration} to $S_N\to f$. Its conclusion also proves convergence of the numerical series of integrals, rather than assuming it.
\end{proof}

\KEntry{K03.12.termwise_differentiation}{Termwise differentiation}
\begin{theorem}
Let $u_n$ be continuous on $[a,b]$ and differentiable inside. If $\sum u_n(x_0)$ converges for one $x_0\in[a,b]$ and $\sum u_n'$ converges uniformly on $(a,b)$, then $\sum u_n$ converges uniformly and its sum satisfies $f'(x)=\sum u_n'(x)$ inside.
\end{theorem}
\begin{proof}
Apply \KRef{K03.12.differentiation} to the finite partial sums. Their derivatives are exactly the finite sums of derivatives by the finite linearity rule, and their values at $x_0$ converge by hypothesis. All hypotheses of that theorem are thus verified.
\end{proof}

\KEntry{K03.12.pointwise_discontinuous}{A discontinuous pointwise limit}
\begin{example}
On $[0,1]$, $f_n(x)=x^n$ tends to $f(x)=0$ for $x<1$, $f(1)=1$. The geometric decay proved in \KRef{K03.11.geometric_p} proves the pointwise claim. For each fixed $n$, values $x<1$ arbitrarily close to $1$ make $x^n$ arbitrarily close to $1$, so $\sup_{[0,1]}\abs{f_n-f}=1$. Convergence is not uniform, and the limit is discontinuous at $1$.
\end{example}

\KEntry{K03.12.integral_failure}{Moving triangles and failure of integral interchange}
\begin{example}
For $n\geq2$, define on $[0,1]$
\[f_n(x)=\begin{cases}n^2x,&0\leq x\leq1/n,\\ n(2-nx),&1/n\leq x\leq2/n,\\0,&2/n\leq x\leq1.\end{cases}\]
The formulas agree at the joins, so $f_n$ is continuous. At $x=0$ it is zero; for fixed $x>0$ it is eventually zero because $2/n<x$. Thus $f_n\to0$ pointwise. Direct integration of the two affine pieces gives $1/2+1/2=1$ for every $n$, whereas the integral of the limit is zero. Also $f_n(1/n)=n$, so the convergence is not uniform.
\end{example}

\KEntry{K03.12.derivative_failure}{Uniform function convergence need not control derivatives}
\begin{example}
After the trigonometric construction in \KRef{K03.07.sin_cos}, take $f_n(x)=\sin(nx)/n$ on $\R$. Its supremum norm is at most $1/n$, so it converges uniformly to zero. But $f_n'(x)=\cos(nx)$, and at $x=\pi$ these values are $(-1)^n$, which do not converge. Thus the derivatives do not even converge pointwise everywhere.
\end{example}

\KEntry{K03.12.whole_line_failure}{Local uniformity need not be global}
\begin{example}
With exp constructed by \KRef{K03.07.exp_construction}, let $P_N(x)=\sum_{n=0}^Nx^n/n!$. For each $A>0$, the ratio of consecutive numerical bounds $A^n/n!$ is $A/(n+1)\to0$; the M-test proves uniform convergence on $[-A,A]$. On $x\geq0$ all omitted terms are nonnegative, so
\[\exp(x)-P_N(x)\geq x^{N+1}/(N+1)!,\]
which is unbounded as $x\to\infty$ for every fixed $N$. There is therefore no uniform convergence on $\R$. This is a downstream example; the M-test and general power-series theory do not use exp.
\end{example}

% BEGIN REVISED SOURCE chapters/M03.13.tex SHA256=86a9b0d911a485f6ab57dffd6da29923e14bfc9f55f4718676b39b96680fe2d4
\chapter{Power series and analytic functions}
\SourceRef{Apostol, Sections 9.14--9.19 (PDF 247--255); Rudin, Sections 3.38--3.40 and Chapter 8; Zorich II, Section 16.3 exercises on Abel and Tauber summation (PDF 410). General power-series arguments precede the elementary-function applications logically.}

\KEntry{K03.13.power_series}{Power series and their values}
\begin{definition}
For coefficients $a_n\in\K$, $\K=\R$ or $\C$, and fixed real center $c$, a power series is the expression $\sum_{n\geq0}a_n(x-c)^n$. At each argument its sum is defined only by convergence of its numerical partial sums. The same expression may be evaluated at complex arguments for estimates. A formal sequence of coefficients is not itself a function on all arguments, and a finite polynomial is distinguished from an infinite convergent sum.
\end{definition}

\KEntry{K03.13.radius}{The radius of convergence}
\begin{theorem}
Let $L=\limsup_{n\to\infty}\abs{a_n}^{1/n}$ and set $R=1/L$, with $1/0=\infty$, $1/\infty=0$. The series converges absolutely for $\abs{x-c}<R$ and diverges for $\abs{x-c}>R$. This defines the unique radius $R\in[0,\infty]$; no assertion about the boundary follows.
\end{theorem}
\begin{proof}
For $t=\abs{x-c}>0$, the limsup of the $n$th roots of the term magnitudes is $tL$. The root test \KRef{K03.11.root} gives the conclusion for $tL<1$ and $tL>1$, including $L=0$ or $\infty$. At $t=0$ only the constant term remains. The asserted interior and exterior determine $R$ uniquely: distinct radii would give both convergence and divergence at some intermediate distance. This proof uses only integer roots, not log, exp, or trigonometric functions.
\end{proof}

\KEntry{K03.13.ratio_radius}{The ratio formula when a ratio limit exists}
\begin{theorem}
If $a_n\ne0$ eventually and $\abs{a_{n+1}/a_n}\to\ell\in[0,\infty]$, then $R=1/\ell$. Equivalently a limit of $\abs{a_n/a_{n+1}}$ gives $R$ directly, with the same reciprocal conventions.
\end{theorem}
\begin{proof}
At positive distance $t$, consecutive term magnitudes have ratio tending to $t\ell$. The numerical ratio test gives convergence when $t\ell<1$ and divergence when $t\ell>1$. Comparing with the unique radius proves the formula, also when $\ell=0$ or $\infty$. If ratios do not exist, the root-limsup formula remains valid and no ratio formula is claimed.
\end{proof}

\KEntry{K03.13.compact_uniform}{Uniform convergence of all derivative series inside the radius}
\begin{theorem}
For $0\leq r<R$, the original series and every formal derivative series converge absolutely and uniformly on $\abs{x-c}\leq r$. In particular, for fixed integer $k\geq0$, the series with terms $n(n-1)\cdots(n-k+1)a_n(x-c)^{n-k}$, $n\geq k$, has this property.
\end{theorem}
\begin{proof}
Choose $r<\rho<R$ (if $r=0$ first replace it by any small positive radius). Absolute convergence at distance $\rho$ makes $\abs{a_n}\rho^n\leq M$. Put $q=r/\rho<1$. On the smaller disk the $k$th derivative term is bounded by $M\rho^{-k}n^kq^{n-k}$ for $n\geq k$. For fixed $k$, the ratio of consecutive bounds tends to $q$, since $(1+1/n)^k\to1$ by finite algebra. Thus these numerical bounds form a convergent series. The M-test proves both conclusions. No asymptotic estimate involving exponentials was needed.
\end{proof}

\KEntry{K03.13.endpoints}{Independent endpoint behavior}
\begin{example}
All three following real power series have radius $1$. The series $\sum_{n\geq0}x^n$ converges at neither endpoint, since its terms do not tend to zero at $\pm1$. The series $\sum_{n\geq1}x^n/n$ converges at $-1$ by the alternating test and diverges at $1$ by harmonic divergence. The series $\sum_{n\geq1}x^n/n^2$ converges absolutely at both endpoints by condensation. Their radius is $1$ by the root test and the elementary root estimates in \KRef{K03.11.root}. Replacing $x$ by $-x$ switches the one-endpoint example. Thus a radius never determines the endpoint verdicts.
\end{example}

\KEntry{K03.13.termwise}{Real-variable calculus of power series}
\begin{theorem}
Inside its radius a power series has derivatives of every order with respect to the real argument, given by formal termwise differentiation. Its primitive vanishing at $c$ is
\[\int_c^x\sum_{n\geq0}a_n(t-c)^n\dd t=\sum_{n\geq0}\frac{a_n}{n+1}(x-c)^{n+1}.\]
The derivative and primitive series have exactly the original radius (including radii $0$ and $\infty$).
\end{theorem}
\begin{proof}
On any compact interval strictly inside the radius, the derivative series is uniformly convergent by \KRef{K03.13.compact_uniform}; the original series converges at $c$. Apply termwise differentiation, \KRef{K03.12.termwise_differentiation}, enlarging the compact interval to include $c$. Repeat for higher derivatives. Termwise integration follows from \KRef{K03.12.termwise_integration} and the elementary integral of a monomial.

For exact equality of radii, the estimates just proved show that differentiation cannot reduce the radius. Conversely if the derivative series converges absolutely at distance $t>0$, then $\sum_{n\geq1}\abs{a_n}t^n\leq t\sum_{n\geq1}n\abs{a_n}t^{n-1}<\infty$; hence it cannot have a larger radius. The primitive series has at least the original radius by division by $n+1$, and its derivative is the original series, so the derivative result gives equality. These comparisons are statements about convergent sums even when the original radius is zero. Complex arguments were used for bounds only; complex differentiation belongs to D12.
\end{proof}

\KEntry{K03.13.uniqueness}{Coefficient uniqueness}
\begin{theorem}
If $f(x)=\sum a_n(x-c)^n$ in a real neighborhood of $c$, then $f^{(k)}(c)=k!a_k$ for every $k\geq0$. Consequently two such expansions of the same function at $c$ have identical coefficients.
\end{theorem}
\begin{proof}
Differentiate $k$ times using \KRef{K03.13.termwise} and set $x=c$. Terms with index below $k$ have vanished, the index $k$ term is $k!a_k$, and all terms with larger index contain a positive power of $x-c$ and vanish. Dividing by the nonzero integer $k!$ proves uniqueness.
\end{proof}

\KEntry{K03.13.products}{Products of power series}
\begin{theorem}
If $f=\sum a_n(x-c)^n$ and $g=\sum b_n(x-c)^n$ have radii $R_f,R_g$, their product for $\abs{x-c}<\min(R_f,R_g)$ is
\[fg=\sum_{n\geq0}\left(\sum_{k=0}^na_kb_{n-k}\right)(x-c)^n.\]
The resulting radius is at least the minimum and may be larger.
\end{theorem}
\begin{proof}
At every such argument both numerical series are absolutely convergent, so the absolute Cauchy product theorem \KRef{K03.11.cauchy_product} applies. Their Cauchy product terms are exactly the displayed monomials. For strict enlargement use $f(x)=\sum_{n\geq0}x^n$, $g(x)=1-x$: the coefficient products cancel to the constant $1$, whose radius is infinite although $R_f=1$.
\end{proof}

\KEntry{K03.13.composition}{Recentering, composition, and reciprocal series}
\begin{theorem}
A power-series sum with positive radius of convergence can be recentered at every interior point. Compositions are locally power-series sums whenever the inner value is inside the outer convergence disk. A power series of positive radius with $a_0\ne0$ has a reciprocal power series in a sufficiently small neighborhood of its center.
\end{theorem}
\begin{proof}
For $F(z)=\sum a_n(z-d)^n$ and $\abs{b-d}<R$, choose $r>0$ with $\abs{b-d}+r<R$. The finite binomial theorem and absolute double summation give
\[F(b+h)=\sum_{k\geq0}A_kh^k,\qquad A_k=\sum_{n\geq k}\binom nk a_n(b-d)^{n-k},\quad \abs h<r.\]
The decisive absolute bound is
\[\sum_n\sum_{k=0}^n\abs{a_n}\binom nk\abs{b-d}^{n-k}\abs h^k=\sum_n\abs{a_n}(\abs{b-d}+\abs h)^n<\infty.\]
Thus all coefficient sums exist and reordering is justified.

Let $g(c)=b$ lie inside that disk and write $g(c+h)=b+u(h)$ with $u(h)=\sum_{j\geq1}b_jh^j$. For sufficiently small $r>0$, the majorant $U(r)=\sum_{j\geq1}\abs{b_j}r^j$ is smaller than the convergence radius of the recentered outer series: indeed $U(r)\to0$, since for fixed $r_0$ with absolute convergence, $U(r)\leq(r/r_0)U(r_0)$. Expand each $u(h)^m$ by repeated absolutely convergent Cauchy products. Its degree-$k$ coefficient is
\[c_{m,k}=\sum_{j_1+\cdots+j_m=k,\ j_i\geq1}b_{j_1}\cdots b_{j_m},\]
with $c_{0,0}=1$, and it is zero for $k<m$. The total absolute sum is bounded by $\sum_m\abs{A_m}U(r)^m<\infty$. Hence
\[F(g(c+h))=\sum_{k\geq0}\left(\sum_{m=0}^kA_mc_{m,k}\right)h^k.\]
This proves the local claim from the value condition; a majorant bound on a whole preassigned disk is not inferred merely from a pointwise value bound.

For $f(c+h)=a_0(1-v(h))$, $v(0)=0$, choose $r$ with the majorant of $v$ less than $1$. The just proved composition with $\sum_{m\geq0}v^m=(1-v)^{-1}$ constructs $1/f$. Its coefficients satisfy $q_0=1/a_0$ and $q_k=-a_0^{-1}\sum_{j=1}^ka_jq_{k-j}$, by multiplying the convergent series and equating coefficients. This supplies both construction and convergence, not merely a formal recurrence.
\end{proof}

\KEntry{K03.13.analyticity}{Real analyticity}
\begin{definition}
A scalar function on an open real interval is real analytic if near each point it equals a power series in the real argument with positive radius. Such a function is $C^\infty$: apply \KRef{K03.13.termwise} in each defining neighborhood. Recentering in \KRef{K03.13.composition} shows that every power-series sum is analytic throughout its open convergence interval. Smoothness alone does not assert equality to the Taylor series; a counterexample is supplied in the Taylor chapter.
\end{definition}

\KEntry{K03.13.identity}{The identity principle on an interval}
\begin{theorem}
Two real-analytic scalar functions on a connected open interval which agree on a nonempty open subinterval, or on a set with an interior accumulation point, agree everywhere.
\end{theorem}
\begin{proof}
Subtract the functions. If an analytic expansion about $c$ has a first nonzero coefficient $a_m$, then near $c$ it factors as $(x-c)^m(a_m+\sum_{j>m}a_j(x-c)^{j-m})$. The parenthesis is continuous and nonzero near $c$, so $c$ is an isolated zero. Consequently an accumulation of zeros at $c$ forces every coefficient to vanish, and the function is zero in a neighborhood of $c$.

Let $J$ be the set of points with a neighborhood on which the function vanishes. It is open. If $c$ belongs to its closure inside the interval, then either $c\in J$ already or distinct zeros from $J$ accumulate at $c$; the previous argument gives $c\in J$. Thus $J$ is also relatively closed. It is nonempty under either stated hypothesis. Connectedness of the interval makes $J$ the whole interval (otherwise a boundary between a point in $J$ and a point outside it contradicts relative openness and closedness). This proves the result.
\end{proof}

\KEntry{K03.13.binomial}{The generalized binomial series}
\begin{theorem}
For real $\alpha$, put $c_0=1$ and $c_n=\alpha(\alpha-1)\cdots(\alpha-n+1)/n!$. For $\abs x<1$,
\[(1+x)^\alpha=\sum_{n\geq0}c_nx^n.\]
For $\alpha$ a nonnegative integer this is a terminating polynomial identity, valid for all real $x$ in its algebraic meaning. Otherwise the radius is $1$. At $x=1$ convergence holds exactly when $\alpha>-1$ and the sum is $2^\alpha$. At $x=-1$ convergence holds exactly when $\alpha>0$, with sum $0$. The terminating cases are assessed separately, including $\alpha=0$.
\end{theorem}
\begin{proof}
For a nonterminating series, $\abs{c_{n+1}/c_n}=\abs{\alpha-n}/(n+1)\to1$, so the radius is $1$. The coefficient relation $(n+1)c_{n+1}=(\alpha-n)c_n$ shows by termwise differentiation that its sum $y$ satisfies $(1+x)y'=\alpha y$, $y(0)=1$. Using the real powers already constructed in \KRef{K03.07.real_powers}, differentiate $(1+x)^{-\alpha}y(x)$: its derivative is zero on $(-1,1)$ and its value at zero is $1$. The zero-derivative theorem therefore identifies the sum with $(1+x)^\alpha$. This identification is a downstream application and is not used to construct real powers.

For the endpoint tests choose $N>\max(\alpha,0)$ and put $t_n=\abs{c_n}$ for $n\geq N$. Then $t_{n+1}=\frac{n-\alpha}{n+1}t_n$, and $c_n$ alternates in sign eventually. If $\alpha>0$, the identity $nt_n-(n+1)t_{n+1}=\alpha t_n$ gives $\sum_{n=N}^Mt_n\leq Nt_N/\alpha$, so both endpoints converge absolutely. If $-1<\alpha<0$, $t_n$ decreases to zero: a positive limit $l$ would give $t_n-t_{n+1}=(\alpha+1)t_n/(n+1)\geq(\alpha+1)l/(n+1)$, contradicting harmonic divergence. Thus $x=1$ converges by alternation. If $\alpha\leq-1$, the $t_n$ do not tend to zero. Endpoint sums at $1$, and at $-1$ when $\alpha>0$, follow from Abel's theorem below and the respective limits of $(1+x)^\alpha$. If $\alpha<0$ and the series at $-1$ converged, the same Abel theorem would force a finite limit as $x\to-1+$, whereas $(1+x)^\alpha\to+\infty$. This proves all nonterminating endpoint assertions. The finite binomial theorem proves the terminating cases directly.
\end{proof}

\KEntry{K03.13.abel}{Abel's endpoint theorem}
\begin{theorem}
If $\sum_{n\geq0}a_n=S$ ordinarily, then $\sum a_nx^n$ converges for $0\leq x<1$ and tends to $S$ as $x\to1-$. By replacing $a_n$ by $(-1)^na_n$, the corresponding assertion holds at the endpoint $-1$.
\end{theorem}
\begin{proof}
The partial sums $s_n$ are bounded, and hence $a_n$ are bounded; absolute convergence for $x<1$ follows by a geometric bound. Finite summation by parts, followed by the vanishing boundary $s_Nx^{N+1}\to0$, gives
\[\sum_{n\geq0}a_nx^n=(1-x)\sum_{n\geq0}s_nx^n.\]
The weights $(1-x)x^n$ are nonnegative and sum to one. Given $\eps>0$, choose $m$ with $\abs{s_n-S}<\eps$ for $n>m$. The weighted error from those indices is at most $\eps$; the finite prefix $(1-x)\sum_{n=0}^m(s_n-S)x^n$ tends to zero. This proves the limit, with no interchange of uncontrolled endpoint sums.
\end{proof}

\KEntry{K03.13.tauber}{Tauber's additional coefficient condition}
\begin{theorem}
Let $a_n\in\R$ or $\C$, suppose $na_n\to0$, and suppose $f(x)=\sum_{n\geq0}a_nx^n\to S$ as $x\to1-$, $0<x<1$. Then the ordinary series $\sum a_n$ converges to $S$.
\end{theorem}
\begin{proof}
The coefficient condition implies bounded coefficients, so the series is absolutely convergent at each such $x$. Put $x_N=1-1/N$, $N\geq2$, and $s_N=\sum_{n=0}^Na_n$. For given $\eps>0$ choose $m$ with $\abs{a_n}\leq\eps/n$ for $n>m$. For $N>m$, use $1-x_N^n\leq n(1-x_N)=n/N$ to estimate
\[\begin{split}
\abs{s_N-f(x_N)}&\leq\sum_{n=0}^m\abs{a_n}(1-x_N^n)
+\eps\sum_{n=m+1}^N\frac{1-x_N^n}{n}
+\eps\sum_{n>N}\frac{x_N^n}{n}\\
&\leq\sum_{n=0}^m\abs{a_n}(1-x_N^n)+\eps+\frac{\eps x_N^{N+1}}{N(1-x_N)}\leq o(1)+2\eps.
\end{split}\]
The finite prefix tends to zero. Letting $N\to\infty$ and then $\eps\downarrow0$ gives $s_N-f(x_N)\to0$. Since $x_N\to1-$, $f(x_N)\to S$; hence $s_N\to S$. This explicit coefficient estimate is the additional argument that Abel summability alone lacks.
\end{proof}

\begin{corollary}[The weaker averaged coefficient condition]
The conclusion of the preceding theorem also holds if its hypothesis $na_n\to0$ is replaced by
\[
 \frac1n\sum_{k=1}^n k a_k\longrightarrow0.
\]
\end{corollary}
\begin{proof}
Put $T_n=\sum_{k=1}^nka_k$, $b_0=a_0$ and $b_n=T_n/[n(n+1)]$ for $n\geq1$. Finite summation by parts gives, for $s_n=\sum_{k=0}^na_k$ and $B_n=\sum_{k=0}^nb_k$,
\[
 s_n-B_n=\frac{T_n}{n+1}\longrightarrow0.
\]
Indeed $a_k=(T_k-T_{k-1})/k$, and telescoping its finite sum yields this identity. Also $nb_n=T_n/(n+1)\to0$, so $b_n=O(1/n)$, hence in particular $B_n=O(n)$ and $s_n=O(n)$. The identity for $a_k$ and $T_n=O(n)$ also gives $a_n=O(1)$. These elementary bounds avoid any logarithmic construction. For $0<r<1$ all ensuing series are absolutely convergent, since $\sum nr^n$ converges by the ratio test, and summing the partial-sum identities gives
\[
 \sum_{n\geq0}a_nr^n=(1-r)\sum_{n\geq0}s_nr^n,
 \qquad \sum_{n\geq0}b_nr^n=(1-r)\sum_{n\geq0}B_nr^n.
\]
The difference tends to zero as $r\uparrow1$: for a given $\eps>0$, bound the finite prefix of $s_n-B_n$ after multiplication by $1-r$, and bound the remaining null-sequence tail by $\eps(1-r)\sum r^n\leq\eps$. Thus the $b$ series has Abel limit $S$. The already proved theorem applies to $b$, since $nb_n\to0$, and yields $B_n\to S$. The exact finite identity now gives $s_n\to S$. This proof does not impose an unstated bound on $na_n$.
\end{proof}
\SourceRef{Zorich II, section 16.3 exercises, exercise 8(b), PDF p.410; stronger source variant retained in addition to the selected catalog theorem}

\KEntry{K03.13.abel_not_converse}{Abel summability alone has no converse}
\begin{example}
For $a_n=(-1)^n$, the geometric identity gives $\sum_{n\geq0}a_nx^n=1/(1+x)$ for $0\leq x<1$, so the Abel limit is $1/2$. The ordinary partial sums alternate between $1$ and $0$ and do not converge. Here $na_n$ does not tend to zero, so Tauber's hypothesis fails exactly where it is needed.
\end{example}

\KEntry{K03.13.standard_series}{Elementary series and their domains}
\begin{theorem}
With the elementary functions defined and identified in \KRef{K03.07.exp_construction}, \KRef{K03.07.sin_cos}, \KRef{K03.07.log_construction}, and \KRef{K03.07.trig_inverses},
\[\begin{gathered}
e^x=\sum_{n=0}^\infty x^n/n!,\quad
\sin x=\sum_{n=0}^\infty(-1)^n x^{2n+1}/(2n+1)!,\quad
\cos x=\sum_{n=0}^\infty(-1)^n x^{2n}/(2n)!,\qquad x\in\R,\\
\log(1+x)=\sum_{n=1}^\infty(-1)^{n-1}x^n/n\quad(-1<x\leq1),\\
\arctan x=\sum_{n=0}^\infty(-1)^nx^{2n+1}/(2n+1)\quad(-1\leq x\leq1).
\end{gathered}\]
The binomial expansion, including its endpoint cases, is \KRef{K03.13.binomial}. These are identifications with previously constructed functions, not independent conflicting definitions.
\end{theorem}
\begin{proof}
The factorial-series convergence follows from the ratio bound on each bounded interval. The first three equalities are the series definitions in the cited entries; their derivative identities follow from the general termwise theorem. For the other two, integrate respectively $1/(1+t)=\sum_{n\geq0}(-t)^n$ and $1/(1+t^2)=\sum_{n\geq0}(-1)^nt^{2n}$ from $0$ to $x$ inside $(-1,1)$, where convergence is uniform on the integration segment. The derivatives and values at zero of log and arctan identify the primitives by \KRef{K03.05.zero_derivative}. At $x=1$ the log series converges by alternation and Abel identifies its sum with $\log2$; at $x=-1$ it is the negative harmonic series and diverges. At $x=\pm1$ the arctan series converges by alternation; Abel and continuity identify its sums with $\pm\pi/4$. For $\abs x>1$ the terms of the log and arctan series fail to tend to zero: the ratio of their magnitudes tends to $\abs x$ and $x^2$, respectively. Thus their indicated domains are exact real convergence domains.
\end{proof}

% BEGIN REVISED SOURCE chapters/M03.07.tex SHA256=e7174e3edc397d24284353aba45fb600fd28102366c7cd7c8b5d512c399acd30
\chapter{Elementary functions from power series}
This chapter uses the general series results of modules 11--13 as prior logical
ingredients, irrespective of the printed chapter order. Those general results
use factorials, finite algebra and completeness, and do not use the exponential
or logarithm constructed here. Throughout, derivatives of complex-valued
functions of a real parameter use real increments.
\SourceRef{Rudin, Chapter 8, PDF pp. 183--191; Zorich I, Sections 5.5 and 5.6.5. The series-first route is fixed here; integral and differential-equation descriptions are comparisons}

\KEntry{K03.07.exp_construction}{The exponential series}
\begin{construction}
For $z\in\C$ define
\[E(z)=\sum_{n=0}^{\infty}\frac{z^n}{n!},\qquad 0!=1.
\]
Its restriction to $\R$ will also be written $\exp$.
\end{construction}
\begin{proof}
For $|z|\le R$, the absolute values are bounded by $R^n/n!$.
When $n+1\ge2R$, successive terms of this majorant have ratio at most
$1/2$, so its tail is bounded by a geometric tail. Thus the series converges
absolutely everywhere and uniformly on every disk $|z|\le R$ by the
uniform majorant argument of \KRef{K03.12.m_test}. The differentiated
series on a real interval is $\sum_{n\ge1}x^{n-1}/(n-1)!$ and obeys the
same bound. The general termwise theorem \KRef{K03.13.termwise} therefore
gives $E'(x)=E(x)$. This reasoning is about a previously established general
power-series theorem, not a definition of differentiation by formal symbols.
\end{proof}

\KEntry{K03.07.exp_addition}{Addition law}
\begin{theorem}
For $z,w\in\C$, $E(z+w)=E(z)E(w)$, $E(0)=1$, and
$E(-z)=1/E(z)$.
\end{theorem}
\begin{proof}
The absolute Cauchy-product theorem \KRef{K03.11.cauchy_product} applies
because both factorial series converge absolutely. Its coefficient of total
degree $n$ is
\[\sum_{k=0}^n\frac{z^kw^{n-k}}{k!(n-k)!}
=\frac1{n!}\sum_{k=0}^n\binom nkz^kw^{n-k}
=\frac{(z+w)^n}{n!}.
\]
The constant value is read from the series. Setting $w=-z$ gives
$E(z)E(-z)=1$, proving both nonvanishing and the reciprocal formula.
\end{proof}

\KEntry{K03.07.exp_properties}{The real exponential}
\begin{theorem}
The real function $\exp$ is positive, differentiable with derivative itself,
strictly increasing, and maps $\R$ bijectively onto $(0,\infty)$. Its limits
at $-\infty$ and $+\infty$ are $0$ and $+\infty$, respectively.
\end{theorem}
\begin{proof}
For $x\ge0$ the series has nonnegative terms and $E(x)\ge1+x>0$.
For $x<0$, $E(x)=1/E(-x)>0$. The derivative proved above is consequently
positive, so the mean value theorem makes $E$ strictly increasing.
The bound $E(x)\ge1+x$ proves divergence to $+\infty$; the reciprocal law
proves the limit at $-\infty$. Continuity and the intermediate value theorem
then show that each positive number is assumed exactly once.
\end{proof}

\KEntry{K03.07.log_construction}{The inverse logarithm}
\begin{definition}
For $y>0$, $\log y$ is the unique real $x$ such that $E(x)=y$.
\end{definition}
\begin{proposition}
The logarithm is continuous and strictly increasing on $(0,\infty)$,
$\log1=0$, and $\log(uv)=\log u+\log v$ for $u,v>0$.
\end{proposition}
\begin{proof}
The inverse-continuity theorem \KRef{K03.03.inverse_continuity} applies to
the continuous strictly increasing surjection just proved. Alternatively,
if $x=\log y$, the images of $x-\eps$ and $x+\eps$ enclose $y$ and trap
the inverse of every sufficiently close positive number between these two
points. Finally $E(\log u+\log v)=uv$ by addition, and inverse uniqueness
proves the product formula. In particular $\log(1/u)=-\log u$.
\end{proof}

\KEntry{K03.07.log_derivative}{Derivative and integral comparison for logarithm}
\begin{theorem}
For $x>0$, $\log'(x)=1/x$. After the fundamental theorem in module 09,
\[\log x=\int_1^x\frac{\dd t}{t}\quad(x>0)
\]
with oriented integration.
\end{theorem}
\begin{proof}
The inverse derivative rule applies because $E'(\log x)=E(\log x)=x\ne0$;
it gives the asserted reciprocal derivative. On every compact interval in
$(0,\infty)$, $1/t$ is continuous, so \KRef{K03.09.ftc_first} makes the
accumulation integral a primitive. Its difference from $\log$ has derivative
zero, hence is constant on $(0,\infty)$. Both functions vanish at $1$.
Only this comparison uses integration; the construction of $\log$ does not.
\end{proof}

\KEntry{K03.07.real_powers}{Real powers and roots}
\begin{definition}
For $x>0$ and $\alpha\in\R$, put $x^\alpha=E(\alpha\log x)$.
\end{definition}
\begin{proposition}
For positive bases,
\[x^{\alpha+\beta}=x^\alpha x^\beta,\quad
(xy)^\alpha=x^\alpha y^\alpha,\quad (x^\alpha)^\beta=x^{\alpha\beta}.
\]
This agrees with integer powers and with the positive algebraic root
interpretation of rational powers.
\end{proposition}
\begin{proof}
The first two laws follow from exponential addition and logarithm products;
the third follows from $\log(E(\alpha\log x))=\alpha\log x$.
For integers the first law, $x^1=x$, and $x^0=1$ give the familiar positive
and negative powers. If $q\ge1$ and $p\in\Z$, then
$(x^{p/q})^q=x^p$ and $x^{p/q}>0$, so uniqueness of the positive $q$th
root from \KRef{K03.01.positive_roots} proves agreement.
For $\alpha>0$ the limit at $x\downarrow0$ is $0$, since $\log x\to-\infty$;
thus $0^\alpha=0$ is a continuous extension. For $\alpha<0$ the limit is
$+\infty$. No definition of $0^0$ follows. Negative bases admit integer
powers, and rational powers with odd reduced denominator by real roots;
there is no real-valued extension for all real exponents at a negative base.
\end{proof}

\KEntry{K03.07.sin_cos}{Sine and cosine series}
\begin{construction}
For real or complex $z$ define
\[S(z)=\sum_{n=0}^{\infty}\frac{(-1)^nz^{2n+1}}{(2n+1)!},\qquad
C(z)=\sum_{n=0}^{\infty}\frac{(-1)^nz^{2n}}{(2n)!}.
\]
Write $\sin=S$ and $\cos=C$. On real inputs $S'=C$ and $C'=-S$.
\end{construction}
\begin{proof}
Absolute values of either series are subseries of $\sum |z|^n/n!$;
on $|z|\le R$ the same majorant proves uniform convergence. Their
differentiated real series are the displayed cosine and negative sine
series, with the same compact majorants. Termwise differentiation yields
the formulas, with $S(0)=0$ and $C(0)=1$.
\end{proof}

\KEntry{K03.07.trig_properties}{Trigonometric identities and bounds}
\begin{theorem}
Sine is odd, cosine is even, and for real $x,y$,
\begin{align*}
S(x+y)&=S(x)C(y)+C(x)S(y),\\
C(x+y)&=C(x)C(y)-S(x)S(y),\\
S(x)^2+C(x)^2&=1.
\end{align*}
Consequently $|S(x)|,|C(x)|\le1$, $|S(x)|\le|x|$, and
$|C(x)-1|\le x^2/2$.
\end{theorem}
\begin{proof}
Parity follows term by term. To prove the addition formulas without assuming
geometric trigonometry, fix $y$ and subtract the proposed expressions from
$S(x+y)$ and $C(x+y)$, calling the differences $u(x),v(x)$.
The derivative formulas give $u'=v$, $v'=-u$, and $u(0)=v(0)=0$.
Hence $(u^2+v^2)'=0$, and its initial value zero implies $u=v=0$.
Likewise $(S^2+C^2)'=2SC-2CS=0$, with value one at zero.
The norm bounds follow, and the mean value inequality gives $|S(x)|\le|x|$.
For the sharper cosine estimate, put $h(x)=x^2/2-1+C(x)$ for $x\ge0$.
Then $h'=x-S(x)\ge0$, so $1-C(x)\le x^2/2$; evenness handles $x<0$.
\end{proof}

\KEntry{K03.07.pi_period}{Pi, periods and quadrants}
\begin{construction}
Let $c$ be the first positive zero of $C$ and define $\pi=2c$.
Then $S(c)=1$, $C(c)=0$, sine and cosine have period $2\pi$,
and their first-quadrant ranges are $S([0,c])=[0,1]$ and
$C([0,c])=[0,1]$ with strict increase and decrease, respectively.
\end{construction}
\begin{proof}
Continuity makes $C>0$ on some $[0,\delta]$, and then $S$ increases there,
with $S(\delta)>0$. If $C$ had no positive zero, continuity would force
$C>0$ everywhere on $(0,\infty)$. Thus $S(x)\ge S(\delta)>0$ for
$x\ge\delta$ and $C'(x)=-S(x)\le-S(\delta)$; the mean value theorem
would make $C$ negative for large $x$, a contradiction. A zero exists.
The zero set in a compact interval from $\delta$ to any such zero is
nonempty and closed, hence has a minimum $c>0$. On $(0,c)$, $C>0$,
so $S$ is positive and strictly increasing. The identity gives $S(c)=1$.
Since $C'=-S<0$ on $(0,c]$, the range statements follow.
Addition gives $S(x+c)=C(x)$, $C(x+c)=-S(x)$; applying these twice gives
$S(x+2c)=-S(x)$, $C(x+2c)=-C(x)$ and four times gives period $4c=2\pi$.
These relations and parity describe all quadrants: sine increases on
$[-c,c]$, decreases on $[c,3c]$, and cosine decreases on $[0,2c]$,
with translated intervals giving all other monotonicity intervals.
They also show that the zeros of sine are $k\pi$ and those of cosine
$c+k\pi$ ($k\in\Z$). There is no common positive period below $2\pi$:
a common period must have sine zero and cosine one, which these quadrant
signs allow only at even multiples of $\pi$.
This analytic definition does not assume a circle-length definition of $\pi$.
\end{proof}

\KEntry{K03.07.trig_inverses}{Principal inverse trigonometric functions}
\begin{definition}
$\arcsin:[-1,1]\to[-\pi/2,\pi/2]$ is the inverse of sine on that range;
$\arccos:[-1,1]\to[0,\pi]$ is the inverse of cosine there.
$\arctan:\R\to(-\pi/2,\pi/2)$ is the inverse of tangent on that interval.
\end{definition}
\begin{proof}[Existence and derivative formulas]
The preceding strict monotonicity and endpoint values prove the first two
bijections and inverse continuity. On $(-\pi/2,\pi/2)$ cosine is positive;
the quotient rule gives $(S/C)'=1/C^2>0$. At the endpoints its limits are
$-\infty,+\infty$ because the numerator tends to $-1,1$ and the
denominator to $0$ through positive values. Thus tangent is a continuous
strictly increasing surjection, establishing the third inverse.
For $|x|<1$, inverse differentiation and the sign of the appropriate sine
or cosine give
\[(\arcsin x)'=\frac1{\sqrt{1-x^2}},\qquad
(\arccos x)'=-\frac1{\sqrt{1-x^2}}.
\]
For all $x\in\R$, $(\arctan x)'=1/(1+x^2)$ because
$C(\arctan x)^{-2}=1+x^2$. At $x=\pm1$ the first two inverse functions
remain continuous, but their finite one-sided derivatives do not exist:
the interior derivative magnitudes tend to infinity and the mean value
theorem on a one-sided secant gives the same divergence of slopes.
\end{proof}

\KEntry{K03.07.tan}{Tangent and poles}
\begin{proposition}
$\tan x=S(x)/C(x)$ is defined exactly for
$x\notin\{\pi/2+k\pi:k\in\Z\}$, has period $\pi$, and derivative
$1/C(x)^2$. At each pole $p=\pi/2+k\pi$, the left and right limits
are $+\infty$ and $-\infty$.
\end{proposition}
\begin{proof}
The zero description for $C$ proves the domain. Both numerator and
denominator change sign after $\pi$, proving periodicity.
The quotient rule gives $(C^2+S^2)/C^2=1/C^2$.
At $\pi/2$, sine tends to $1$ while cosine is positive on the left and
negative on the right. Periodicity translates these one-sided limits to
every pole. A pole is excluded from the domain; the limits are not values.
\end{proof}

\KEntry{K03.07.hyperbolic}{Hyperbolic functions and their real inverses}
\begin{definition}
Put $\sinh x=(E(x)-E(-x))/2$, $\cosh x=(E(x)+E(-x))/2$ and
$\tanh x=\sinh x/\cosh x$ for $x\in\R$.
\end{definition}
\begin{proposition}
$\cosh^2x-\sinh^2x=1$, $(\sinh)'=\cosh$, $(\cosh)'=\sinh$,
and $(\tanh)'=1/\cosh^2$. The respective principal inverses are
\begin{align*}
\operatorname{arsinh}y&=\log(y+\sqrt{y^2+1})&& (y\in\R),\\
\operatorname{arcosh}y&=\log(y+\sqrt{y^2-1})&& (y\ge1),\\
\operatorname{artanh}y&=\tfrac12\log\frac{1+y}{1-y}&& (|y|<1),
\end{align*}
with ranges $\R,[0,\infty),\R$. Their derivatives on the interiors
of these domains are $(1+y^2)^{-1/2}$, $(y^2-1)^{-1/2}$,
and $(1-y^2)^{-1}$.
\end{proposition}
\begin{proof}
Expanding the squares and using $E(x)E(-x)=1$ proves the identity;
differentiation proves the three derivative formulas. Since $\cosh x>0$,
$\sinh$ increases strictly and tends to both real infinities. Cosine
hyperbolic decreases to $1$ on $(-\infty,0]$ and increases from $1$ on
$[0,\infty)$. The quotient expression
$\tanh x=(E(2x)-1)/(E(2x)+1)$ gives strict increase and limits $\pm1$.
Solving each equation using $u=E(x)>0$ gives the displayed logarithms;
for example $2y=u-u^{-1}$ has the unique positive root
$u=y+\sqrt{y^2+1}$. For cosh select $u\ge1$ to fix the nonnegative branch.
Inverse differentiation gives the stated derivatives from the identity,
with the positive square root on the chosen branch. At $y=1$ arcosh
is continuous with infinite right derivative; no finite derivative is claimed.
\end{proof}

\KEntry{K03.07.euler}{Imaginary inputs and Euler's identity}
\begin{theorem}
For real $t$, $E(it)=C(t)+iS(t)$, $|E(it)|=1$, and
$\overline{E(z)}=E(\bar z)$ for every $z\in\C$. For fixed $z\in\C$,
the real-parameter derivative of $E(tz)$ is $zE(tz)$.
\end{theorem}
\begin{proof}
Absolute convergence permits separating the even and odd terms in $E(it)$,
which gives the identity. The squared modulus is $C^2+S^2=1$.
Conjugation commutes with finite partial sums and limits, giving the second
identity. Uniform convergence of the differentiated series on bounded real
$t$-intervals gives
$\sum_{n\ge1}nz^nt^{n-1}/n!=zE(tz)$ componentwise. This asserts a real
parameter derivative; it does not invoke complex differentiability theory.
\end{proof}

\KEntry{K03.07.elementary_derivatives}{A derivative table with domains}
\begin{proposition}
For $x>0$ and $\alpha\in\R$,
\[(x^\alpha)'=\alpha x^{\alpha-1},\qquad
(a^x)'=(\log a)a^x\quad(a>0).
\]
Together with the formulas proved above this gives the derivatives of the
elementary functions on their stated domains. In particular, $\sec=1/\cos$
has derivative $\sec\tan$ where $\cos\ne0$, and $\cot=\cos/\sin$ has
derivative $-1/\sin^2$ where $\sin\ne0$.
\end{proposition}
\begin{proof}
The chain rule gives $(E(\alpha\log x))'=E(\alpha\log x)\alpha/x$,
equal to $\alpha x^{\alpha-1}$ by the power laws. The second formula is
the chain rule for $E(x\log a)$. The reciprocal and quotient rules give
the last two formulas. Thus every entry follows from an actual derivative
calculation with its denominator condition, rather than a formal table
applied across zeros or branch endpoints.
\end{proof}

\KEntry{K03.07.elementary_domains}{Domains and exceptional points}
\begin{definition}
In these notes $\log$ has real domain $(0,\infty)$, $x^\alpha$ initially
has $x>0$, $\sqrt{x}$ is the nonnegative root for $x\ge0$, and the three
inverse trigonometric ranges are exactly those fixed above. Reciprocal
trigonometric functions exclude their denominator zeros.
\end{definition}
\begin{example}
The equation $\sin y=x$ has a unique solution in $[-\pi/2,\pi/2]$, but
not a unique real solution without that restriction: for $|x|<1$ both
$\arcsin x$ and $\pi-\arcsin x$ solve it, along with their $2\pi$
translates. Also $(\sqrt{x})'=1/(2\sqrt{x})$ holds for $x>0$;
the right quotient at zero is $1/\sqrt{h}\to+\infty$.
These computations explain the exclusions. A complex logarithm and its
branches require a different domain theory and are outside this chapter.
\end{example}

\KEntry{K03.07.standard_limits}{Reference limits}
\begin{theorem}
As $h\to0$ through real nonzero values,
\[\frac{\sin h}{h}\to1,\qquad\frac{E(h)-1}{h}\to1,\qquad
\frac{\log(1+h)}h\to1\quad(1+h>0).
\]
More generally, if $u\to0$, $1+u>0$, and $v u\to L\in\R$, then
$(1+u)^v\to E(L)$, provided at zeros $u=0$ the expression is interpreted
as $1$ and the asserted product limit is consistent there.
\end{theorem}
\begin{proof}
The three quotients are derivative quotients at $0,0,1$, respectively,
whose derivatives have been proved to be one. For the last assertion,
at nonzero $u$,
$v\log(1+u)=(vu)\log(1+u)/u\to L$.
At zeros the exponent is zero; if zeros accumulate, existence of the
product limit forces $L=0$ along that subsequence. Continuity of $E$
completes the proof, including this case.
In particular $\lim_{h\to0}(1+h)^{1/h}=E(1)$ on the positive-base
domain, by taking $u=h$ and $v=1/h$ for nonzero $h$.
\end{proof}

\KEntry{K03.07.construction_bridges}{Sequence, differential equation and integral comparisons}
\begin{theorem}
For every fixed $x\in\R$,
\[\lim_{n\to\infty}(1+x/n)^n=E(x).
\]
On an interval containing zero, $E$ is the unique differentiable scalar
solution of $y'=y$, $y(0)=1$. The integral logarithm agrees with the
inverse logarithm.
\end{theorem}
\begin{proof}
For all sufficiently large $n$, $1+x/n>0$. If $x\ne0$, then
$n\log(1+x/n)=x\,\log(1+x/n)/(x/n)\to x$ by the preceding limit;
exponentiating proves the sequence assertion. For $x=0$ every term is one.
For the equation, existence is $E'=E$, $E(0)=1$. If $y$ is another
real or complex solution, the real derivative of $y(t)E(-t)$ is
$y'(t)E(-t)-y(t)E(-t)=0$. The zero-derivative theorem, componentwise
for complex values, gives $y(t)E(-t)=1$ on the interval.
Finally the integral comparison is precisely the primitive-uniqueness
argument of \KRef{K03.07.log_derivative}. These are proved identifications,
not simultaneous definitions that would create a circular construction.
\end{proof}

% BEGIN REVISED SOURCE chapters/M03.08.tex SHA256=15e2fccad382d1c4db02a0e0561f7afaad5f0b1f0f8a1c34095e89b0f97890e0
\chapter{The proper Riemann integral}
\SourceRef{Zorich I, Section 6.1, PDF pp. 351--363, for mesh-tagged sums; Rudin, Chapter 6, for Darboux sums with integrator $x$; Apostol, Section 7.26, PDF pp. 182--185, for countable covers and oscillation; Tao I, Chapter 11, for step-function bounds}
The basic real definition is bounded Darboux integration on a compact
nondegenerate interval. Integrability is a separate predicate, not a consequence
of writing an integral symbol or assigning a default value to nonintegrable data.

\KEntry{K03.08.partitions}{Partitions and tags}
\begin{definition}
For $a<b$, a partition is $P=(a=x_0<x_1<\cdots<x_n=b)$.
Its mesh is $|P|=\max_i\Delta x_i$, where $\Delta x_i=x_i-x_{i-1}$.
A tagging chooses $\xi_i\in[x_{i-1},x_i]$, and its sum is
$S(f;P,\xi)=\sum_i f(\xi_i)\Delta x_i$. A refinement inserts partition
points. Two partitions have the common refinement formed by their union.
For $[a,a]$ we use the empty sum, equal to zero, and no positive mesh.
\end{definition}

\KEntry{K03.08.darboux_sums}{Upper and lower sums}
\begin{definition}
For bounded $f:[a,b]\to\R$, set
\[M_i=\sup_{[x_{i-1},x_i]}f,\quad m_i=\inf_{[x_{i-1},x_i]}f,\quad
 U(f,P)=\sum_iM_i\Delta x_i,\quad L(f,P)=\sum_im_i\Delta x_i.
\]
\end{definition}
\begin{lemma}
Refinement increases $L$ and decreases $U$. For arbitrary $P,Q$,
$L(f,P)\le U(f,Q)$, and every tagged sum for $P$ lies between its
lower and upper sums.
\end{lemma}
\begin{proof}
On a subinterval inserted into an old cell, its infimum is at least the old
infimum and its supremum at most the old supremum. Multiplication by the
positive new lengths and their additivity gives the refinement inequalities.
For the common refinement $R$, $L(P)\le L(R)\le U(R)\le U(Q)$.
The last assertion is the termwise inequality $m_i\le f(\xi_i)\le M_i$.
\end{proof}

\KEntry{K03.08.integral}{The integrability predicate}
\begin{definition}
The lower and upper integrals are
\[\underline{I}(f)=\sup_P L(f,P),\qquad
\overline{I}(f)=\inf_P U(f,P).
\]
Boundedness ensures that these are finite and $\underline I\le\overline I$.
The function is Riemann integrable if they are equal; their common value
is $\int_a^bf$. A complex function is integrable if its real and imaginary
parts are, and its integral is their componentwise integral.
\end{definition}
For comparison with step-function definitions, a partition provides step
minorants and majorants with constants $m_i,M_i$ on the open cells,
assigning the finitely many endpoints their actual values. Their integrals
are $L(P),U(P)$. Conversely, for a finite step minorant $s\le f$, refine
at its finitely many breaks. Off those breaks its constants are bounded by
the infima of $f$ on the open cells. Shrinking the cells adjacent to breaks
changes a sum by at most a fixed bound times their total length. Consequently
$\int s\le\underline I(f)$ on letting that length tend to zero. The same
argument for majorants gives $\overline I(f)\le\int t$ for $t\ge f$.
Taking suprema and infima, together with the partition-provided bounds,
proves equality of the two definitions. Endpoint values have zero contribution
because finitely many neighboring lengths can be made arbitrarily small;
this is a proof of the comparison, not an identification of notation.

\KEntry{K03.08.criterion}{The Darboux criterion}
\begin{theorem}
A bounded real function is integrable if and only if for each $\eps>0$
some partition $P$ satisfies $U(f,P)-L(f,P)<\eps$.
\end{theorem}
\begin{proof}
If $\underline I=\overline I=I$, choose $P,Q$ with
$U(P)<I+\eps/2$ and $L(Q)>I-\eps/2$. Their common refinement
has gap less than $\eps$. Conversely every such partition gives
$0\le\overline I-\underline I\le U(P)-L(P)<\eps$, hence equality.
\end{proof}

\KEntry{K03.08.tagged_equivalence}{Equivalence with all fine tagged sums}
\begin{theorem}
For bounded real $f$, Darboux integrability with value $I$ is equivalent to:
for every $\eps>0$ there is $\delta>0$ such that
$|S(f;Q,\xi)-I|<\eps$ for every tagging of every $Q$ with $|Q|<\delta$.
The analogous statement for complex functions is equivalent componentwise.
\end{theorem}
\begin{proof}
Let $|f|\le M$ and fix $P$ with Darboux gap $<\eps/2$.
If $P$ has $m$ interior points, cells of $Q$ whose interiors contain
such a point have total length at most $m|Q|$; even counting cells
touching such points gives the harmless bound $2m|Q|$.
Every other $Q$ cell lies in a single $P$ cell. Insert the $P$ points
into $Q$ and compare sums: the good terms are unchanged, while the total
change on the exceptional cells has magnitude at most $4Mm|Q|$.
The refined tagged sum is between $L(P)$ and $U(P)$, as is $I$.
Thus $|S(Q,\xi)-I|\le\eps/2+4Mm|Q|$. Choose $\delta$ to make the
second term $<\eps/2$; if $M m=0$ it is already zero.

Conversely suppose all fine sums approach $I$. Choose a fine $Q$ such that
every sum differs from $I$ by less than $\eps/4$. In each cell choose a
tag with value within $\eps/(8(b-a))$ of its supremum and, independently,
a tag equally close to its infimum. Such tags exist by the definition of
supremum, without assuming extrema are attained. The two sums differ by
at most $\eps/2$, and together approximate $U(Q),L(Q)$ with total error
at most $\eps/4$. Hence $U(Q)-L(Q)<\eps$. The criterion proves
integrability, and the first direction plus uniqueness of a limit proves
the value is $I$. For complex values use $|\Re z|,|\Im z|\le|z|$ in one
direction and $|z|\le|\Re z|+|\Im z|$ in the other.
\end{proof}

\KEntry{K03.08.cauchy}{Cauchy criterion for sums}
\begin{theorem}
For bounded scalar $f$, a common fine-mesh limit exists if and only if
for every $\eps>0$ all two tagged sums of mesh below some $\delta$
differ in norm by less than $\eps$.
\end{theorem}
\begin{proof}
A common limit implies the assertion by the triangle inequality.
Conversely choose any sequence of tagged partitions with meshes tending
to zero (uniform partitions suffice). Their sums are a Cauchy sequence,
so scalar completeness gives a limit $I$. Given $\eps$, choose the
threshold for mutual distance $\eps/2$, then one of these sums below that
threshold and within $\eps/2$ of $I$. Every sum of mesh below the same
threshold is within $\eps$ of $I$. The preceding equivalence identifies
this limit with the Darboux integral.
\end{proof}

\KEntry{K03.08.continuous}{Continuous scalar functions}
\begin{theorem}
Every continuous real or complex function on $[a,b]$ is integrable.
\end{theorem}
\begin{proof}
Compactness gives boundedness and uniform continuity. For real $f$, choose
$\delta$ so that $|f(x)-f(y)|<\eps/(b-a)$ when $|x-y|<\delta$.
Any partition of mesh below $\delta$ then has each oscillation
$M_i-m_i\le\eps/(b-a)$, and hence gap at most $\eps$.
Apply the criterion with a smaller input tolerance for a strict bound.
For complex $f$, its two continuous real components are integrable.
\end{proof}

\KEntry{K03.08.monotone}{Monotone, step and piecewise continuous functions}
\begin{theorem}
Bounded monotone real functions, finite step functions, and functions
continuous on each of finitely many open cells with finite one-sided
boundary limits are integrable, provided all assigned boundary values are finite.
\end{theorem}
\begin{proof}
For increasing $f$ and a uniform partition of mesh $(b-a)/n$,
\[U-L=\frac{b-a}{n}\sum_{i=1}^n(f(x_i)-f(x_{i-1}))
=\frac{b-a}{n}(f(b)-f(a))\longrightarrow0.
\]
Decreasing functions follow on changing sign. A step function has zero
oscillation outside cells touching its finitely many breaks; their total
length is at most twice the number of breaks times the mesh, making the
gap tend to zero. A function in the last class extends continuously to
each closed cell using its one-sided limits. Those finitely many continuous
pieces are integrable by the previous theorem, and changing finitely many
endpoint values has zero effect by \KRef{K03.08.finite_changes}.
Combining their partitions proves integrability on the whole interval;
the same partition argument proves the additivity used here.
\end{proof}

\KEntry{K03.08.linearity}{Linearity and splitting intervals}
\begin{theorem}
Finite scalar linear combinations of integrable functions are integrable
and their integrals are the same combinations. For $c\in[a,b]$, a bounded
function is integrable on $[a,b]$ if and only if both restrictions to
$[a,c]$ and $[c,b]$ are, and then
$\int_a^bf=\int_a^cf+\int_c^bf$.
\end{theorem}
\begin{proof}
For real functions the oscillation inequality
$\osc(\alpha f+\beta g,J)\le|\alpha|\osc(f,J)+|\beta|\osc(g,J)$,
on common refinements of their small-gap partitions, proves integrability.
Every tagged sum is linear, so its limit is linear. For splitting, insert
$c$ into a partition. Its Darboux gap is the sum of the two nonnegative
restricted gaps. Small whole gaps therefore give small restricted gaps;
conversely combine two small-gap partitions. Tagged sums for partitions
containing $c$ split exactly, and their limits give the asserted equality.
Complex versions follow by components and complex scalar multiplication.
\end{proof}

\KEntry{K03.08.order}{Order, modulus and norm estimates}
\begin{theorem}
For integrable real $f\le g$, $\int f\le\int g$. For integrable real
or complex $f$, $|f|$ is integrable and
\[\left|\int_a^bf\right|\le\int_a^b|f|\le(b-a)\sup_{[a,b]}|f|.
\]
\end{theorem}
\begin{proof}
All paired tagged sums preserve a real pointwise inequality, so their
limits preserve it. For real values $||u|-|v||\le|u-v|$ bounds the
Darboux gap for $|f|$ by that for $f$. For complex values the oscillation
of $|f|$ is at most the sum of those of the two components, so common
partitions prove its integrability as well. Put $z=\int f$. If $z\ne0$,
choose $\lambda=\bar z/|z|$. Then
$|z|=\int\Re(\lambda f)\le\int|f|$; if $z=0$ the assertion follows
from positivity. Comparing $|f|$ with its supremum proves the last bound.
\end{proof}

\KEntry{K03.08.products}{Products and continuous composition}
\begin{theorem}
Products of bounded integrable scalar functions are integrable. If $f$ is
real integrable and $\Phi$ is continuous on a compact interval containing
its range, then $\Phi\circ f$ is integrable. In particular $1/f$ is
integrable when $|f|\ge d>0$.
\end{theorem}
\begin{proof}
If $|f|\le A$, $|g|\le B$, then
$\osc(fg,J)\le A\osc(g,J)+B\osc(f,J)$; take common small-gap
partitions. Complex products reduce to four real products.
For composition let $|\Phi|\le H$. Given $\eta>0$, uniform continuity
provides $d>0$ such that $|u-v|<d$ gives
$|\Phi(u)-\Phi(v)|<\eta$. On cells with $\osc(f,J)<d$ the new
oscillation is at most $\eta$. The total length of the remaining cells
is at most $(U(f,P)-L(f,P))/d$. Thus the composition gap is at most
$\eta(b-a)+2H(U-L)/d$, which can be made arbitrarily small in that order.
For reciprocals, even when the range crosses both signs, use
$|1/u-1/v|\le|u-v|/d^2$ directly on the actual range. The same estimate
works for a complex reciprocal with modulus at least $d$.
\end{proof}

\KEntry{K03.08.finite_changes}{Finite changes}
\begin{theorem}
If bounded $f,g$ differ at only finitely many points and $f$ is integrable,
then $g$ is integrable with the same integral.
\end{theorem}
\begin{proof}
It suffices to treat $h=g-f$, supported on $N$ points, with $|h|\le M$.
Only partition cells touching those points contribute nonzero oscillation;
their total length is at most $2N|P|$. Thus the Darboux gap is at most
$4MN|P|\to0$, and every tagged sum has magnitude at most
$2MN|P|\to0$. Hence $h$ is integrable with value zero, proving the
claim by linearity, componentwise if necessary. This argument also
justifies all finite endpoint assignments in step-function comparisons.
\end{proof}

\KEntry{K03.08.uniform_stability}{Uniform limits preserve integrability}
\begin{theorem}
If integrable $f_n:[a,b]\to\K$ converge uniformly to $f$, then $f$ is
bounded and integrable, and
$\int f_n\to\int f$ with error at most $(b-a)\|f_n-f\|_\infty$.
\end{theorem}
\begin{proof}
A uniform error at most one from any bounded $f_N$ bounds $f$.
In the real case,
$U(f,P)-L(f,P)\le U(f_N,P)-L(f_N,P)+2(b-a)\|f-f_N\|_\infty$.
First choose $N$ with small uniform error and then a small-gap partition
for $f_N$. This proves integrability; the norm estimate of the previous
entry gives convergence of integrals. Apply the argument to components
for complex values, and then use the complex integral norm estimate for
the same final error bound.
\end{proof}

\KEntry{K03.08.zero_covers}{Length-zero sets using countable covers}
\begin{definition}
$E\subset\R$ has length zero if, for each $\eps>0$, it is covered by
countably many open bounded intervals $I_j$ with $\sum_j|I_j|<\eps$.
Subsets of length-zero sets have length zero.
\end{definition}
\begin{theorem}
Countable unions of length-zero sets have length zero. No nondegenerate
interval has length zero.
\end{theorem}
\begin{proof}
For $E=\bigcup_{n\ge1}E_n$, cover $E_n$ with intervals of total length
less than $\eps2^{-n-1}$. Enumerate the pairs of indices in the resulting
countable family. Finite sums of lengths are bounded by the sum of the
row bounds, so the total is less than $\eps$.
If a countable open cover covers $[c,d]$, compactness supplies a finite
subcover. Its total length is at least $d-c$: sort all its endpoints
together with $c,d$; each resulting cell inside $[c,d]$ is covered by
at least one interval, and summing its length counts at most the total
length of the covering intervals. Thus $[c,d]$ cannot have length zero
when $c<d$. Every nondegenerate interval contains such a compact interval.
\end{proof}

\KEntry{K03.08.zero_vs_content}{Countable length zero versus finite content}
\begin{example}
Enumerate $\Q\cap[0,1]=\{r_n:n\ge1\}$. Intervals centered at $r_n$
of length $\eps2^{-n-1}$ cover this set with total length less than
$\eps$, so it has length zero. However, any finite interval cover of it
has total length at least one. Indeed the union of the closures of those
finitely many intervals is closed and contains a dense subset of $[0,1]$,
hence contains all of $[0,1]$. The elementary finite-length inequality
just proved gives the lower bound one. Thus countable-cover length zero
cannot be replaced by arbitrarily small finite-cover content.
\end{example}

\KEntry{K03.08.cantor_set}{The middle-thirds set}
\begin{construction}
Let $C_0=[0,1]$, and obtain $C_{n+1}$ by deleting the open middle third
of each component of $C_n$. Put $C=\bigcap_{n\ge0}C_n$.
\end{construction}
\begin{proof}[Compactness, perfection, empty interior and length zero]
$C_n$ consists of $2^n$ disjoint closed intervals of length $3^{-n}$.
They are nested, so $C$ is closed in $[0,1]$, compact and nonempty
(it contains $0,1$). Every component endpoint remains in all later sets.
For $x\in C$ and every $n$, the component of $C_n$ containing $x$ has
two endpoints in $C$, at least one different from $x$, at distance at most
$3^{-n}$. Hence $x$ is a limit point of $C$: the set is perfect.
If an open interval were contained in $C$, it would be contained in every
$C_n$, impossible once its length exceeds $3^{-n}$, since a connected
interval cannot cross the gaps between distinct components. Thus the
interior is empty. The total length of the $C_n$ components is $(2/3)^n$.
Enlarge their finitely many closed intervals slightly to open intervals,
adding less than any prescribed further length. Choosing $n$ large
produces arbitrarily small open covers, so $C$ has length zero.
Equivalently, $C$ consists of numbers admitting a ternary expansion using
only digits $0,2$: membership chooses a left or right component at each
stage, whose nested intervals shrink to that digit expansion, and conversely
such digits always select surviving components. At a deleted endpoint
use the expansion with only $0,2$, for example $1/3=0.0222\ldots_3$;
the alternative $0.1000\ldots_3$ is not the membership convention.
\end{proof}

\KEntry{K03.08.oscillation_criterion}{The discontinuity criterion}
\begin{theorem}[Lebesgue criterion, proved with interval covers]
A bounded real function on $[a,b]$ is Riemann integrable if and only if
its set of discontinuities has length zero.
\end{theorem}
\begin{proof}
For $x\in[a,b]$ define its relative oscillation by
\[\omega(x)=\inf_{r>0}\osc(f,(x-r,x+r)\cap[a,b]).\]
It is zero exactly when $f$ is continuous at $x$: continuity makes all
nearby values close, and conversely oscillation below $\eps$ bounds
$|f(y)-f(x)|$ for all sufficiently close $y$. The set
$E_\alpha=\{x:\omega(x)\ge\alpha\}$ is closed, since a neighborhood
with oscillation less than $\alpha$ also supplies such a neighborhood
at every sufficiently close point. It is consequently compact.
The discontinuity set is $D=\bigcup_{n\ge1}E_{1/n}$.

Suppose $f$ integrable. Choose a partition with gap below $\alpha\eta$.
Every $x\in E_\alpha$ in the interior of a cell forces that cell's
oscillation at least $\alpha$. The total length of these bad cells
is therefore below $\eta$. Their interiors cover $E_\alpha$ except
possibly for finitely many partition endpoints. Enlarge the bad cells
and cover those endpoints by tiny open intervals, adding less than $\eta$
to total length. Since $\eta$ is arbitrary, $E_\alpha$ has length zero,
and the countable-union theorem makes $D$ length zero.

Conversely suppose $D$ length zero and $|f|\le M$.
Choose $\alpha>0$ and cover the compact $E_\alpha\subset D$ by finitely
many open intervals of total length $<\eta$; let their union be $G$.
For every $x\notin E_\alpha$ select a relative open neighborhood $V_x$
on which oscillation is $<\alpha$. Together with $G$ these cover $[a,b]$.
There is a mesh $\delta>0$ such that every cell of mesh below $\delta$
is contained in one member of this cover. To see this elementary cover
fact, choose about each point a ball whose doubled radius lies in one
cover member, take finitely many of the smaller balls covering the compact
interval, and use the least of their radii. A cell of smaller length
meeting a selected smaller ball lies in its doubled ball.
On cells contained in a $V_x$ the oscillation is $<\alpha$.
All other cells lie in $G$, and their total length is at most $\eta$:
their disjoint interiors form a subset of the finite union of covering
intervals, so the finite endpoint length calculation applies.
Hence $U-L\le\alpha(b-a)+2M\eta$. Choose $\alpha$ and then $\eta$
arbitrarily small. The Darboux criterion proves integrability.
This argument constructed no Lebesgue integral and used no measure-theoretic
convergence theorem.
\end{proof}

\KEntry{K03.08.dirichlet}{The rational indicator}
\begin{example}
Let $D(x)=1$ on rational $x\in[0,1]$ and $D(x)=0$ otherwise.
Every nondegenerate interval contains both a rational and an irrational:
Archimedean rational density gives the first, and translating a sufficiently
small nonzero irrational multiple into a rational subinterval gives the
second. Thus every Darboux cell has supremum one and infimum zero.
For every partition $U(D,P)=1$, $L(D,P)=0$, so $D$ is not integrable.
It is discontinuous at every point for the same density reason.
\end{example}

\KEntry{K03.08.thomae}{Thomae's function}
\begin{example}
For a reduced fraction $p/q\in[0,1]$ with $q>0$ set $T(p/q)=1/q$,
and put $T(x)=0$ on irrationals. It is continuous at every irrational,
discontinuous at every rational, and integrable with integral zero.
\end{example}
\begin{proof}
For fixed $N$, only finitely many reduced fractions in $[0,1]$ have
$q\le N$, since $0\le p\le q$. Near an irrational $x$, choose a
neighborhood avoiding this finite set; then $0\le T<1/N$ there, proving
continuity. At a rational $p/q$, irrational points approaching it have
value zero, different from $1/q$, so continuity fails (one-sided at
endpoints). For an integral proof independent of the discontinuity criterion,
put tiny intervals around the finitely many rationals of denominator
$\le N$ and choose a partition for which all cells touching them have
total length less than $\eta$. On every other cell $T\le1/N$.
Thus $U(T,P)\le\eta+1/N$. Every cell contains an irrational, so
$L(T,P)=0$. Let $N\to\infty$ and $\eta\downarrow0$ to obtain the result.
\end{proof}

\KEntry{K03.08.countable_changes}{Countable modifications can fail}
\begin{example}
The zero function is integrable with value zero, while the rational
indicator is not integrable. They differ exactly on the countable set
$\Q\cap[0,1]$, which has length zero. Thus equality off a length-zero
set does not by itself preserve Riemann integrability. If both functions
are already bounded and Riemann integrable, such equality does imply
equal integrals: their difference $h$ is zero off a length-zero set;
at every continuity point of $h$, density of that complement forces
$h=0$. Alternatively, cover its level sets by the oscillation argument.
For a direct value proof, every cell has a point outside the zero set,
so one can choose tags there to make all sums for $h$ equal zero; the
all-tagged limit then equals zero. The distinction is the separate
existence hypothesis on both integrals.
\end{example}

\KEntry{K03.08.oriented}{Orientation and finite-dimensional values}
\begin{definition}
For $a>b$ define $\int_a^bf=-\int_b^af$, requiring integrability on
$[b,a]$; define $\int_a^af=0$. For a finite-vector function
$f=(f_1,\ldots,f_d)$, integrability means integrability of every component,
and $\int f=(\int f_1,\ldots,\int f_d)$.
\end{definition}
\begin{proposition}
Interval additivity holds for arbitrary ordered triples whenever $f$ is
integrable on the compact interval containing them. For Euclidean norm,
$\|\int_a^bf\|\le\int_{\min(a,b)}^{\max(a,b)}\|f\|$.
\end{proposition}
\begin{proof}
Sort the three points; ordinary additivity on the resulting adjacent
intervals and the sign convention prove every ordering. The norm of $f$
is integrable because its oscillation is bounded by the sum of component
oscillations. If $z=\int f\ne0$, the real linear functional
$v\mapsto\langle z/\|z\|,v\rangle$ has norm one and takes value
$\|z\|$ at $z$. Componentwise linearity and the pointwise bound by
$\|f\|$ prove the estimate; the case $z=0$ is immediate. Complex
vectors use the real part of the Hermitian pairing in the same argument.
\end{proof}

\KEntry{K03.08.bridge_lebesgue}{Outgoing interface to D11}
\begin{remark}
This is a routing-only record, outside the proved D03 scope. The future
D11 interface must specify its Lebesgue integral and prove that a bounded
Riemann-integrable function on a compact interval is Lebesgue integrable
with the same oriented value. No such assertion is assumed as a lemma
in these notes, and no matching formal-library declaration is claimed.
The countable-cover criterion proved above is an elementary D03 theorem;
it does not silently supply the full Lebesgue integral or this comparison.
\end{remark}

% BEGIN REVISED SOURCE chapters/M03.09.tex SHA256=56ac24dacf7f99e5c8c105d93b7ae80a406199295351710303fca6bb0593fa5b
\chapter{Primitives and finite integration formulas}
\SourceRef{Zorich I, Sections 6.2--6.4, especially PDF pp. 377--386; Apostol, Section 7.20; Tao I, Section 11.9. For rationalizing substitutions compare Zorich I, Sections 5.7.4--5.7.5}
Unless stated otherwise $a<b$ and integration is the proper Riemann
integration of module 08. A scalar may be real or complex; order and exact
mean-value conclusions require real values.

\KEntry{K03.09.primitive}{Primitives are functions}
\begin{definition}
On an open interval $I$, a primitive of $f:I\to\K$ is a differentiable
function $F$ with $F'=f$ at every point. On a general interval impose
this on its interior, and specify continuity or one-sided derivatives at
endpoints separately. The indefinite integral denotes the family of all
primitives, not a number; the definite integral denotes a value under an
integrability hypothesis.
\end{definition}

\KEntry{K03.09.uniqueness}{Uniqueness up to a constant}
\begin{theorem}
Two primitives of the same function on an interval differ by a constant.
If endpoints are included, assume their continuous extensions there.
\end{theorem}
\begin{proof}
Their difference has derivative zero. For any two interior points, the
real mean value theorem applied between them gives equal values.
For complex values apply this separately to real and imaginary parts.
Continuity passes the equality to included endpoints. On a disconnected
domain the argument applies to each interval component, with independent
constants: for example a function equal to $0$ on $(-\infty,0)$ and $1$
on $(0,\infty)$ has zero derivative on its domain but is not constant there.
\end{proof}

\KEntry{K03.09.accumulation}{Lipschitz continuity of accumulation}
\begin{theorem}
If $f$ is bounded and integrable on $[a,b]$, then
$F(x)=\int_a^xf$ satisfies
$|F(y)-F(x)|\le M|y-x|$ for $M=\sup|f|$.
\end{theorem}
\begin{proof}
Restrictions are integrable by interval additivity. For $x\le y$,
$F(y)-F(x)=\int_x^yf$, whose modulus is at most $M(y-x)$ by
\KRef{K03.08.order}. Reversing $x,y$ proves the other case. This
argument uses bounded integrability and does not assume continuity of $f$.
\end{proof}

\KEntry{K03.09.ftc_first}{The first fundamental theorem}
\begin{theorem}
For integrable $f$ on $[a,b]$, its accumulation $F(x)=\int_a^xf$
is differentiable at every interior continuity point $x$ of $f$, with
$F'(x)=f(x)$. At a relatively continuous endpoint the corresponding
one-sided derivative equals that endpoint value.
\end{theorem}
\begin{proof}
For $h\ne0$ with $x+h\in[a,b]$, additivity and orientation give
\[\frac{F(x+h)-F(x)}h-f(x)=\frac1h\int_x^{x+h}(f(t)-f(x))\dd t.
\]
Its modulus is bounded by
$\sup_{t\text{ between }x\text{ and }x+h}|f(t)-f(x)|$, which tends to
zero by continuity at $x$. The same estimate with one-sided $h$ proves
the endpoint statements, with no two-sided derivative claimed there.
\end{proof}

\KEntry{K03.09.ftc_second}{Newton--Leibniz with integrable derivative}
\begin{theorem}
Suppose $F:[a,b]\to\K$ is continuous, differentiable on $(a,b)$,
and its derivative, extended by any finite endpoint values, is bounded
and Riemann integrable. Then
\[\int_a^bF'(x)\dd x=F(b)-F(a).
\]
\end{theorem}
\begin{proof}
For real $F$ and any partition the mean value theorem supplies
$\xi_i\in(x_{i-1},x_i)$ with
$F(x_i)-F(x_{i-1})=F'(\xi_i)\Delta x_i$. Summing telescopes to
$F(b)-F(a)$. As the mesh tends to zero, these sums tend to $\int F'$
by the all-tagged equivalence. For complex $F$ apply the real argument
to its components; their mean-value tags may differ, which causes no
problem because each real integral has the all-tagged limit.
Finite endpoint assignments do not change the integral.
\end{proof}

\KEntry{K03.09.continuous_primitive}{Existence for continuous functions}
\begin{theorem}
Every continuous scalar function on an interval $I$ has a primitive on
its interior, continuously defined at every included finite endpoint.
\end{theorem}
\begin{proof}
Fix $x_0\in I$ and define $F(x)=\int_{x_0}^xf$ on each compact
interval between $x_0$ and $x$. Continuity supplies proper integrability
there. Additivity ensures the definitions agree as the chosen compact
interval grows. The first fundamental theorem gives $F'=f$ at every
interior point. Local boundedness gives continuity at each included
endpoint by the accumulation estimate. No assertion of a finite value
at an excluded infinite or singular endpoint is involved.
\end{proof}

\KEntry{K03.09.parts}{Integration by parts}
\begin{theorem}
For $u,v\in C^1([a,b],\K)$,
\[\int_a^buv'=u(b)v(b)-u(a)v(a)-\int_a^bu'v.
\]
More generally it suffices that $u,v$ are continuous on $[a,b]$,
differentiable inside, and their derivatives admit bounded integrable
endpoint extensions.
\end{theorem}
\begin{proof}
The continuous $u,v$ are bounded integrable, so products with their
integrable derivatives are integrable by \KRef{K03.08.products}.
The product rule gives $(uv)'=u'v+uv'$ in the interior. The second
fundamental theorem applied to $uv$ and linearity give the formula.
The $C^1$ case satisfies all these hypotheses directly.
\end{proof}

\KEntry{K03.09.substitution}{Substitution for continuous integrands}
\begin{theorem}
If $\phi\in C^1([a,b],\R)$ and $f$ is continuous on an interval
containing $\phi([a,b])$, then
\[\int_a^bf(\phi(t))\phi'(t)\dd t
=\int_{\phi(a)}^{\phi(b)}f(x)\dd x.
\]
No monotonicity of $\phi$ is required in this version.
\end{theorem}
\begin{proof}
Let $m=\min\phi([a,b])$ and $M=\max\phi([a,b])$. Extend the restriction of $f$ to $[m,M]$ constantly beyond its endpoints: use $f(m)$ for $x<m$ and $f(M)$ for $x>M$. This gives a continuous scalar function $\widetilde f$ on $\R$, also when $m=M$. Choose its accumulation primitive $H$ on $\R$. This open-domain extension is needed even for a turning point such as $\phi(t)=t^2$ when the original $f$ was defined only on a closed interval.
The chain rule gives $(H\circ\phi)'=(\widetilde f\circ\phi)\phi'=(f\circ\phi)\phi'$, continuous
on $[a,b]$ with the one-sided endpoint convention. Newton--Leibniz gives
the left side as $H(\phi(b))-H(\phi(a))$. The same theorem for $H$
identifies this difference with the oriented integral on the right, because $\widetilde f=f$ between the two image endpoints.
\end{proof}

\KEntry{K03.09.monotone_substitution}{Substitution for Riemann-integrable integrands}
\begin{theorem}
If $\phi\in C^1([a,b],\R)$ is monotone and $f$ is bounded and Riemann
integrable on the interval between $\phi(a),\phi(b)$, then
$(f\circ\phi)\phi'$ is integrable and the same substitution formula holds.
\end{theorem}
\begin{proof}
First let $\phi$ be nondecreasing and nonconstant, $|f|\le M$, and
$L=\sup|\phi'|$. For a tagged partition $P$ put
$d_i=(\phi(t_i)-\phi(t_{i-1}))/\Delta t_i$.
The mean value theorem and uniform continuity of $\phi'$ imply
$|\phi'(\xi_i)-d_i|\le\omega_{\phi'}(|P|)$, where this uniform
modulus tends to zero. Thus the tagged sum for $g=(f\circ\phi)\phi'$
differs by at most $M(b-a)\omega_{\phi'}(|P|)$ from
\[\sum_i f(\phi(\xi_i))(\phi(t_i)-\phi(t_{i-1})).\]
After deleting repeated image points, this is a tagged sum for $f$ on
$[\phi(a),\phi(b)]$: terms on constant image cells are zero, and each
other image tag lies in its image cell. Its mesh is at most $L|P|$.
All such sums consequently approach $\int_{\phi(a)}^{\phi(b)}f$.
The displayed error tends to zero uniformly in tags, so $g$, which is
bounded, is integrable by \KRef{K03.08.tagged_equivalence} with the
required value. For constant $\phi$, $\phi'=0$ and both sides are zero.
For nonincreasing $\phi$, apply the increasing result to $-\phi$ and
$x\mapsto f(-x)$, then reverse orientation. Complex $f$ follows by
components. This proves existence of the transformed integral rather
than adding it as an extra assumption.
\end{proof}

\KEntry{K03.09.mean_value}{The weighted integral mean value theorem}
\begin{theorem}
For continuous real $f$ and integrable real $g\ge0$ on $[a,b]$, some
$c\in[a,b]$ satisfies $\int fg=f(c)\int g$.
\end{theorem}
\begin{proof}
Let $m=\min f$, $M=\max f$ and $G=\int g\ge0$.
Order gives $mG\le\int fg\le MG$. If $G=0$ these bounds make
$\int fg=0$, and any $c$ works. If $G>0$, the ratio $(\int fg)/G$
lies in $[m,M]$, and the intermediate value theorem supplies $c$.
The exact point conclusion is real; for complex $f$ the analogous
norm bounds hold but a point of equality need not exist.
\end{proof}

\KEntry{K03.09.second_mean_value}{The normalized second mean value theorem}
\begin{theorem}
If $f$ is real integrable and $g$ is nonnegative and nonincreasing on
$[a,b]$, then some $\xi\in[a,b]$ satisfies
\[\int_a^bg(x)f(x)\dd x=g(a)\int_a^\xi f(x)\dd x.
\]
\end{theorem}
\begin{proof}
The bounded monotone $g$ and its product with $f$ are integrable.
Let $F(x)=\int_a^xf$, continuous by the accumulation estimate, with
minimum $m$ and maximum $M$; $m\le F(a)=0\le M$.
For a partition put $g_i=g(x_i)$. The difference between $\int gf$ and
$\sum_{i=1}^ng_{i-1}(F(x_i)-F(x_{i-1}))$ is bounded by
\[\|f\|_\infty\sum_i\osc(g,[x_{i-1},x_i])\Delta x_i
\le\|f\|_\infty |P|(g(a)-g(b))\longrightarrow0.
\]
Finite summation by parts writes that sum as
\[g_{n-1}F(b)+\sum_{i=1}^{n-1}(g_{i-1}-g_i)F(x_i).
\]
All weights are nonnegative and sum to $g(a)$. It is consequently
between $mg(a)$ and $Mg(a)$, and the same holds for its limit $\int gf$.
If $g(a)=0$, monotonicity and nonnegativity give $g=0$. Otherwise
$(\int gf)/g(a)\in[m,M]$ and continuity of $F$ gives the required
$\xi$. The proof allows jumps and retains the actual endpoint value
$g(a)$, as in Zorich I, formula (6.33); it does not silently assume $C^1$.
\end{proof}

\KEntry{K03.09.elementary_primitives}{Elementary primitives with domains}
\begin{proposition}
The following primitive formulas hold on any interval in the indicated domain:
\begin{align*}
\int x^\alpha\dd x&=x^{\alpha+1}/(\alpha+1)+C
 &&(x>0,\ \alpha\ne-1),\\
\int x^{-1}\dd x&=\log|x|+C&&(x\ne0),\\
\int E(ax)\dd x&=E(ax)/a+C&&(a\ne0),\\
\int\cos x\dd x&=\sin x+C,&
\int\sin x\dd x&=-\cos x+C,\\
\int\frac{\dd x}{1+x^2}&=\arctan x+C,&
\int\frac{\dd x}{\sqrt{1-x^2}}&=\arcsin x+C\quad(|x|<1),\\
\int\tan x\dd x&=-\log|\cos x|+C&&(\cos x\ne0).
\end{align*}
For integer $\alpha\ge0$ the power formula holds on all of $\R$.
\end{proposition}
\begin{proof}
Differentiate each right side using module 07. In the logarithmic formula,
on $x<0$ the chain rule gives $(\log(-x))'=(-1)/(-x)=1/x$;
on $x>0$ the usual logarithm derivative applies. Likewise
$(\log|\cos x|)'=-\sin x/\cos x$ on each pole-free interval.
The other displayed derivatives are respectively $x^\alpha$, $E(ax)$,
$\cos x$, $\sin x$, $(1+x^2)^{-1}$ and $(1-x^2)^{-1/2}$.
Uniqueness of primitives gives the additive constants on each interval.
Constants on different sides of a pole are independent. The excluded
case $a=0$ is the constant integrand one, with primitive $x+C$.
\end{proof}

\KEntry{K03.09.trig_substitution}{The half-angle rationalization}
\begin{proposition}
On $I_k=((2k-1)\pi,(2k+1)\pi)$, put $t=\tan(x/2)$.
This is a $C^1$ increasing bijection onto $\R$, with inverse
$x=2\arctan t+2k\pi$, and
\[\sin x=\frac{2t}{1+t^2},\qquad
\cos x=\frac{1-t^2}{1+t^2},\qquad
\dd x=\frac{2\dd t}{1+t^2}.
\]
\end{proposition}
\begin{proof}
The tangent branch results give the bijection and its inverse.
For $u=x/2$, divide $\sin2u=2\sin u\cos u$ and
$\cos2u=\cos^2u-\sin^2u$ by
$1=\cos^2u(1+\tan^2u)$ to obtain the identities.
Differentiating the inverse gives $\dd x/\dd t=2/(1+t^2)$.
The proper substitution theorem applies on compact subintervals avoiding
the endpoints and all zeros of the original rational denominator.
Thus a rational expression $R(\sin x,\cos x)$ becomes rational in $t$,
but a primitive must be continued separately across the excluded branch
endpoints or denominator poles if continuation is possible.
\end{proof}

\KEntry{K03.09.algebraic_substitution}{Selected real square-root substitutions}
\begin{example}
For $A>0$, on $-A<x<A$, $x=A\sin u$ with
$u\in(-\pi/2,\pi/2)$ gives
$\sqrt{A^2-x^2}=A\cos u>0$ and $\dd x=A\cos u\dd u$.
On $x>A$, $x=A\cosh u$, $u>0$, gives
$\sqrt{x^2-A^2}=A\sinh u>0$ and $\dd x=A\sinh u\dd u$.
On all real $x$, $x=A\sinh u$ gives
$\sqrt{x^2+A^2}=A\cosh u>0$ and $\dd x=A\cosh u\dd u$.
\end{example}
\begin{proof}
Each radicand identity follows from the trigonometric or hyperbolic
square identity. The branch restrictions fix the sign of the square root
and make the substitution a bijection. The differentials follow by
differentiation. For a linear root $y=\sqrt{\alpha x+\beta}$ with
$\alpha\ne0$ and $\alpha x+\beta>0$, use
$x=(y^2-\beta)/\alpha$, $y>0$, and $\dd x=2y\dd y/\alpha$.
Completing the square, when the residual constant is nonzero, reduces
a quadratic with nonzero leading coefficient to a positive scaling of
one of the three displayed radicands on intervals where it is positive,
or reveals that no positive real branch exists. For a zero residual
constant, $\sqrt{A(x-h)^2}=\sqrt A\,|x-h|$ if $A>0$; on each side
of $h$ this is a fixed linear branch. If $A<0$ there is no open interval
of positive radicand. For instance
$\int\dd x/\sqrt{x^2+A^2}=u+C=\operatorname{arsinh}(x/A)+C$.
All reductions use the substitution theorem on regular compact
subintervals. These selected reductions do not assert a universal
algorithm for elementary antiderivatives.
\end{proof}

\KEntry{K03.09.holder}{H\"older's finite integral inequality}
\begin{theorem}
Let $p,q>1$, $1/p+1/q=1$, and let $f,g$ be integrable real or complex
functions on $[a,b]$. Then
\[\int_a^b|fg|\le
\left(\int_a^b|f|^p\right)^{1/p}
\left(\int_a^b|g|^q\right)^{1/q}.
\]
\end{theorem}
\begin{proof}
All absolute powers are integrable by continuous composition on bounded
ranges. Young's pointwise inequality $uv\le u^p/p+v^q/q$ for $u,v\ge0$
follows as follows: for fixed $v>0$ the derivative of
$u^p/p-uv+v^q/q$ is $u^{p-1}-v$, so its minimum is at
$u=v^{1/(p-1)}$ and has value zero because $q=p/(p-1)$.
The case $v=0$ is immediate.
Put $A=(\int|f|^p)^{1/p}$, $B=(\int|g|^q)^{1/q}$.
If $A,B>0$, apply Young to $|f|/A,|g|/B$ and integrate, obtaining
$\int|fg|/(AB)\le1/p+1/q=1$.
If $A=0$, Young applied to $|f|/\delta$ and $|g|$ gives
$\int|fg|\le\delta\int|g|^q/q$ for every $\delta>0$, so the left
side is zero. If $B=0$, interchange the roles. This handles zero
seminorms without claiming the functions vanish at every point.
The case $p=q=2$ is Cauchy--Schwarz.
\end{proof}

\KEntry{K03.09.minkowski}{Minkowski's finite integral inequality}
\begin{theorem}
For $p\ge1$, write $\|f\|_p=(\int_a^b|f|^p)^{1/p}$.
For integrable scalar $f,g$, $\|f+g\|_p\le\|f\|_p+\|g\|_p$.
\end{theorem}
\begin{proof}
For $p=1$ integrate $|f+g|\le|f|+|g|$.
For $p>1$ put $H=\|f+g\|_p$. All functions below are integrable
by products and continuous composition. If $H=0$ there is nothing to
prove. Otherwise the triangle inequality followed by H\"older with
$q=p/(p-1)$ gives
\begin{align*}
H^p&\le\int(|f|+|g|)|f+g|^{p-1}\\
&\le(\|f\|_p+\|g\|_p)
\left(\int|f+g|^{(p-1)q}\right)^{1/q}
=(\|f\|_p+\|g\|_p)H^{p-1}.
\end{align*}
Divide by $H^{p-1}$. This is a seminorm on actual functions: a function
nonzero at just one point can have zero seminorm. No quotient space is
being constructed here.
\end{proof}

\KEntry{K03.09.partial_fractions}{Rational integration and the algebraic interface}
\begin{construction}
The shared algebra interface used here is the following precise result:
for real polynomials $P,Q$, $Q\ne0$, Euclidean division and factorization
of $Q$ into real linear factors and irreducible real quadratic factors
provide a finite identity
\[\frac P Q=A(x)+\sum_{r,k}\frac{a_{rk}}{(x-r)^k}
+\sum_{j,k}\frac{b_{jk}x+c_{jk}}{(x^2+u_jx+v_j)^k},
\qquad u_j^2-4v_j<0.
\]
This factorization and partial-fraction algebra are accepted from the
fixed shared algebra foundation (D07), not asserted as newly proved D03
theorems. The identity holds where $Q\ne0$.
\end{construction}
\begin{proof}[The analytic reduction]
Polynomial terms integrate by powers. Linear terms integrate to
$a\log|x-r|$ for $k=1$ and
$a(x-r)^{1-k}/(1-k)$ for $k>1$ on intervals excluding $r$.
For a quadratic write $y=x+u/2$, $d=v-u^2/4>0$ and split the numerator
as $By+C$. The $By$ term has primitive
$(B/2)\log(y^2+d)$ if $k=1$, and
$B(y^2+d)^{1-k}/(2(1-k))$ otherwise.
For the constant term define $J_k=\int\dd y/(y^2+d)^k$.
$J_1=d^{-1/2}\arctan(y/\sqrt d)$, and differentiating
$y/(y^2+d)^{k-1}$ gives
\[\left(\frac{y}{(y^2+d)^{k-1}}\right)'
=\frac{2(k-1)d}{(y^2+d)^k}-\frac{2k-3}{(y^2+d)^{k-1}}.
\]
Therefore
\[J_k=\frac{y}{2(k-1)d(y^2+d)^{k-1}}
+\frac{2k-3}{2(k-1)d}J_{k-1}\quad(k\ge2).
\]
This finite recurrence completes the reduction to rational functions,
logarithms and arctangents, and differentiating verifies every term.
Complex partial fractions can instead be reduced on real intervals
componentwise, or by this real decomposition of their real and imaginary
rational parts; a complex logarithm branch is not introduced implicitly.
\end{proof}

\KEntry{K03.09.area_length}{Area and length of a scalar graph}
\begin{theorem}
For continuous $f\ge0$ on $[a,b]$, the area of its subgraph is $\int_a^bf$
when area is defined by inner and outer finite axis-parallel rectangle
contents. For $f\in C^1([a,b])$, its graph has polygonal length
\[L=\int_a^b\sqrt{1+f'(x)^2}\dd x.
\]
\end{theorem}
\begin{proof}
For area, a partition gives interior columns of heights $\inf_J f$ and
covering columns of heights $\sup_J f$. Their total areas are $L(f,P)$
and $U(f,P)$. Boundaries can be covered by finitely many rectangles
of arbitrarily small total area; a common finite grid proves the elementary
inequality that every inner finite rectangle sum is at most every outer
one (decompose all overlaps by their finitely many horizontal and vertical
edges). The inner content is thus bounded above by every $U(f,P)$ and
the outer content below by every $L(f,P)$, while these column constructions
give the reverse bounds in the limit. The Darboux gap tends to zero,
so both contents equal $\int f$.

Define graph length as the supremum over partitions of the sum of chord
lengths for $\gamma(x)=(x,f(x))$. Since
$\gamma(y)-\gamma(x)=\int_x^y(1,f')$, the vector norm estimate gives
every chord length at most $\int_x^y\sqrt{1+f'^2}$. Hence all polygonal
lengths are bounded above by the displayed integral. Conversely, for a
tagged fine partition, uniform continuity of $f'$ gives
\[\left\|\gamma(x_i)-\gamma(x_{i-1})
-(1,f'(\xi_i))\Delta x_i\right\|
\le\omega_{f'}(|P|)\Delta x_i.
\]
The reverse triangle inequality bounds the difference of polygonal length
and the tagged sum for $\sqrt{1+f'^2}$ by
$(b-a)\omega_{f'}(|P|)\to0$. These polygonal lengths approach the
integral, proving the opposite inequality for their supremum.
\end{proof}

\KEntry{K03.09.derivative_not_integrable}{A primitive with nonintegrable derivative}
\begin{example}
Define $F(x)=x^2\sin(1/x^2)$ for $x\ne0$ and $F(0)=0$ on $[-1,1]$.
Then $F'(0)=\lim_{h\to0}h\sin(1/h^2)=0$, while for $x\ne0$,
\[F'(x)=2x\sin(1/x^2)-\frac2x\cos(1/x^2).
\]
At $x_n=(2\pi n)^{-1/2}$ this equals $-2/x_n\to-\infty$.
Thus $F$ is a primitive of its everywhere defined derivative, but the
derivative is unbounded on $[-1,1]$ and is not properly Riemann integrable.
This is the unbounded-derivative example of \KRef{K03.04.unbounded_derivative},
and shows why existence of a primitive does not remove the integrability
condition from Newton--Leibniz.
\end{example}

% BEGIN REVISED SOURCE chapters/M03.06.tex SHA256=23aa70b7abf4d108724f4fe22223811c8154a61becffd2e0d83d1336b3e9678d
\chapter{Taylor formulas and controlled remainders}
A Taylor polynomial is a finite object. Its approximation properties depend on local differentiability, and a conclusion about an infinite Taylor series needs additional information. We retain separately the Peano, real single-point, and integral versions of the remainder.
\SourceRef{Zorich I, 5.3.3 and 6.3.3; Apostol, 5.12, PDF 126--127 (printed 113--114)}

\KEntry{K03.06.taylor_polynomial}{The Taylor polynomial}
\begin{definition}
When the derivatives through order $n$ exist at $a$, define
\[ T_n(f,a)(x)=\sum_{k=0}^n\frac{f^{(k)}(a)}{k!}(x-a)^k,\qquad R_n(f,a)(x)=f(x)-T_n(f,a)(x). \]
Here $0!=1$ and $f^{(0)}=f$. This polynomial is the unique polynomial of degree at most $n$ whose derivatives of orders $0,\ldots,n$ agree with those of $f$ at $a$. A first nonzero coefficient describes a finite vanishing order; the definition makes no infinite-series assertion.
\end{definition}
\begin{proof}[Proof of uniqueness]
For a polynomial expressed in powers of $x-a$, differentiating $k$ times and evaluating at $a$ gives $k!$ times its $k$th coefficient. Thus all coefficients are forced to be the displayed ones, which also proves existence of this matching polynomial.
\end{proof}

\KEntry{K03.06.peano}{Taylor's theorem with Peano remainder}
\begin{theorem}
Let $n\ge1$. Suppose $f:I\to\K$ has derivatives through order $n-1$ in a neighborhood of interior $a$, and $f^{(n-1)}$ is differentiable at $a$. Then
\[ f(a+h)=\sum_{k=0}^n\frac{f^{(k)}(a)}{k!}h^k+o(h^n). \]
For $n=0$, the corresponding assertion $f(a+h)=f(a)+o(1)$ is exactly continuity at $a$.
\end{theorem}
\begin{proof}
Subtract $T_n(f,a)$ and write the result as $r$. Then $r^{(j)}(a)=0$ for $0\le j\le n$, and differentiability of $r^{(n-1)}$ at $a$ gives $r^{(n-1)}(a+t)=o(t)$. Given $\eps>0$, choose $\delta>0$ with $|r^{(n-1)}(a+t)|\le\eps|t|$ for $|t|<\delta$. If $n=1$ this is the conclusion. For each lower order, apply the norm mean inequality on the segment from $a$ to $a+t$. If already $|r^{(j+1)}(a+s)|\le\eps|s|^{n-j-1}$ for all sufficiently small $s$, then on that segment its derivative is bounded by $\eps|t|^{n-j-1}$; since $r^{(j)}(a)=0$,
\[ |r^{(j)}(a+t)|\le\eps|t|^{n-j}. \]
The required endpoint continuity follows from the locally existing lower derivatives, and at the highest step from differentiability of $r^{(n-1)}$ at $a$. Descending to $j=0$ gives $|r(a+h)|\le\eps|h|^n$. Arbitrarily small $\eps$ proves the result for both scalar targets without assuming the nth derivative exists away from $a$.
\end{proof}

\KEntry{K03.06.lagrange}{The real Lagrange remainder}
\begin{theorem}
Let $a\ne x$, $n\ge0$, and let real $f$ belong to $C^n$ on the closed segment joining $a$ and $x$, with $f^{(n+1)}$ existing at every interior point. Endpoint derivatives in $T_n$ use their continuous one-sided values. Then some $\xi$ strictly between $a$ and $x$ satisfies
\[ R_n(f,a)(x)=\frac{f^{(n+1)}(\xi)}{(n+1)!}(x-a)^{n+1}. \]
\end{theorem}
\begin{proof}
Define
\[ F(t)=f(t)-T_n(f,a)(t)-\frac{R_n(f,a)(x)}{(x-a)^{n+1}}(t-a)^{n+1}. \]
It has $F(a)=F(x)=0$ and $F^{(j)}(a)=0$ for $1\le j\le n$. Rolle gives a zero $c_1$ of $F'$ strictly between the endpoints. Apply Rolle to $F'$ between $a$ and $c_1$, using its zero value at $a$, to get a zero of $F''$; repeat. After $n+1$ applications there is an interior $\xi$ with $F^{(n+1)}(\xi)=0$. All intermediate functions are continuous on the relevant closed segments by $C^n$, and differentiable inside at the required orders. Substituting the derivative of $(t-a)^{n+1}$ gives the formula. The same nested-segment argument works when $x<a$ and retains the signed power $(x-a)^{n+1}$. For $n=0$ it is just the mean value theorem. As its $n=0$ case already fails for complex values by \KRef{K03.05.lagrange}, this single-point equality has no unrestricted complex version.
\end{proof}

\KEntry{K03.06.cauchy}{The real Cauchy remainder}
\begin{theorem}
Under the real hypotheses of \KRef{K03.06.lagrange}, an interior $\xi$ also gives
\[ R_n(f,a)(x)=\frac{f^{(n+1)}(\xi)}{n!}(x-\xi)^n(x-a). \]
The point $\xi$ may differ from the Lagrange point.
\end{theorem}
\begin{proof}
With $x$ fixed define $F(t)=\sum_{k=0}^n f^{(k)}(t)(x-t)^k/k!$ on the joining segment. It is continuous there, differentiable inside, and the product rule makes all neighboring terms cancel:
\[ F'(t)=\frac{f^{(n+1)}(t)}{n!}(x-t)^n. \]
Since $F(x)=f(x)$ and $F(a)=T_n(f,a)(x)$, the real mean value theorem gives $F(x)-F(a)=F'(\xi)(x-a)$. This theorem on the unordered segment yields exactly the displayed signed formula also for $x<a$; no orientation sign has been omitted.
\end{proof}

\KEntry{K03.06.integral}{The integral remainder}
\begin{theorem}
If $f\in C^{n+1}([u,v];\K)$ and $a,x\in[u,v]$, then
\[ R_n(f,a)(x)=\frac1{n!}\int_a^x(x-t)^n f^{(n+1)}(t)\dd t. \]
The integral is oriented, so the formula includes $x<a$ and $x=a$.
\end{theorem}
\begin{proof}
This result uses the later integral construction and the proved formulas \KRef{K03.09.ftc_second} and \KRef{K03.09.parts}; those results do not require Taylor's theorem. For $n=0$ it is the fundamental theorem, $\int_a^x f'=f(x)-f(a)$. For $n\ge1$ integration by parts gives
\[ \frac1{n!}\int_a^x(x-t)^n f^{(n+1)}(t)\dd t
=-\frac{(x-a)^n f^{(n)}(a)}{n!}+\frac1{(n-1)!}\int_a^x(x-t)^{n-1}f^{(n)}(t)\dd t. \]
Induction subtracts precisely the remaining Taylor coefficients. All integrands are continuous, so their proper integrals exist. Both integration by parts and the fundamental theorem hold with oriented endpoints; for complex values they hold in each real component, giving the stated scalar identity.
\end{proof}

\KEntry{K03.06.uniform_bound}{Uniform error estimates}
\begin{theorem}
Let $f:[u,v]\to\K$ belong to $C^n$ and have an interior $(n+1)$st derivative satisfying $|f^{(n+1)}(t)|\le M_n$. Then for all $a,x\in[u,v]$,
\[ |f(x)-T_n(f,a)(x)|\le\frac{M_n}{(n+1)!}|x-a|^{n+1}. \]
The constant is uniform on this interval for the stated order $n$; it may depend on $n$.
\end{theorem}
\begin{proof}
The real case follows from the Lagrange remainder. For complex values fix $a,x$ and let $z=R_n(f,a)(x)$. If $z=0$ there is nothing to prove. Otherwise put $w=z/|z|$ and apply the real remainder theorem to $F(t)=\RePart(\bar w f(t))$. Its Taylor remainder at $x$ is $|z|$, and $|F^{(n+1)}|\le|f^{(n+1)}|\le M_n$. This gives the same bound. In the stronger $C^{n+1}$ situation it also follows from the integral remainder and
\[ \left|\int_a^x(x-t)^n f^{(n+1)}(t)\dd t\right|\le M_n\int_{\min(a,x)}^{\max(a,x)}|x-t|^n\dd t=\frac{M_n|x-a|^{n+1}}{n+1}. \]
An assertion that remainders tend to zero as $n\to\infty$ needs control of the entire sequence $M_n$, not merely finiteness of each member.
\end{proof}

\KEntry{K03.06.leading_term}{The first nonzero derivative}
\begin{proposition}
Under the order-$n$ Peano hypotheses, $n\ge1$, suppose $f^{(j)}(a)=0$ for $1\le j<n$ and $c=f^{(n)}(a)/n!\ne0$. Then
\[ f(a+h)-f(a)=ch^n+o(h^n)\sim ch^n\quad(h\to0,\ h\ne0). \]
For real $f$, even $n$ and $c>0$ give a strict local minimum, even $n$ and $c<0$ a strict local maximum, and odd $n$ gives neither.
\end{proposition}
\begin{proof}
All lower Taylor coefficients vanish, so the expansion is the Peano formula. Division by $ch^n\ne0$ gives a ratio tending to one, proving the equivalence. For real values the remainder is smaller than $|c||h|^n/2$ for sufficiently small nonzero $h$. The increment therefore has the sign of $ch^n$. For even $n$ this sign is fixed on both sides; for odd $n$ the two sides have opposite signs.
\end{proof}

\KEntry{K03.06.composition}{Finite composition and division}
\begin{proposition}
Let $n\ge1$. Suppose $u(h)=b+P(h)+o(h^n)$ with a polynomial $P$ having zero constant term, and $F(b+z)=Q(z)+o(z^n)$ for a polynomial $Q$ of degree at most $n$ on the relevant scalar domain, with $F(b)=Q(0)$. Then $F(u(h))$ is obtained by substituting $P(h)$ into $Q$, discarding powers of $h$ above $n$, and adding $o(h^n)$. Likewise finite expansions through order $n$ can be divided whenever the denominator has a nonzero constant term. If all supplied remainders have the stronger next-order big-O bounds, the resulting remainder is $O(h^{n+1})$.
\end{proposition}
\begin{proof}
Put $z=u(h)-b=O(h)$. The outer remainder satisfies $|r(z)|\le\eps|z|^n$ for small $z$, with $r(0)=0$, hence is $o(h^n)$. For each fixed power, the factorization $z^j-P^j=(z-P)\sum_{\ell=0}^{j-1}z^\ell P^{j-1-\ell}$ shows that replacement by $P$ changes it by $o(h^n)$, since all factors are bounded near zero. Discarded polynomial terms of degrees at least $n+1$ are $O(h^{n+1})=o(h^n)$.
For division write the denominator $q(h)=q_0(1+v(h))$ with $q_0\ne0$ and $v(h)=O(h)$. Eventually $|v|<1/2$, and the exact finite identity
\[ \frac1{1+v}=\sum_{j=0}^n(-v)^j+\frac{(-v)^{n+1}}{1+v} \]
has final term $O(h^{n+1})$. Substitute the known polynomial for $v$ by the first argument, then multiply by the numerator and truncate. For next-order big-O input errors, every factorization error above is $O(h^{n+1})$, and the outer error is bounded by a constant times $|z|^{n+1}=O(h^{n+1})$; the stronger conclusion follows. These are finite polynomial operations, requiring no convergence theorem for an infinite series.
\end{proof}

\KEntry{K03.06.reference_expansions}{Reference finite expansions}
\begin{proposition}
For each fixed $n\ge0$, as $x\to0$,
\begin{align*}
\exp x&=\sum_{k=0}^n\frac{x^k}{k!}+O(x^{n+1}),\\
\log(1+x)&=\sum_{k=1}^n\frac{(-1)^{k+1}}k x^k+O(x^{n+1}),\\
(1+x)^\alpha&=\sum_{k=0}^n\frac{\alpha(\alpha-1)\cdots(\alpha-k+1)}{k!}x^k+O(x^{n+1})\quad(\alpha\in\R).
\end{align*}
The empty product is one. The logarithm and power formulas use $x>-1$. For fixed $m\ge0$,
\begin{align*}
\sin x&=\sum_{j=0}^m\frac{(-1)^jx^{2j+1}}{(2j+1)!}+O(x^{2m+3}),\\
\cos x&=\sum_{j=0}^m\frac{(-1)^jx^{2j}}{(2j)!}+O(x^{2m+2}).
\end{align*}
All constants can be chosen uniformly for $|x|\le r$ with a fixed $0<r<1$ for logarithms and powers, and any fixed $0<r<\infty$ for exponential and trigonometric functions.
\end{proposition}
\begin{proof}
Use the established derivatives from \KRef{K03.07.elementary_derivatives} and \KRef{K03.04.elementary_rules}. Repeated differentiation gives $\exp^{(k)}(0)=1$; for $k\ge1$, $(\log(1+x))^{(k)}=(-1)^{k-1}(k-1)!/(1+x)^k$; and the kth derivative of $(1+x)^\alpha$ is the falling product $\alpha(\alpha-1)\cdots(\alpha-k+1)(1+x)^{\alpha-k}$, proved by induction. On $[-r,r]$, the corresponding $(n+1)$st derivatives are bounded by $e^r$, $n!/(1-r)^{n+1}$, and
\[ |\alpha(\alpha-1)\cdots(\alpha-n)|\max\{(1-r)^{\alpha-n-1},(1+r)^{\alpha-n-1}\}, \]
respectively. The uniform remainder theorem gives the displayed orders (for the logarithm, even the explicit bound $|x|^{n+1}/((n+1)(1-r)^{n+1})$). Derivatives of sine and cosine cycle, their values at zero alternate as shown, and their moduli are at most one. Apply Taylor to degree $2m+2$ for sine, whose even coefficients vanish, and to degree $2m+1$ for cosine, whose odd coefficients vanish. This gives their stated next nonzero orders and uniform bounds. These finite applications occur after construction of the elementary functions and do not construct them circularly from Taylor's theorem.
\end{proof}

\KEntry{K03.06.smooth_flat}{A smooth function with every derivative zero at the origin}
\begin{construction}
Define $f(0)=0$ and $f(x)=\exp(-1/x^2)$ for $x\ne0$. Then $f\in C^\infty(\R)$, $f^{(j)}(0)=0$ for every integer $j\ge0$, and $f(x)>0$ whenever $x\ne0$.
\end{construction}
\begin{proof}
For any integer $N\ge0$,
\[ |x|^{-N}\exp(-1/x^2)\longrightarrow0\quad(x\to0), \]
because with $y=1/x^2\to\infty$ it is $y^{N/2}\exp(-y)$, which tends to zero by the proved growth estimate \KRef{K03.02.growth_scale}. Away from zero induction gives
\[ f^{(j)}(x)=P_j(1/x)\exp(-1/x^2),\quad P_0(t)=1,\quad P_{j+1}(t)=-t^2P_j'(t)+2t^3P_j(t). \]
The recurrence follows from the product and chain rules, since $(1/x)'=-1/x^2$ and $(-1/x^2)'=2/x^3$. Each $P_j$ is a polynomial. The decay estimate makes this expression tend to zero, so extending it by zero is continuous. Its difference quotient at zero is $P_j(1/h)\exp(-1/h^2)/h$, which also tends to zero by the same estimate. Inductively the extended jth expression is differentiable at zero with derivative zero, and its derivative away from zero is the next expression. Thus all derivatives exist and are continuous everywhere, with their stated values. Positivity away from zero is positivity of the exponential.
\end{proof}

\KEntry{K03.06.smooth_not_analytic}{Smoothness does not imply equality with the Taylor series}
\begin{example}
For the preceding function every Taylor coefficient at zero is zero. Its Taylor series is therefore the zero series, which converges for every real input; nevertheless it differs from $f$ at every nonzero input. Hence there is no neighborhood of zero on which $f$ equals its Taylor series. In the terminology of \KRef{K03.13.analyticity}, $f$ is smooth but not real analytic at zero. This example satisfies both existence of all derivatives and convergence of the Taylor series, while equality with the function still fails. A finite Peano estimate at every order therefore cannot, by itself, justify identifying a function with an infinite series: its neighborhood and constants may depend on the chosen order.
\end{example}

\KEntry{K03.06.error_notation}{Comparing error orders}
\begin{proposition}
As $h\to0$, $O(h^{n+1})\subset o(h^n)$ for every integer $n\ge0$, but the inclusion is strict.
\end{proposition}
\begin{proof}
If $|r(h)|\le C|h|^{n+1}$ eventually, then $|r(h)|/|h|^n\le C|h|\to0$ for $h\ne0$. For $0<\alpha<1$, the function $r(h)=|h|^{n+\alpha}$ instead satisfies $|r(h)|/|h|^n=|h|^\alpha\to0$, but $|r(h)|/|h|^{n+1}=|h|^{\alpha-1}\to\infty$. It is little-o of the nth power and not big-O of the next integer power. These assertions retain absolute values for odd powers and two-sided approaches.
\end{proof}

% BEGIN REVISED SOURCE chapters/M03.10.tex SHA256=6864dae43615d2de83f642a6679003c185d82f1de086f2921d83430dc9a5f72c
\chapter{Improper integrals and tails}
\SourceRef{Zorich I, Section 6.5, PDF pp. 413--421; Apostol, Section 10.13; Zorich II, Section 17.2 for the quantitative tail form}
All improper integrals below require proper Riemann integrability on each
regular compact truncation. Infinite endpoints and finite singular endpoints
are treated by limits, not by assigning an integrand value at a singularity.

\KEntry{K03.10.one_endpoint}{One singular endpoint}
\begin{definition}
Let $b\in\R\cup\{+\infty\}$ and $a<b$. If $f$ is properly integrable
on every $[a,t]$, $a<t<b$, define
\[\int_a^bf=\lim_{t\to b-}\int_a^tf
\]
only when this limit is finite in $\R$ or $\C$. For $b=+\infty$ the
notation means $t\to+\infty$. Define a singular left endpoint by the
corresponding right-hand limit, with orientation as in module 08.
Infinite divergence is a separate conclusion, not convergence to a scalar.
\end{definition}

\KEntry{K03.10.several_endpoints}{Independent singular endpoints}
\begin{definition}
For $f$ locally integrable on $(a,b)$, choose $c\in(a,b)$ and define
$\int_a^bf=\int_a^cf+\int_c^bf$ only if both terms converge
independently. If there are finitely many interior singularities, split
at each and require both adjacent one-sided integrals separately.
\end{definition}
\begin{proposition}
Existence and value are independent of all regular splitting points.
\end{proposition}
\begin{proof}
For $c<d$ in a regular region and a left truncation $u<c$,
$\int_u^df=\int_u^cf+\int_c^df$. The last term is a fixed proper
integral, so the limit for one splitting point exists if and only if
the other does, with their values differing by that proper integral.
The same calculation for right truncations cancels this difference in
the total. Repeat for every regular component. This argument never
links two singular truncation parameters, so it does not authorize
cancellation between divergent endpoint terms.
\end{proof}

\KEntry{K03.10.cauchy}{The Cauchy tail criterion}
\begin{theorem}
$\int_a^bf$ converges if and only if for every $\eps>0$ there is a
regular $r<b$ such that
\[\left|\int_u^vf\right|<\eps\qquad(r<u<v<b).
\]
At infinity replace $r<b$ by a sufficiently large $r$.
\end{theorem}
\begin{proof}
Write $H(t)=\int_a^tf$. Then $\int_u^vf=H(v)-H(u)$ by proper
additivity. A finite limit of $H$ makes these differences small.
Conversely choose any increasing cofinal sequence $t_n\to b$.
The tail condition makes $H(t_n)$ Cauchy, hence convergent to a scalar
$L$. Given $\eps$, choose a threshold for differences below $\eps/2$
and then a $t_n$ beyond it with $|H(t_n)-L|<\eps/2$. Every sufficiently
late $t$ has $|H(t)-H(t_n)|<\eps/2$ (reverse the order if necessary),
so $H(t)\to L$. Left endpoint versions follow by reversing orientation.
\end{proof}

\KEntry{K03.10.nonnegative}{Nonnegative improper integrals}
\begin{theorem}
For real locally integrable $f\ge0$ on $[a,b)$, the partial integrals
$H(t)=\int_a^tf$ increase. They converge finitely if and only if they
are bounded above; otherwise they tend to $+\infty$.
\end{theorem}
\begin{proof}
Order and additivity give $H(v)-H(u)=\int_u^vf\ge0$.
If the set of partial values has finite supremum $L$, choose a partial
value above $L-\eps$; all later values lie in $[L-\eps,L]$, proving
convergence. If unbounded, for each $A$ choose one value above $A$;
all later values stay above $A$, proving divergence to $+\infty$.
\end{proof}

\KEntry{K03.10.comparison}{Comparison and limit comparison}
\begin{theorem}
For eventually $0\le f\le g$, convergence of $\int g$ implies
convergence of $\int f$. If $f,g$ are eventually positive and
$f(x)/g(x)\to c\in(0,\infty)$ at the singular endpoint, their
improper integrals converge or diverge together.
\end{theorem}
\begin{proof}
Discard a regular initial segment, which contributes only a finite proper
value. Order bounds partial integrals of $f$ by those of $g$, so the
nonnegative boundedness criterion proves the first assertion. For the
second, sufficiently late points satisfy
$(c/2)g\le f\le(3c/2)g$. Apply comparison in both directions.
Eventual positivity is needed for this order argument; the ratio test
does not make the same conclusion for arbitrary sign-changing functions.
\end{proof}

\KEntry{K03.10.absolute}{Absolute and conditional convergence}
\begin{definition}
An improper integral is absolutely convergent if $\int|f|$ converges.
It is conditionally convergent if $\int f$ converges and $\int|f|$
does not. These definitions apply separately at every singular endpoint.
\end{definition}
\begin{theorem}
Absolute convergence implies convergence for real or complex scalars,
and $|\int f|\le\int|f|$.
\end{theorem}
\begin{proof}
On every regular compact interval $|f|$ is integrable by module 08 and
$|\int_u^vf|\le\int_u^v|f|$. Cauchy tails of the latter therefore
give Cauchy tails of the former. Scalar completeness proves convergence;
passing to the limit in the finite norm inequality proves the value bound.
The oscillatory example below proves that the converse fails.
\end{proof}

\KEntry{K03.10.p_test}{Power and logarithmic tests}
\begin{theorem}
For real $p$,
\[\int_1^\infty x^{-p}\dd x<\infty\iff p>1,\qquad
\int_0^1 x^{-p}\dd x<\infty\iff p<1.
\]
For $A>1$, $\int_A^\infty\dd x/(x(\log x)^q)$ converges exactly
when $q>1$.
\end{theorem}
\begin{proof}
For $p\ne1$ a truncation from $u$ to $v$ equals
$(v^{1-p}-u^{1-p})/(1-p)$, by the proved power derivative and
Newton--Leibniz. At infinity $v^{1-p}$ tends to zero for $p>1$
and to infinity for $p<1$; the resulting finite value for $p>1$ is
$1/(p-1)$. At zero $u^{1-p}$ tends to zero exactly for $p<1$,
giving $1/(1-p)$, and diverges for $p>1$.
For $p=1$, truncations are $\log v-\log u$, divergent at both
relevant endpoints. For the logarithmic test make the proper monotone
substitution $y=\log x$, $\dd y=\dd x/x$ on each $[A,R]$.
Since $\log R\to\infty$ this becomes the power test
$\int_{\log A}^\infty y^{-q}\dd y$. The starting point $A>1$
ensures positive logarithm and no extra singularity at $1$.
\end{proof}

\KEntry{K03.10.integral_test}{Series and integrals}
\begin{theorem}
If $f:[1,\infty)\to[0,\infty)$ is nonincreasing, then
$\sum_{n=1}^{\infty}f(n)$ converges if and only if $\int_1^\infty f$
does. If they converge, for every integer $N\ge1$,
\[\int_{N+1}^{\infty}f\le\sum_{n=N+1}^{\infty}f(n)
\le\int_N^{\infty}f.
\]
\end{theorem}
\begin{proof}
Monotone bounded functions are integrable on each finite interval.
For integer $n$, monotonicity gives
$f(n+1)\le\int_n^{n+1}f\le f(n)$.
Sum over consecutive unit intervals. Boundedness of the nonnegative
partial sums is equivalent to boundedness of partial integrals (an
arbitrary truncation lies between two neighboring integer truncations).
The respective monotone convergence criteria prove equivalence.
Summing the shifted inequalities from $N$ to a finite upper integer
and then taking its limit gives both stated remainder bounds.
\end{proof}

\KEntry{K03.10.dirichlet}{Dirichlet's oscillation test}
\begin{lemma}[Proper parts formula for an accumulation primitive]
If $f$ is Riemann integrable on $[r,s]$, $F(x)=\int_r^xf$, and
$g\in C^1([r,s])$, then
\[\int_r^sfg=F(s)g(s)-\int_r^sFg'.
\]
\end{lemma}
\begin{proof}
First suppose $g$ is real-valued; $f$ and $F$ may still be complex-valued. For a partition write $F_i=F(x_i)$, $g_i=g(x_i)$.
The difference between $\int fg$ and
$\sum_i g_{i-1}(F_i-F_{i-1})$ is at most
$\|f\|_\infty(s-r)\omega_g(|P|)\to0$.
Finite summation by parts writes the latter as
$F(s)g(s)-\sum_iF_i(g_i-g_{i-1})$, since $F(r)=0$.
By the mean value theorem $g_i-g_{i-1}=g'(\eta_i)\Delta x_i$.
Replace $F_i$ by $F(\eta_i)$ with error at most
$(s-r)\|g'\|_\infty\omega_F(|P|)\to0$.
The remaining sum tends to $\int Fg'$, proving the formula.
For complex-valued $g=u+iv$, apply the proved real-$g$ identity separately to $u$ and $v$ and combine the two identities by complex linearity of the proper integral. This supplies the complex case without a false exact complex mean-value tag.
This proof is necessary because $F'=f$ need not hold at discontinuities
of $f$; a merely integrable $f$ does not make $F$ a $C^1$ primitive.
\end{proof}
\begin{theorem}
Let $f$ be locally Riemann integrable on $[a,\infty)$ with bounded
accumulation $A(x)=\int_a^xf$, $|A(x)|\le C$. If real $g$ is $C^1$,
monotone and tends to zero, then $\int_a^\infty fg$ converges.
\end{theorem}
\begin{proof}
A monotone function tending to zero has constant sign, and $h=|g|$
is nonincreasing with $h\to0$. For $R<S$ the primitive
$F_R(x)=\int_R^xf=A(x)-A(R)$ satisfies $|F_R|\le2C$.
The lemma, with $h'\le0$, gives
\[\left|\int_R^Sfh\right|
\le2Ch(S)+2C\int_R^S(-h')=2Ch(R).
\]
The right side tends to zero uniformly for $S\ge R\to\infty$,
so the Cauchy criterion proves convergence of $fh$ and hence $fg$.
The proof also works at a finite singular endpoint with $R,S$ tending
to that endpoint in the specified direction.
\end{proof}

\KEntry{K03.10.abel}{Abel's test}
\begin{theorem}
If $\int_a^\infty f$ converges and real $g$ is bounded, $C^1$ and
monotone, then $\int_a^\infty fg$ converges.
\end{theorem}
\begin{proof}
Let $|g|\le B$. For given $\eps>0$, the Cauchy criterion gives
$|F_R(x)|=|\int_R^xf|<\eps$ for all $x\ge R$ and all sufficiently
large $R$. The proper accumulation parts formula gives
\[\left|\int_R^Sfg\right|
\le\eps\left(|g(S)|+\int_R^S|g'|\right)
=\eps(|g(S)|+|g(S)-g(R)|)\le3B\eps.
\]
This proves Cauchy tails. For $B=0$ the product is zero.
Bounded monotonicity also gives a finite limit $\ell$ of $g$; the
identity $fg=\ell f+f(g-\ell)$ explains the distinct limits, but the
displayed estimate proves convergence without an unproved interchange.
\end{proof}

\KEntry{K03.10.parts}{Improper integration by parts}
\begin{theorem}
On regular truncations let the proper parts formula hold for $u,v$.
If $u(t)v(t)$ has a finite singular-endpoint limit and either of
$\int uv'$ and $\int u'v$ converges, then the other converges and
\[\int_a^buv'=[uv]_a^b-\int_a^bu'v.
\]
For several singular endpoints require these limits separately at each.
\end{theorem}
\begin{proof}
For regular $t<b$, the proper formula is
$\int_a^tuv'=u(t)v(t)-u(a)v(a)-\int_a^tu'v$.
If the boundary product and one integral have finite limits, this
identity expresses the remaining partial integral as a difference of
convergent quantities and proves its finite limit and the formula.
Repeat separately at a singular left endpoint or on either side of an
interior singularity. If the product has no justified limit, subtracting
formal infinite boundary terms is not a valid application of this theorem.
\end{proof}

\KEntry{K03.10.substitution}{Improper substitutions}
\begin{theorem}
Suppose $\phi:[\alpha,\beta)\to[a,b)$ is a monotone $C^1$
parameterization with $\phi(\alpha)=a$, $\phi(t)\to b$ as
$t\to\beta-$, and is onto the indicated interval. If $f$ is locally
properly integrable, then
\[\int_a^bf=\int_\alpha^\beta f(\phi(t))\phi'(t)\dd t
\]
whenever either improper limit converges, and existence of one is
equivalent to existence of the other. Orientation handles decreasing
parameterizations and left endpoints.
\end{theorem}
\begin{proof}
On each regular $[\alpha,T]$ the monotone substitution theorem proves
integrability of the transformed product and gives
$\int_\alpha^T(f\circ\phi)\phi'=\int_a^{\phi(T)}f$.
Cofinality $\phi(T)\to b$ proves the transformed limit if the original
limit exists. Conversely, for each regular $r\in[a,b)$ choose
$T_r$ with $\phi(T_r)=r$, using surjectivity; as $r\to b$, such
$T_r$ leave every compact subinterval of $[\alpha,\beta)$, since its
image has maximum strictly below $b$. Thus $T_r\to\beta$ and the
same identity proves the original limit. Flat portions cause no problem.
A nonmonotone parameterization is not covered by this cofinal argument.
\end{proof}

\KEntry{K03.10.principal_value}{Symmetric principal values}
\begin{definition}
At an interior singularity $c$ define, when finite,
\[\operatorname{pv}\int_a^bf
=\lim_{\eps\downarrow0}\left(\int_a^{c-\eps}f+
\int_{c+\eps}^bf\right),
\]
assuming the other endpoints are regular or separately controlled.
On the whole line a symmetric value is
$\lim_{R\to\infty}\int_{-R}^Rf$, with any interior singularities
handled explicitly. The center and linked truncation rule are part of
the definition. Independent improper convergence implies the matching
principal value by addition of limits, but not conversely.
\end{definition}

\KEntry{K03.10.pv_counterexample}{A principal value without an improper integral}
\begin{example}
For $f(x)=1/x$ on $[-1,1]\setminus\{0\}$,
\[\int_{-1}^{-\eps}\frac{\dd x}{x}=\log\eps,\qquad
\int_\eps^1\frac{\dd x}{x}=-\log\eps.
\]
Their sum is zero for each $\eps$, so the symmetric principal value
is zero. The two terms separately tend to $-\infty$ and $+\infty$,
so the ordinary improper integral does not exist.
\end{example}

\KEntry{K03.10.oscillatory_example}{A conditionally convergent oscillatory integral}
\begin{example}
$\int_1^\infty\sin x\,\dd x/x$ converges, but not absolutely.
\end{example}
\begin{proof}
The primitive $\int_1^x\sin t\dd t=\cos1-\cos x$ is bounded by two,
and $1/x$ decreases differentiably to zero, so Dirichlet applies.
Since $\sin(\pi/2)=1$ and sine is continuous, choose
$0<\delta<\pi/2$ such that $\sin x\ge1/2$ on
$[\pi/2-\delta,\pi/2+\delta]$. For integer $k\ge1$, translate this
interval by $k\pi$ to obtain $I_k$; the sign-change identity for sine
gives $|\sin x|\ge1/2$ on every $I_k$. Thus
\[\int_{I_k}\frac{|\sin x|}{x}\dd x
\ge\frac{\delta}{(k+1)\pi}.
\]
The harmonic series diverges (group terms between consecutive powers of
two, each group having sum at least $1/2$). Partial integrals of the
absolute integrand therefore grow without bound. This proves conditional
convergence by an actual recurring-interval lower bound.
\end{proof}

\KEntry{K03.10.tail_bounds}{Quantitative oscillatory remainders}
\begin{theorem}
If $|A(x)|=|\int_a^xf|\le C$ and real $g\ge0$ is $C^1$,
nonincreasing and tends to zero, then for $S\ge R\ge a$,
\[\left|\int_R^Sfg\right|\le2Cg(R),\qquad
\left|\int_R^\infty fg\right|\le2Cg(R).
\]
\end{theorem}
\begin{proof}
The proof of Dirichlet already gives the first bound, with primitive
$A(x)-A(R)$ bounded by $2C$. Dirichlet proves the infinite integral
exists, so letting $S\to\infty$ proves the second bound by continuity
of modulus. If instead one assumes directly
$|\int_R^xf|\le C_R$ for all $x\ge R$, the same calculation improves
both bounds to $C_Rg(R)$. These estimates specify the tail, rather than
merely invoking convergence by name, and will be used uniformly in
parameters in module 14.
\end{proof}

\KEntry{K03.10.endpoint_counterexample}{Linked endpoint truncations are insufficient}
\begin{example}
On $(0,1)$ let $f(x)=1/(1-x)-1/x$. A primitive on this interval is
$H(x)=-\log(1-x)-\log x$. Since $H(1-\eps)=H(\eps)$,
$\int_\eps^{1-\eps}f=0$ for every $0<\eps<1/2$.
Nevertheless
$\int_\eps^{1/2}f=H(1/2)-H(\eps)\to-\infty$,
and
$\int_{1/2}^{1-\eps}f=H(1-\eps)-H(1/2)\to+\infty$.
Both independent endpoint conditions fail. A finite linked limit on two
moving endpoints therefore does not establish ordinary improper convergence.
\end{example}

% BEGIN REVISED SOURCE chapters/M03.14.tex SHA256=252f7e805550cce6ae7b5539192b61da8ff85423bc9f60dcc3474c61b75172db
\chapter{Integrals depending on a parameter}
\SourceRef{Zorich II, Sections 17.1--17.2, PDF pp. 421--442, especially the uniform difference-quotient proof and the base-point version of improper differentiation; Apostol, Section 7.23}
The integration variable is real and the parameter is real. Scalar values may
be real or complex, with all complex derivative estimates understood through
the real-parameter norm inequality. The results here concern proper Riemann
integrals and uniform improper tails; no Lebesgue dominated-convergence or
general Fubini theorem is an implicit prerequisite.

\KEntry{K03.14.proper_definition}{A proper integral family}
\begin{definition}
Let $T\subset\R$ and let $f:[a,b]\times T\to\K$. If each section
$x\mapsto f(x,t)$ is Riemann integrable, define
$F(t)=\int_a^bf(x,t)\dd x$. The formula defines $F$ only on the
parameters where that section is integrable; the symbol alone supplies
no existence or parameter regularity.
\end{definition}

\KEntry{K03.14.proper_continuity}{Continuity under a finite integral}
\begin{theorem}
Suppose the sections are integrable near $t_0$ and
$\sup_{x\in[a,b]}|f(x,t)-f(x,t_0)|\to0$ as $t\to t_0$ in $T$.
Then $F$ is continuous at $t_0$. In particular joint continuity on a
compact rectangle $[a,b]\times[c,d]$ gives a continuous $F$ on $[c,d]$.
\end{theorem}
\begin{proof}
The finite norm estimate gives
\[|F(t)-F(t_0)|\le(b-a)\sup_x|f(x,t)-f(x,t_0)|\longrightarrow0.
\]
For joint continuity, compact uniform continuity gives the stated
uniform-in-$x$ estimate whenever $|t-t_0|$ is sufficiently small.
The same compactness gives continuous, hence integrable, $x$-sections.
At parameter endpoints continuity is relative to the parameter interval.
\end{proof}

\KEntry{K03.14.proper_derivative}{Differentiation under a finite integral}
\begin{theorem}
Let $J$ be an open parameter interval. Suppose $f$ and its real partial
derivative $f_t$ are jointly continuous on $[a,b]\times J$.
Then $F\in C^1(J)$ and
\[F'(t)=\int_a^bf_t(x,t)\dd x.
\]
It is enough to have these conditions on every compact parameter
neighborhood on which the conclusion is used.
\end{theorem}
\begin{proof}
Fix $t_0$ and a compact neighborhood $[c,d]\subset J$.
Uniform continuity of $f_t$ there gives a modulus $\omega(\delta)\to0$
for changes of parameter distance at most $\delta$, uniformly in $x$.
For $h\ne0$ with $t_0+h\in[c,d]$, the norm mean-value inequality
applied to $s\mapsto f(x,s)-s f_t(x,t_0)$ gives
\[\left|\frac{f(x,t_0+h)-f(x,t_0)}h-f_t(x,t_0)\right|
\le\sup_{s\text{ between }t_0,t_0+h}|f_t(x,s)-f_t(x,t_0)|
\le\omega(|h|).
\]
This is valid for complex scalars as a norm inequality, without claiming
a single complex Lagrange equality. Integrate and use the finite norm
bound to obtain difference-quotient error at most $(b-a)\omega(|h|)$.
It tends to zero, proving the derivative formula. The preceding continuity
theorem applied to $f_t$ proves continuity of the resulting derivative.
\end{proof}

\KEntry{K03.14.moving_endpoints}{Leibniz's rule with moving endpoints}
\begin{theorem}
Let $\alpha,\beta\in C^1(J,\R)$. Near a parameter $t_0$, suppose
$f,f_t$ are jointly continuous on a rectangle containing all interval
endpoints $\alpha(t),\beta(t)$ and the intervals between them. Then
\begin{align*}
\frac{\dd}{\dd t}\int_{\alpha(t)}^{\beta(t)}f(x,t)\dd x
={}&f(\beta(t),t)\beta'(t)-f(\alpha(t),t)\alpha'(t)\\
&+\int_{\alpha(t)}^{\beta(t)}f_t(x,t)\dd x.
\end{align*}
Orientation permits the two endpoints to cross.
\end{theorem}
\begin{proof}
Choose a fixed regular $x_*$ in the spatial interval and set
$H(y,t)=\int_{x_*}^yf(x,t)\dd x$. The first fundamental theorem gives
$H_y(y,t)=f(y,t)$, and the preceding differentiation theorem gives
$H_t(y,t)=\int_{x_*}^yf_t(x,t)\dd x$ for fixed $y$.
For completeness along a moving endpoint $y(t)$, split
\begin{align*}
H(y(t+h),t+h)-H(y(t),t)
={}&\int_{y(t)}^{y(t+h)}f(x,t+h)\dd x\\
&+\int_{x_*}^{y(t)}(f(x,t+h)-f(x,t))\dd x.
\end{align*}
The first quotient tends to $f(y(t),t)y'(t)$: subtract that constant
integrand and bound the error by
$|y(t+h)-y(t)|/|h|$ times the local uniform continuity error.
The second quotient tends to $\int_{x_*}^{y(t)}f_t$ by the uniform
difference-quotient estimate. Subtract these formulas for $y=\beta$
and $y=\alpha$ and use oriented additivity. This establishes the rule
using one-dimensional estimates only.
\end{proof}

\KEntry{K03.14.proper_order}{Interchanging two finite integrations}
\begin{theorem}
For continuous scalar $f$ on $[a,b]\times[c,d]$ both iterated proper
integrals exist and
\[\int_c^d\left(\int_a^bf(x,t)\dd x\right)\dd t
=\int_a^b\left(\int_c^df(x,t)\dd t\right)\dd x.
\]
\end{theorem}
\begin{proof}
The inner sections are integrable and their integral functions are
continuous by the continuity theorem, so the outer integrals exist.
Let $\omega(\delta)$ be a uniform continuity modulus on the rectangle.
Take tagged partitions $P,Q$ in $x,t$ with meshes below $\delta$ and
form the finite double sum
$D=\sum_{i,j}f(\xi_i,\tau_j)\Delta x_i\Delta t_j$.
For every fixed $t$, replacing $\int_a^bf(x,t)\dd x$ by
$\sum_i f(\xi_i,t)\Delta x_i$ has error at most
$(b-a)\omega(\delta)$, uniformly in $t$, by integrating the
cellwise difference. Integrating that error in $t$, and then replacing
each $\int_c^df(\xi_i,t)\dd t$ by its tagged sum, shows the first
iterated integral differs from $D$ by at most
$2(b-a)(d-c)\omega(\delta)$. Reverse the order to obtain the same
bound for the other iterated integral. Their difference tends to zero
as $\delta\downarrow0$. This is an equality of limits of finite sums,
not an application of a general multiple-integration theorem.
\end{proof}

\KEntry{K03.14.uniform_improper}{Uniform improper convergence}
\begin{definition}
Assume local proper integrability in $x$ for every $t\in T$ and write
$F_R(t)=\int_a^Rf(x,t)\dd x$. The family converges uniformly
improperly on $T$ if there is $F:T\to\K$ with
$\sup_{t\in T}|F_R(t)-F(t)|\to0$ as $R\to\infty$.
Equivalently, for every $\eps>0$ there is $R_0$ such that
\[\left|\int_R^Sf(x,t)\dd x\right|<\eps
\quad(t\in T,\ S\ge R\ge R_0).
\]
For a finite right singular endpoint $b$, require the same condition
for $r_0<R<S<b$ with $r_0$ sufficiently close to $b$; for a left
endpoint reverse the truncations. Several singularities require
independent uniform conditions at every endpoint.
\end{definition}
\begin{proof}[Equivalence]
Uniform convergence implies the displayed estimate by the triangle
inequality. Conversely, for each fixed $t$ the ordinary Cauchy criterion
gives a finite limit $F(t)$. Hold $R\ge R_0$ and let $S\to\infty$
in the displayed inequality. This gives $|F(t)-F_R(t)|\le\eps$
simultaneously for all $t\in T$. Apply the condition with an arbitrarily
smaller tolerance to obtain uniform convergence. The same argument works
at a finite singular endpoint.
\end{proof}

\KEntry{K03.14.uniform_comparison}{Uniform domination of tails}
\begin{theorem}
Suppose $g\ge0$ is locally properly integrable and $\int_A^\infty g$
converges. If $|f(x,t)|\le g(x)$ for $x\ge A$ and all $t\in T$,
and each $f$ section is locally properly integrable, then its improper
integrals converge uniformly on $T$, with
\[\sup_{t\in T}\left|\int_R^\infty f(x,t)\dd x\right|
\le\int_R^\infty g(x)\dd x\quad(R\ge A).
\]
\end{theorem}
\begin{proof}
On every regular $[R,S]$, the finite norm and order inequalities give
$|\int_R^Sf(x,t)\dd x|\le\int_R^Sg$ for every $t$.
The right-hand tails are uniformly small in $S\ge R$ by convergence
of $\int g$, so the uniform Cauchy criterion applies. Let $S\to\infty$
in the finite inequality to obtain the quantitative bound. Integrals
over the fixed initial interval exist separately by local integrability;
they need not obey a single domination bound to prove tail convergence.
Finite-endpoint versions have exactly the same proof on the corresponding
truncations.
\end{proof}

\KEntry{K03.14.uniform_dirichlet}{Uniform Dirichlet test}
\begin{theorem}
Let $f(\cdot,t)$ be locally Riemann integrable and
$A(x,t)=\int_a^xf(u,t)\dd u$. Suppose $|A(x,t)|\le C$ for all
$x\ge a$, $t\in T$. Let $g(\cdot,t)\ge0$ be $C^1$ and
nonincreasing, and suppose $\sup_{t\in T}g(R,t)\to0$.
Then $\int_a^\infty f(x,t)g(x,t)\dd x$ converges uniformly on $T$.
\end{theorem}
\begin{proof}
Products are locally integrable. The finite accumulation-parts proof
in \KRef{K03.10.dirichlet} works for each fixed $t$, with the same
constant: the primitive anchored at $R$ has modulus at most $2C$.
Thus
\[\left|\int_R^Sf(x,t)g(x,t)\dd x\right|
\le2C g(R,t)\le2C\sup_{u\in T}g(R,u).
\]
The right side tends to zero independently of $S\ge R$ and $t$.
The uniform Cauchy criterion proves the claim. Signed amplitudes of a
fixed sign are covered by replacing $g$ by its modulus when that modulus
has the same monotonicity and uniform limit property.
\end{proof}

\KEntry{K03.14.improper_continuity}{Continuity from uniform tails}
\begin{theorem}
If every proper truncation $F_R$ is continuous at $t_0$, and the
improper tails converge uniformly on a relative neighborhood $T_0$
of $t_0$, then $F(t)=\int_a^\infty f(x,t)\dd x$ is continuous at $t_0$.
\end{theorem}
\begin{proof}
Choose $R$ with $\sup_{t\in T_0}|F(t)-F_R(t)|<\eps/3$.
Continuity of this fixed $F_R$ gives
$|F_R(t)-F_R(t_0)|<\eps/3$ for all sufficiently close $t\in T_0$.
The triangle inequality gives
\[|F(t)-F(t_0)|\le|F(t)-F_R(t)|+
|F_R(t)-F_R(t_0)|+|F_R(t_0)-F(t_0)|<\eps.
\]
Joint continuity on each regular compact rectangle is one sufficient
condition for the truncation continuity, by the proper theorem.
\end{proof}

\KEntry{K03.14.improper_derivative}{Differentiation from uniform derivative tails}
\begin{theorem}
Let $[c,d]$ be a compact parameter interval. Suppose every truncation
$F_R$ is continuous on $[c,d]$ and differentiable inside, and proper
differentiation is valid there:
\[F_R'(t)=\int_a^R f_t(x,t)\dd x.
\]
Assume these derivative integrals converge uniformly on $[c,d]$ to
$G(t)$ (using endpoint extensions where needed), and $F_R(t_*)$ has a
finite limit at one $t_*\in[c,d]$. Then $F_R$ converges uniformly on
$[c,d]$ to an improper integral family $F$, differentiable inside with
$F'=G$. Joint continuity of $f,f_t$ on every regular compact rectangle
is a sufficient condition for the proper differentiation assumptions.
\end{theorem}
\begin{proof}
For any two late truncations, the norm mean-value inequality gives
\[\sup_{t\in[c,d]}|F_R(t)-F_S(t)|
\le|F_R(t_*)-F_S(t_*)|+(d-c)\sup_{u\in[c,d]}|F_R'(u)-F_S'(u)|.
\]
Both terms tend to zero, so completeness pointwise and the same uniform
Cauchy argument as above yield a uniform limit $F$. In particular all
original improper integrals now exist, not just the one at $t_*$.
Fix an interior $t$ and a nonzero $h$ with $t+h\in[c,d]$.
Apply the norm mean-value inequality to $F_S-F_R$ between these points
and let $S\to\infty$. Uniform convergence of derivatives gives
\[\left|\frac{F(t+h)-F(t)}h-
\frac{F_R(t+h)-F_R(t)}h\right|
\le\sup_{u\in[c,d]}|G(u)-F_R'(u)|.
\]
Hence the difference of the first quotient from $G(t)$ is at most this
uniform error plus
$|(F_R(t+h)-F_R(t))/h-F_R'(t)|+|F_R'(t)-G(t)|$.
Choose $R$ first to make the first and third terms small, then let
$h\to0$ for that fixed differentiable $F_R$. This proves $F'=G$.
It is the derivative-limit argument of \KRef{K03.12.differentiation}
spelled out for a continuous truncation parameter; uniform convergence
of the original integrands alone would not justify this conclusion.
\end{proof}

\KEntry{K03.14.parameter_order}{Integrating an improper family in the parameter}
\begin{theorem}
Suppose $f$ is jointly continuous on $[a,R]\times[c,d]$ for every
finite $R>a$, and $\int_a^\infty f(x,t)\dd x$ converges uniformly
for $t\in[c,d]$. Then
\[\int_c^d\left(\int_a^\infty f(x,t)\dd x\right)\dd t
=\int_a^\infty\left(\int_c^df(x,t)\dd t\right)\dd x,
\]
with a proper outer integral on the left and a convergent improper
outer integral on the right.
\end{theorem}
\begin{proof}
The truncations $F_R$ are continuous and converge uniformly, so $F$ is
continuous and properly integrable on $[c,d]$ by
\KRef{K03.08.uniform_stability}. Let
$H(x)=\int_c^df(x,t)\dd t$. Proper parameter continuity shows $H$
is continuous on each regular $[a,R]$, hence locally properly integrable.
For every $R$, the proper order theorem gives
$\int_a^RH(x)\dd x=\int_c^dF_R(t)\dd t$.
The right side tends to $\int_c^dF$ with error at most
$(d-c)\|F_R-F\|_\infty$. Thus the left partial integrals have a
finite limit, proving existence of the improper right side and the
claimed equality simultaneously. More generally the same proof works
if truncation order interchange and integrability are established by
other explicit proper hypotheses.
\end{proof}

\KEntry{K03.14.local_uniform}{Local uniformity is sufficient}
\begin{proposition}
On an open parameter interval $J$, uniform tail convergence on each
compact parameter interval implies continuity whenever truncations are
continuous. Uniform convergence of the derivative-integral tails on
each compact interval, with convergence of the original family at one
point of each such interval, gives differentiability there. If $J$ is
an interval and convergence is known at just one point, the base-point
condition propagates to all compact subintervals by overlapping intervals.
\end{proposition}
\begin{proof}
Every $t_0\in J$ is in the interior of some $[c,d]\subset J$.
Apply the preceding continuity or differentiation theorem on that
compact neighborhood. On overlaps the limiting functions agree because
they are limits of the same truncations; so do their derivative formulas.
For propagation from a fixed $t_*$ to any compact $K\subset J$, the
closed interval spanning $K\cup\{t_*\}$ lies in $J$. Uniform derivative
tails there and the one base point give convergence throughout by the
preceding theorem. No single bound uniform over an unbounded $J$ is
required for these local conclusions.
\end{proof}

\KEntry{K03.14.failure}{Pointwise convergence without parameter continuity}
\begin{example}
For $t\ge0$, $x\ge0$, put $f(x,t)=tE(-tx)$, so $f(x,0)=0$.
For each fixed $x$ this is continuous at zero (indeed jointly continuous
on finite rectangles). For $t>0$,
\[F_R(t)=\int_0^RtE(-tx)\dd x=1-E(-tR),\qquad
F(t)=1.
\]
At $t=0$ every partial integral and the improper value are zero, so $F$
is discontinuous. The tail for $t>0$ equals $E(-tR)$, whose supremum
over $0<t\le\delta$ is one for every $R$. Thus uniform convergence
fails on every neighborhood of zero, exactly where the continuity
theorem requires it. The finite computations use the proved exponential
derivative and the infinite computation uses $E(-tR)\to0$ for fixed $t>0$.
\end{example}

\KEntry{K03.14.boundary_terms}{A moving boundary supplies the derivative}
\begin{example}
Let $f(x,t)=1$ and $F(t)=\int_0^t1\dd x=t$, using orientation for
$t<0$. Thus $F'(t)=1$. The integrand parameter derivative is zero;
Leibniz's rule supplies the value entirely from the upper boundary term
$f(t,t)\,\frac{\dd t}{\dd t}=1$. Omitting that term would give the
incorrect derivative zero, even for this continuous constant integrand.
\end{example}

% BEGIN REVISED SOURCE chapters/M03.15.tex SHA256=b9bb3e206744bea87048824cb488bf5c1ca2780a44eceae9cf9c6341524a55c4
\chapter{Bounded variation and the refinement Stieltjes integral}
\SourceRef{Apostol, Chapters 6--7, especially Definition 7.1 (PDF 154--155); Rudin, Chapter 6. The refinement definition is retained exactly. Its comparison with a mesh definition has a genuine continuity hypothesis on the integrator.}

\KEntry{K03.15.variation}{Total variation}
\begin{definition}
For $f:[a,b]\to\K$, $a<b$, put
\[\Var_{[a,b]}f=\sup_{a=x_0<\cdots<x_m=b}\sum_{i=1}^m\abs{f(x_i)-f(x_{i-1})}\in[0,\infty].\]
The function has bounded variation, abbreviated BV, when this is finite. Set $\Var_{[a,a]}f=0$. Values at endpoints and at individual interior points are part of the definition; changing a single value can change variation. For example, a function zero except for value $1$ at an interior point has variation $2$.
\end{definition}

\KEntry{K03.15.variation_rules}{Variation rules and boundedness}
\begin{theorem}
For $a\leq c\leq b$, variation is additive:
\[\Var_{[a,b]}f=\Var_{[a,c]}f+\Var_{[c,b]}f.\]
Also $\Var(f+g)\leq\Var f+\Var g$ and $\Var(\lambda f)=\abs\lambda\Var f$ for finite variation (with the zero scalar case interpreted directly). A BV function satisfies $\abs{f(x)}\leq\abs{f(a)}+\Var f$.
\end{theorem}
\begin{proof}
Inserting $c$ into any partition cannot decrease its sum, by the triangle inequality. The refined sum is the sum of two partition sums on the subintervals, proving the upper bound for additivity. Conversely combine two partitions approaching the respective suprema; their union gives the reverse bound, including the infinite case. The triangle inequality on each increment proves the sum rule, and homogeneity on each increment proves the scalar rule. The partition containing $a,x,b$ bounds $\abs{f(x)-f(a)}$ by the total variation, proving boundedness.
\end{proof}

\KEntry{K03.15.monotone}{Monotone functions have bounded variation}
\begin{theorem}
For real monotone $f$, $\Var f=\abs{f(b)-f(a)}$. Consequently finite real or complex linear combinations of BV functions remain BV.
\end{theorem}
\begin{proof}
For nondecreasing $f$ each increment is nonnegative, so every partition sum telescopes to $f(b)-f(a)$. For a nonincreasing function apply this to $-f$. The assertion about linear combinations follows repeatedly from the preceding sum and scalar bounds.
\end{proof}

\KEntry{K03.15.jordan}{Jordan decomposition of real BV functions}
\begin{theorem}
A real function $f$ on $[a,b]$ is BV if and only if it is the difference of two nondecreasing real functions. One choice is
\[V(x)=\Var_{[a,x]}f,\qquad p(x)=\tfrac12(V(x)+f(x)-f(a)),\quad q(x)=\tfrac12(V(x)-f(x)+f(a)),\]
so $f=f(a)+p-q$.
\end{theorem}
\begin{proof}
For $y>x$, $V(y)-V(x)=\Var_{[x,y]}f\geq\abs{f(y)-f(x)}$. Both $p(y)-p(x)$ and $q(y)-q(x)$ are therefore nonnegative. Conversely the difference of two bounded monotone functions has finite variation by the preceding entries. This is a decomposition of functions; it does not invoke signed measures or their Jordan decomposition.
\end{proof}

\KEntry{K03.15.one_sided}{One-sided limits and discontinuities of BV functions}
\begin{theorem}
A real BV function has finite left and right limits at every appropriate interior or endpoint approach and has at most countably many discontinuities.
\end{theorem}
\begin{proof}
For bounded nondecreasing $h$, $h(c-)=\sup_{x<c}h(x)$ and $h(c+)=\inf_{x>c}h(x)$ give the finite one-sided limits by the defining supremum and infimum properties. At an interior discontinuity, the interval $(h(c-),h(c+))$ has positive length, since $h(c-)\leq h(c)\leq h(c+)$. For distinct $c$ these open intervals are disjoint. Choose in each a rational number; disjointness makes this an injection into $\Q$, hence there are at most countably many discontinuities. Endpoints add at most two. Apply the two conclusions to the monotone functions in Jordan decomposition: their difference has one-sided limits, and can be discontinuous only in the union of their discontinuity sets.
\end{proof}

\KEntry{K03.15.regulated}{Regulated functions and step approximation}
\begin{theorem}
Suppose a scalar $f$ on $[a,b]$ has finite left limits at every point of $(a,b]$ and finite right limits at every point of $[a,b)$. Then for every $\eps>0$ there is a finite partition and a function $s$ constant on each open partition interval, with separately prescribed node values, such that $\sup\abs{f-s}<\eps$. In particular $f$ is bounded and Riemann integrable.
\end{theorem}
\begin{proof}
At each $c$, choose a relative neighborhood such that the oscillation of $f$ on its portion strictly to the left of $c$ is $<\eps$ and likewise on its portion strictly to the right; the appropriate one-sided limits give this. These neighborhoods cover $[a,b]$. Choose a finite subcover with centers $c_1,\ldots,c_r$ and a Lebesgue number $\delta>0$: every set of diameter $<\delta$ lies in a member. To see this elementary compactness fact, if no such number existed choose sets of diameter $<1/n$ not lying in any member and points $x_n$ in them; a convergent subsequence of $x_n$ has a limit in one of the open members, which then contains the corresponding small sets, a contradiction.

Choose a partition of mesh $<\delta$ containing all $c_i$. Each open partition interval lies in one of the selected neighborhoods and on one side of its center, since every center is a node. Its oscillation is $<\eps$. Define $s$ there by a value of $f$ in that interval, and put $s=f$ at nodes. Then $\sup\abs{f-s}<\eps$. Taking one such finite-valued $s$ proves boundedness of $f$.

A finite step function with arbitrary values at its finitely many nodes is Riemann integrable: put those nodes in small intervals of total length $<\eta$; outside these intervals its upper and lower sums agree, and inside their gap is at most $2\norm{s}_\infty\eta$. The Darboux criterion applies as $\eta\to0$. Taking step approximations with errors tending to zero and using uniform stability of the Riemann integral proves integrability of $f$.
\end{proof}

\KEntry{K03.15.bv_riemann}{BV functions are Riemann integrable}
\begin{theorem}
Every real BV function on $[a,b]$ is bounded, regulated, and Riemann integrable. The same holds for complex BV functions.
\end{theorem}
\begin{proof}
For real functions combine \KRef{K03.15.one_sided} and \KRef{K03.15.regulated}. For complex $f$, the variation of each real component is at most $\Var f$ because each component increment has magnitude at most that of the complex increment. Apply the real result to the two components. These are ordinary Riemann statements; integrability against a discontinuous Stieltjes integrator is a separate condition.
\end{proof}

\KEntry{K03.15.cantor_function}{The Cantor staircase}
\begin{construction}
There exists a continuous nondecreasing $C:[0,1]\to[0,1]$ with $C(0)=0$, $C(1)=1$, constant on every complementary interval of the middle-thirds Cantor set $D$. The set $D$ has length zero and $C'(x)=0$ for every $x\notin D$, although $C$ is not constant.
\end{construction}
\begin{proof}
Start with $C_0(x)=x$. On continuous functions $h$ with endpoint values $0,1$ define
\[(Th)(x)=\begin{cases}\tfrac12h(3x),&0\leq x\leq1/3,\\1/2,&1/3\leq x\leq2/3,\\\tfrac12+\tfrac12h(3x-2),&2/3\leq x\leq1.\end{cases}\]
The pieces agree at their joins, and $T$ preserves continuity, the endpoint values, and nondecreasingness. Put $C_n=T^nC_0$. Since $\norm{Th-Tg}_\infty\leq\frac12\norm{h-g}_\infty$, we have $\norm{C_{n+1}-C_n}_\infty\leq2^{-n}\norm{C_1-C_0}_\infty$. The geometric bound makes $C_n$ uniformly Cauchy. Its uniform limit $C$ is continuous, nondecreasing, and has the endpoint values. Passing to the limit in $C_{n+1}=TC_n$ gives $C=TC$.

At stage $m$, delete the open middle thirds from each of the $2^{m-1}$ surviving intervals. The equation $C=TC$, iterated along the left or right affine branch, shows that $C$ is constant on each interval deleted at any finite stage: the branch reaches the middle interval, where its value is fixed. The remaining Cantor set $D$ is contained in $2^m$ closed intervals of length $3^{-m}$. Their total length is $(2/3)^m\to0$. Slightly enlarge this finite collection to open intervals with any prescribed additional total length, obtaining arbitrarily short covers; equivalently one may choose such covers with total length $<\eps$ to meet the length-zero definition in \KRef{K03.08.zero_covers}. Every point outside $D$ is in an open deleted interval, so $C$ is constant in a neighborhood of that point and its derivative there is zero. The differing endpoint values show that it is not constant; its variation is $1$. This construction needs no measure differentiation theorem.
\end{proof}

\KEntry{K03.15.continuous_variation}{Continuity of the variation function}
\begin{theorem}
If a scalar $f$ is continuous and BV on $[a,b]$, then $V(x)=\Var_{[a,x]}f$ is continuous. For real $f$, its Jordan decomposition can consequently be chosen with continuous monotone components.
\end{theorem}
\begin{proof}
Fix $c<b$ and $\eta>0$. Choose a partition $c=t_0<t_1<\cdots<t_m=b$ whose variation sum exceeds $\Var_{[c,b]}f-\eta$. For $c<x<t_1$, using the partition $x,t_1,\ldots,t_m$ on $[x,b]$ gives
\[\begin{split}\Var_{[c,x]}f&=\Var_{[c,b]}f-\Var_{[x,b]}f\\
&\leq\eta+\abs{f(t_1)-f(c)}-\abs{f(t_1)-f(x)}\leq\eta+\abs{f(x)-f(c)}.
\end{split}\]
Continuity makes the limsup as $x\downarrow c$ at most $\eta$, hence zero. The same argument on $[a,c]$, deleting its last partition interval, proves left continuity. It also handles the endpoints in their appropriate directions. Additivity of variation translates these bounds into continuity of $V$. The Jordan formulas now give continuous components. Merely right continuity of $f$ only yields the corresponding right continuity here; it cannot imply full continuity of $V$, as a right-continuous step with a left jump demonstrates.
\end{proof}

\KEntry{K03.15.c1_variation}{Variation of a continuously differentiable scalar path}
\begin{theorem}
If $f:[a,b]\to\R$ or $\C$ is $C^1$ (with continuous one-sided extensions of $f'$ at the endpoints), then
\[\Var_{[a,b]}f=\int_a^b\abs{f'(t)}\dd t.\]
\end{theorem}
\begin{proof}
The fundamental theorem of calculus, applied componentwise, gives $f(v)-f(u)=\int_u^v f'$. Summing its norm bound over a partition proves $\Var f\leq\int\abs{f'}$. Conversely let $\omega(\delta)=\sup_{\abs{s-t}\leq\delta}\abs{f'(s)-f'(t)}\to0$. For each interval of a tagged partition,
\[\abs{f(x_i)-f(x_{i-1})-f'(t_i)\Delta x_i}\leq\omega(\norm P)\Delta x_i.\]
Thus its polygonal sum is at least $\sum_i\abs{f'(t_i)}\Delta x_i-\omega(\norm P)(b-a)$. Let the mesh tend to zero; the Riemann sum tends to $\int\abs{f'}$ and the error to zero. This proves the opposite inequality, without an invalid complex scalar mean value equality.
\end{proof}

\KEntry{K03.15.rectifiable}{Rectifiable continuous paths}
\begin{definition}
A continuous scalar path $\gamma:[a,b]\to\K$ has length equal to the supremum of the lengths of its inscribed polygonal paths, namely $\sup_P\sum_i\abs{\gamma(x_i)-\gamma(x_{i-1})}$. It is rectifiable if this is finite. By the identical partition formulas its length is precisely $\Var\gamma$; for a $C^1$ path it is the preceding integral. A discontinuous BV function is not a continuous path under this definition.
\end{definition}

\KEntry{K03.15.rs_sums}{The refinement Stieltjes definition}
\begin{definition}
Let $f:[a,b]\to\K$ and $\alpha:[a,b]\to\R$ be bounded. For $P=(x_0,\ldots,x_m)$ and arbitrary tags $t_i\in[x_{i-1},x_i]$, set
\[S(P,t;f,\alpha)=\sum_{i=1}^mf(t_i)(\alpha(x_i)-\alpha(x_{i-1})).\]
The refinement integral $\int_a^b f\dd\alpha=I$ exists when for every $\eps>0$ there is a fixed partition $P_\eps$ such that for every refinement $P\supset P_\eps$ and every choice of tags, $\abs{S(P,t;f,\alpha)-I}<\eps$. Set the degenerate integral to zero and reverse orientation by a minus sign. Endpoint values are retained. This is Apostol's Definition 7.1, not an unrestricted small-mesh definition.
\end{definition}
\begin{lemma}
This integral is unique. It exists if and only if for each $\eps>0$ there is $P_\eps$ such that any two tagged sums whose partitions refine $P_\eps$ differ by $<\eps$.
\end{lemma}
\begin{proof}
Two proposed limits are compared using a tagged sum on the union of their controlling partitions; their distance is $<2\eps$. Existence implies the Cauchy condition by the triangle inequality. Conversely choose controlling partitions for errors $2^{-n}$, replace them by their successive unions to make them nested, and choose one tagged sum $S_n$ on each. They are scalar Cauchy and have a limit $I$. Every tagged sum on a refinement of the $n$th partition differs from $S_m$ by $<2^{-n}$ for $m\geq n$; let $m\to\infty$ to bound its distance to $I$ by $2^{-n}$. This is the required refinement limit.
\end{proof}

\KEntry{K03.15.rs_darboux}{Monotone integrators: the Darboux bridge}
\begin{theorem}
For bounded real $f$ and nondecreasing $\alpha$, put
\[U_\alpha(f,P)=\sum_i\sup_{[x_{i-1},x_i]}f\,\Delta\alpha_i,\qquad
L_\alpha(f,P)=\sum_i\inf_{[x_{i-1},x_i]}f\,\Delta\alpha_i.\]
The refinement integral exists if and only if $\inf_P U_\alpha(f,P)=\sup_P L_\alpha(f,P)$, equivalently for every $\eps>0$ some $P$ satisfies $U_\alpha-L_\alpha<\eps$. Its value is the common number. No continuity of $\alpha$ is required.
\end{theorem}
\begin{proof}
Refinement decreases upper sums and increases lower sums because the weights are nonnegative and additive. A common refinement therefore shows every lower sum is at most every upper sum. Their infimum and supremum are finite. If they agree, choose partitions with upper and lower sums within $\eps/2$ of the common value and take their union. All tags on all refinements are squeezed between those bounds.

Conversely, suppose the refinement limit is $I$. On its controlling partition every tagged sum lies within $\eta$ of $I$. The supremum and infimum over tags on that fixed partition are exactly $U_\alpha$ and $L_\alpha$: approximate each interval supremum and infimum separately; finiteness of the partition and nonnegative weights allow arbitrary accuracy in their weighted sums (zero weights need no approximation). Thus $U_\alpha-L_\alpha\leq2\eta$. Since $\eta$ is arbitrary, the epsilon criterion follows, and that criterion forces equality of the upper and lower integrals by their ordering. This proves the precise equivalence with Rudin's monotone Darboux formulation.
\end{proof}

\KEntry{K03.15.rs_existence}{Continuous integrands and BV integrators}
\begin{theorem}
If $f:[a,b]\to\K$ is continuous and real $\alpha$ is BV, the refinement integral exists and
\[\abs{\int_a^b f\dd\alpha}\leq\norm f_\infty\Var\alpha.\]
\end{theorem}
\begin{proof}
For a partition $P$ of mesh $<\delta$ and a refinement $Q$, distribute each increment $\Delta\alpha$ of $P$ into its subinterval increments in $Q$. The original tag and each new tag lie in the same original interval, so
\[\abs{S(P,t)-S(Q,u)}\leq\omega_f(\delta)\sum_{J\in Q}\abs{\Delta\alpha_J}\leq\omega_f(\delta)\Var\alpha,\]
where uniform continuity gives $\omega_f(\delta)\to0$. For two refinements of one sufficiently fine fixed partition, compare them on their common refinement; their tagged sums differ by at most $2\omega_f(\delta)\Var\alpha$. The refinement Cauchy criterion proves existence. Every sum has norm at most $\norm f_\infty\Var\alpha$, giving the integral bound by passage to the limit. The same common-refinement estimate actually proves mesh convergence for continuous integrands even when $\alpha$ jumps; this special sufficient condition does not identify the two definitions for arbitrary bounded integrands.
\end{proof}

\KEntry{K03.15.rs_algebra}{Linearity and interval additivity}
\begin{theorem}
Where the indicated integrals exist, the refinement integral is scalar linear in $f$ and real linear in its real integrator. For $a<c<b$, integrability on $[a,b]$ is equivalent to integrability on both $[a,c]$ and $[c,b]$, and
\[\int_a^b f\dd\alpha=\int_a^c f\dd\alpha+\int_c^b f\dd\alpha.\]
\end{theorem}
\begin{proof}
The finite tagged sums are linear in both slots separately. Take the union of controlling partitions for finitely many integrals and apply the triangle inequality, with the scalar coefficients multiplying their error bounds, to pass these identities to their limits.

If both subinterval integrals exist, take the union of their controlling partitions, including $c$. Every refinement splits into the two subinterval sums, and the errors add; this proves whole-interval integrability and its value. Conversely take a whole-interval Cauchy-controlling partition and insert $c$. Two arbitrary tagged sums on refinements of its left restriction can be extended using identical right partitions and tags. Their difference is the difference of the whole sums, hence is small. The Cauchy criterion gives the left integral; the identical argument gives the right one. The already proved implication identifies their sum with the whole integral. Shared endpoint tags are permitted throughout and no endpoint values have been removed.
\end{proof}

\KEntry{K03.15.rs_ordinary}{Reduction to ordinary Riemann integration}
\begin{theorem}
For $\alpha(x)=x$ and bounded scalar $f$, the refinement Stieltjes integral exists exactly when the ordinary Riemann integral exists, and their values agree.
\end{theorem}
\begin{proof}
For real $f$, the weighted upper and lower sums in \KRef{K03.15.rs_darboux} are literally the ordinary Darboux sums because $\Delta\alpha_i=\Delta x_i$. Thus their equality is the ordinary integrability criterion with the same value. For complex $f$, existence of either tagged limit is equivalent to existence of its two component limits; apply the real conclusion. The ordinary Darboux/mesh comparison in \KRef{K03.08.tagged_equivalence} then supplies the ordinary mesh interpretation.
\end{proof}

\KEntry{K03.15.rs_density}{A smooth integrator}
\begin{theorem}
For continuous $f$ and $\alpha\in C^1([a,b],\R)$,
\[\int_a^b f\dd\alpha=\int_a^bf(x)\alpha'(x)\dd x.\]
\end{theorem}
\begin{proof}
The left integral exists since $\alpha$ is BV, and the right integrand is continuous. The fundamental theorem gives $\Delta\alpha_i=\int_{x_{i-1}}^{x_i}\alpha'$. Hence the difference between a Stieltjes sum and the ordinary tagged sum for $f\alpha'$ is bounded by
\[\norm f_\infty\,\omega_{\alpha'}(\norm P)(b-a)\longrightarrow0.\]
The latter sums tend to the ordinary integral. Choose partitions of arbitrarily small mesh also refining any controlling Stieltjes partition; uniqueness identifies the two integrals. The estimate in fact gives a mesh limit here directly.
\end{proof}

\KEntry{K03.15.rs_parts}{Stieltjes integration by parts}
\begin{theorem}
If real $f$ and $\alpha$ are continuous BV functions on $[a,b]$, then
\[\int_a^bf\dd\alpha+\int_a^b\alpha\dd f=f(b)\alpha(b)-f(a)\alpha(a).\]
\end{theorem}
\begin{proof}
Both integrals exist by \KRef{K03.15.rs_existence}. For every partition, finite multiplication and telescoping give
\[\sum_i f(x_{i-1})\Delta\alpha_i+\sum_i\alpha(x_i)\Delta f_i=f(b)\alpha(b)-f(a)\alpha(a).\]
The first and second sums are permitted tagged sums for the two integrals, respectively. Choose a common refinement of their controlling partitions for any given accuracy and use these left and right tags. Each sum differs from its integral by an arbitrarily small amount, so the finite identity proves the integral identity. No jump correction formula is being assumed.
\end{proof}

\KEntry{K03.15.rs_jump}{Step integrators and endpoint atoms}
\begin{theorem}
Let $\alpha$ be a real step function with finitely many nodes $c\in(a,b)$, allowing arbitrary values at those nodes. Let bounded $f$ be continuous at every node where either $\alpha(c)-\alpha(c-)$ or $\alpha(c+)-\alpha(c)$ is nonzero, and at any endpoint carrying a jump. Then
\[\int_a^b f\dd\alpha=f(a)\bigl(\alpha(a+)-\alpha(a)\bigr)
+\sum_{a<c<b}f(c)\bigl(\alpha(c+)-\alpha(c-)\bigr)
+f(b)\bigl(\alpha(b)-\alpha(b-)\bigr).\]
Zero net weights may be omitted from the displayed numerical sum, but do not remove the continuity requirement if the separate one-sided increments are nonzero. Only a node where both increments vanish carries no such requirement.
\end{theorem}
\begin{proof}
Decompose $\alpha$, up to a constant, into finitely many multiples of $H_c^-(x)=\mathbf1_{[c,b]}(x)$ and $H_c^+(x)=\mathbf1_{(c,b]}(x)$ at interior $c$, with coefficients $\alpha(c)-\alpha(c-)$ and $\alpha(c+)-\alpha(c)$ respectively; use $H_a^+$ for the right endpoint-value jump at $a$ and $H_b^-$ for the left jump at $b$. This decomposition agrees at every interval and every individual node.

For $H_c^-$ the only nonzero increment on a partition containing $c$ is the interval immediately to its left; for $H_c^+$ it is the interval immediately to its right. Insert nodes on the appropriate side sufficiently close to $c$. All tags in the resulting nonzero interval have $f$-value within $\eps$ of $f(c)$ by continuity; the same remains true for every refinement. Each elementary integral is therefore $f(c)$. The same argument gives $f(a)$ and $f(b)$ for the endpoint terms. Linearity proves the formula; the two interior coefficients add to the full jump. In particular a right-continuous step $\mathbf1_{[c,b]}$ gives an atom at an interior $c$ or at $b$, while its value already present at $a$ contributes no atom at $a$. These conventions will determine the outgoing $(a,b]$ measure interface.
\end{proof}

\begin{remark}
A zero net jump can still conceal an obstruction. Set $\alpha(c)=f(c)=1$ at one interior node $c$, and set both functions equal to zero elsewhere. Both one-sided limits of $\alpha$ are zero, but the adjacent increments on a partition containing $c$ are $+1$ and $-1$. Independently choosing a tag at $c$ or away from it in the two adjacent intervals produces sums $-1,0,1$ on every such refinement. Every proposed controlling partition can be refined to include $c$, so there is no refinement integral. This proves why cancellation of the displayed net weight does not justify dropping continuity at that node.
\end{remark}

\KEntry{K03.15.rs_common_jumps}{Same-sided jumps can obstruct existence}
\begin{proposition}
Let $\alpha$ be nondecreasing. If $\alpha(c+)-\alpha(c)>0$ and $f$ is not right-continuous at $c$ in the sense that $\limsup_{x\downarrow c}\abs{f(x)-f(c)}>0$, for bounded real $f$, then the refinement integral does not exist. The analogous assertion holds for positive left jump $\alpha(c)-\alpha(c-)$ and failure of left continuity of $f$. In particular common nonzero jumps on the same side obstruct existence; opposite-sided discontinuities need not.
\end{proposition}
\begin{proof}
Choose $\eta>0$ such that every sufficiently small right neighborhood contains a point with $\abs{f(x)-f(c)}>\eta$. In any partition containing $c$, its first interval $[c,u]$ to the right has oscillation at least $\eta$ and weight $\alpha(u)-\alpha(c)\geq\alpha(c+)-\alpha(c)=J>0$. Thus $U_\alpha-L_\alpha\geq\eta J$. Every partition can be refined to include $c$; refinement decreases this gap, so no partition could have had a gap below $\eta J$. The Darboux criterion rules out integrability. Use the interval to the left for the other assertion. This argument uses one-sided oscillation, not merely a label of ``common discontinuity.''
\end{proof}

\KEntry{K03.15.rs_mesh_bridge}{The mesh bridge for continuous monotone integrators}
\begin{theorem}
For bounded scalar $f$ and continuous nondecreasing $\alpha$, refinement integrability is equivalent to existence of the unrestricted tagged mesh limit: there is $I$ such that every tagged sum with sufficiently small $\norm P$ lies near $I$. Their values agree.
\end{theorem}
\begin{proof}
A mesh limit implies a refinement limit by choosing a controlling partition whose mesh is below the mesh threshold. For the converse first take real $f$, $\abs f\leq M$. Choose a fixed partition $Q$ with $U_\alpha(f,Q)-L_\alpha(f,Q)<\eps$ by the Darboux criterion. For any small-mesh $P$, let $T=P\cup Q$. Intervals of $P$ not containing an interior node of $Q$ require no splitting; choose the same tags there in $T$. On an interval that is split, the difference of the original contribution and all new contributions is at most $2M\Delta\alpha$.

There are at most $r$ such intervals, where $r$ is the number of interior nodes of $Q$, and every one has length at most $\norm P$. Uniform continuity of $\alpha$ makes their total $\alpha$-weight at most $r\omega_\alpha(\norm P)$. Thus
\[\abs{S(P,t)-S(T,u)}\leq2Mr\omega_\alpha(\norm P)\longrightarrow0.\]
Every sum on $T$ lies between $L_\alpha(f,Q)$ and $U_\alpha(f,Q)$, as does the refinement integral, so its distance to that integral is $<\eps$. The displayed bound proves mesh convergence. If $r=0$ or $M=0$ the bound is zero. For complex $f$ apply the proof to real and imaginary components and combine the errors. Continuity of $\alpha$ is precisely what controls the crossing intervals; the next example shows why it cannot be deleted from this general bridge.
\end{proof}

\KEntry{K03.15.rs_mesh_counterexample}{Opposite-sided jumps separate the definitions}
\begin{example}
On $[-1,1]$ set $\alpha(x)=0$ for $x\leq0$, $\alpha(x)=1$ for $x>0$, and $f(x)=0$ for $x<0$, $f(x)=1$ for $x\geq0$. Every refinement of the partition $(-1,0,1)$ has exactly one nonzero $\alpha$-increment, namely $1$ on the first interval to the right of zero. All its tags have $f=1$. Hence every such sum is $1$, and the refinement integral is $1$.

For any $\delta>0$, take a partition of mesh $<\delta$ with an interval $[-h,h]$, $0<h<\delta/2$, crossing zero and with no node at zero; fill both remaining sides with intervals of length $<\delta$. Its only nonzero increment is the crossing interval. A negative tag gives sum $0$, while tag zero or any positive tag gives sum $1$. These choices persist at arbitrarily small mesh. The mesh limit therefore does not exist. This is an explicit opposite-sided-jump witness, consistent with both Apostol's refinement definition and Rudin's Darboux criterion.
\end{example}

\KEntry{K03.15.rs_uniform}{Uniform stability for a fixed BV integrator}
\begin{theorem}
For fixed real BV $\alpha$, if bounded $f_n$ are refinement integrable and $f_n\to f$ uniformly, then $f$ is refinement integrable and
\[\abs{\int f\dd\alpha-\int f_n\dd\alpha}\leq\norm{f-f_n}_\infty\Var\alpha\longrightarrow0.\]
\end{theorem}
\begin{proof}
Every tagged sum for a bounded $h$ has magnitude at most $\norm h_\infty\Var\alpha$. If $h$ is integrable, pass to the limit to obtain the same integral bound. By linearity the integral values $I_n$ therefore satisfy $\abs{I_n-I_m}\leq\norm{f_n-f_m}_\infty\Var\alpha$; they have a scalar limit $I$. Given $\eps>0$, choose $n$ making both $\norm{f-f_n}_\infty\Var\alpha$ and $\abs{I_n-I}$ less than $\eps/3$, then a controlling partition for $f_n$ with error $<\eps/3$. On every refinement,
\[\abs{S(f)-I}\leq\abs{S(f)-S(f_n)}+\abs{S(f_n)-I_n}+\abs{I_n-I}<\eps.\]
Thus $f$ is integrable with value $I$. Apply the integral bound to $f-f_n$ for the stated estimate. If $\Var\alpha=0$, all sums are zero and the assertions are immediate.
\end{proof}

\KEntry{K03.15.stieltjes_measure_bridge}{Outgoing interface to Lebesgue--Stieltjes measures}
\begin{remark}
This is a routing record to D11, not a theorem proved or assumed in D03. The intended future comparison starts from a right-continuous nondecreasing $\alpha$, constructs a measure satisfying $\mu_\alpha((u,v])=\alpha(v)-\alpha(u)$, and for continuous $f$ compares the refinement integral on $[a,b]$ with the measure integral over $(a,b]$. Right continuity removes the right endpoint-value jump at $a$, and the half-open interval includes the atom at $b$, matching \KRef{K03.15.rs_jump}. Measure construction, measurable integrands, signed measures, and proof of this comparison remain in D11. No measure assertion is used in any proof in this chapter.
\end{remark}

% BEGIN REVISED SOURCE chapters/M03.16.tex SHA256=1156e6fc11301242fbb9cf04f2b7e335a009aeabe6983c26fc551bcb0e2f8663
\chapter{Compactness and approximation in scalar continuous-function spaces}
\SourceRef{Rudin, Chapter 7, including Theorem 7.18 (PDF 165, formula checked against the scan); Pugh, Sections 4.3--4.4; Zorich II, Section 16.4. The domain here is a compact subset of the real line and the target is scalar.}

\KEntry{K03.16.sup_space}{The scalar space \texorpdfstring{$C(K)$}{C(K)}}
\begin{definition}
Let $K\subset\R$ be nonempty and compact. The real or complex vector space $C(K,\K)$ consists of continuous scalar functions with pointwise operations and norm
\[\norm f_\infty=\max_{x\in K}\abs{f(x)}.\]
The maximum exists by compactness and continuity. The triangle inequality and homogeneity follow pointwise and then by taking maxima, and zero norm implies that every value is zero. The distance is $\norm{f-g}_\infty$, whose convergence is uniform convergence. For $K=\varnothing$, the space has only its zero (empty) function, with norm zero; all subsequent approximation and compactness assertions are then trivial.
\end{definition}

\KEntry{K03.16.complete}{Completeness of \texorpdfstring{$C(K)$}{C(K)}}
\begin{theorem}
Every sup-norm Cauchy sequence in $C(K,\K)$ converges in that norm to a member of $C(K,\K)$.
\end{theorem}
\begin{proof}
Norm Cauchy convergence is exactly uniform Cauchy convergence. By \KRef{K03.12.cauchy}, scalar completeness gives a uniform scalar limit $f$. By \KRef{K03.12.continuity}, $f$ is continuous. It is bounded on compact $K$, hence belongs to $C(K)$, and the supremum criterion gives convergence in the norm.
\end{proof}

\KEntry{K03.16.equicontinuity}{Equicontinuity and boundedness}
\begin{definition}
A family $\mathcal F\subset C(K,\K)$ is pointwise equicontinuous if for every $x\in K$ and $\eps>0$ there is $\delta_x>0$ such that $\abs{f(y)-f(x)}<\eps$ for every $f\in\mathcal F$ whenever $y\in K$, $\abs{x-y}<\delta_x$. It is uniformly equicontinuous if a single $\delta$ works for every pair $x,y\in K$. Uniform boundedness means $\sup_{f\in\mathcal F,x\in K}\abs{f(x)}<\infty$.
\end{definition}
\begin{theorem}
On compact $K$, pointwise equicontinuity implies uniform equicontinuity. For such a family, boundedness at each individual point implies uniform boundedness.
\end{theorem}
\begin{proof}
For each $x$ choose $r_x>0$ such that $\abs{f(y)-f(x)}<\eps/2$ whenever $\abs{y-x}<2r_x$, for every $f$. Finitely many balls $B(x_i,r_{x_i})$ cover $K$. Put $\delta=\min_i r_{x_i}$. If $\abs{y-z}<\delta$ and $y\in B(x_i,r_{x_i})$, then both $y,z$ lie in $B(x_i,2r_{x_i})$; the triangle inequality gives $\abs{f(y)-f(z)}<\eps$ for every $f$. For boundedness choose a finite net with this common modulus at error $1$. At each of its finitely many centers the family values are bounded; the largest center bound plus $1$ bounds every value.
\end{proof}

\KEntry{K03.16.ascoli}{Arzela--Ascoli for scalar functions on compact real sets}
\begin{theorem}
A uniformly bounded equicontinuous sequence in $C(K,\K)$ has a uniformly convergent subsequence. A family $\mathcal F\subset C(K,\K)$ is compact in sup norm if and only if it is closed, uniformly bounded, and equicontinuous.
\end{theorem}
\begin{proof}
Take finite $1/j$-nets of $K$ and let $D$ be their countable union, dense in $K$. At the first point of $D$ use bounded scalar subsequence compactness (\KRef{K03.01.bolzano}); at the second point extract a further subsequence, and continue. The diagonal subsequence $g_n$ converges at every point of $D$. For $\eps>0$, uniform equicontinuity supplies $\delta>0$ giving error $\eps/3$ between points less than $\delta$ apart, simultaneously for all $g_n$. Choose a finite $\delta$-net with centers in $D$. Their finitely many value sequences are Cauchy, so for large $n,m$, all center differences are $<\eps/3$. For any $x\in K$, compare $g_n(x)$ and $g_m(x)$ through a nearby center: the three errors sum to $<\eps$. Thus $g_n$ is uniformly Cauchy and converges in $C(K)$ by completeness.

For sufficiency in the family statement this argument gives a subsequence converging in $C(K)$ from every sequence in $\mathcal F$; closedness puts its limit in $\mathcal F$. For clarity, sequential compactness implies compactness here as follows. A sequentially compact metric space has a finite $\eta$-net for each $\eta>0$, since otherwise inductively choosing points at mutual distances at least $\eta$ produces a sequence with no convergent subsequence. Given an open cover, it has a Lebesgue number: otherwise choose points $x_n$ such that their relative $1/n$-balls lie in no single cover member; a convergent subsequence $x_{n_j}\to x$ and a cover member containing a ball about $x$ contradict this choice for large $j$. A finite net with radius less than half the Lebesgue number now yields a finite subcover. Hence $\mathcal F$ is compact.

Conversely a compact family has a finite sup-norm $\eta$-net for each $\eta$. It is bounded, since the norms of a finite $1$-net plus $1$ bound all its norms. For equicontinuity use a finite $\eps/3$-net of continuous functions; each is uniformly continuous on $K$, and the minimum of their continuity radii controls the middle term in
\[\abs{f(x)-f(y)}\leq\abs{f(x)-h(x)}+\abs{h(x)-h(y)}+\abs{h(y)-f(y)}<\eps.\]
Finally a compact subset of a metric space is closed: for a point $g$ outside it, the balls $B(f,\norm{g-f}/2)$ cover the compact family; a finite subcover yields a positive lower bound for its distances to $g$, so a ball about $g$ misses it. This completes both directions.
\end{proof}

\KEntry{K03.16.ascoli_hypotheses}{Why neither bound nor equicontinuity alone suffices}
\begin{example}
The constant functions $f_n(x)=n$ on any nonempty $K$ are equicontinuous, but their values at every fixed point are unbounded; no subsequence converges uniformly to a finite-valued function. On $[0,2\pi]$, $g_n(x)=\sin(nx)$ has common bound $1$, but $g_n(0)=0$ and $g_n(\pi/(2n))=1$, with $\pi/(2n)\to0$. Thus it is not equicontinuous at zero. No subsequence $g_{n_j}$ can converge uniformly: its continuous limit $g$ would satisfy $g(0)=0$ and $g(\pi/(2n_j))\to0$, whereas uniform convergence would also make those latter values tend to $1$.
\end{example}
\begin{remark}
Pugh PDF 238 lists ODE and holomorphic-family applications without displaying all necessary hypotheses; they are not imported here. For example, $y_n=n$ solves $y'=0$ but has no finite-valued convergent subsequence. Zorich II PDF 412 includes a parenthesis suggesting local compactness of a target suffices to equate boundedness with total boundedness. For a general metric this fails: $\N$ with distance $1$ between distinct points is locally compact (singletons are compact neighborhoods) and bounded, yet no finite collection of balls of radius $1/2$ covers it. The scalar and finite-dimensional target cases used above are valid; no broader target claim is needed.
\end{remark}

\KEntry{K03.16.piecewise_linear}{Piecewise-linear uniform approximation}
\begin{theorem}
For real continuous $f$ on $[a,b]$ and a partition $P$, let $L_P$ interpolate $f$ linearly between successive nodes. If $\omega_f(\delta)=\sup_{\abs{x-y}\leq\delta}\abs{f(x)-f(y)}$, then
\[\norm{f-L_P}_\infty\leq\omega_f(\norm P)\longrightarrow0\quad\text{as }\norm P\to0.\]
\end{theorem}
\begin{proof}
For $x\in[u,v]$ write $x=(1-t)u+tv$, $0\leq t\leq1$. Then $L_P(x)=(1-t)f(u)+tf(v)$, so
\[\abs{f(x)-L_P(x)}\leq(1-t)\abs{f(x)-f(u)}+t\abs{f(x)-f(v)}\leq\omega_f(v-u).\]
Uniform continuity gives $\omega_f(\delta)\to0$. The degenerate interval is handled by its constant interpolant.
\end{proof}

\KEntry{K03.16.bernstein}{Bernstein polynomials and finite moment identities}
\begin{construction}
For $f\in C([0,1],\K)$ and $n\geq1$ define
\[B_nf(x)=\sum_{k=0}^nf(k/n)w_{n,k}(x),\qquad w_{n,k}(x)=\binom nkx^k(1-x)^{n-k}.\]
The weights are nonnegative and satisfy
\[\sum_kw_{n,k}=1,\quad\sum_k(k/n)w_{n,k}=x,\quad\sum_k(k/n-x)^2w_{n,k}=x(1-x)/n.\]
\end{construction}
\begin{proof}
The first identity is the binomial theorem for $(x+1-x)^n$. Using $k\binom nk=n\binom{n-1}{k-1}$ and shifting the finite index gives $\sum k w_{n,k}=nx$. Likewise $k(k-1)\binom nk=n(n-1)\binom{n-2}{k-2}$ gives $\sum k(k-1)w_{n,k}=n(n-1)x^2$; for $n=1$ this second factorial moment is zero directly. Add the first moment to get the second moment and expand the square:
\[\sum(k/n-x)^2w_{n,k}=\frac{n(n-1)x^2+nx}{n^2}-2x^2+x^2=\frac{x(1-x)}n.\]
All identities hold at $0,1$ as polynomial identities, with zeroth powers interpreted as $1$.
\end{proof}

\KEntry{K03.16.bernstein_convergence}{Bernstein's uniform estimate}
\begin{theorem}
For $M=\norm f_\infty$ and any $\delta>0$,
\[\norm{B_nf-f}_\infty\leq\omega_f(\delta)+\frac{M}{2n\delta^2}.\]
In particular $B_nf\to f$ uniformly on $[0,1]$.
\end{theorem}
\begin{proof}
At fixed $x$, write the error as $\sum_k(f(k/n)-f(x))w_{n,k}(x)$. On the indices $\abs{k/n-x}\leq\delta$ its magnitude is bounded by $\omega_f(\delta)\sum_kw_{n,k}\leq\omega_f(\delta)$. On the other indices it is bounded by $2M$ times their total weight, which is at most
\[\delta^{-2}\sum_k(k/n-x)^2w_{n,k}=x(1-x)/(n\delta^2)\leq1/(4n\delta^2).\]
This gives the estimate uniformly in $x$. First choose $\delta$ making $\omega_f(\delta)$ small, then choose $n$ to make the second term small. Alternatively $\delta=n^{-1/4}$ gives an explicit sequence tending to zero. This proof uses only finite algebra and uniform continuity, not a probability convergence theorem.
\end{proof}

\KEntry{K03.16.weierstrass}{Weierstrass approximation on an interval}
\begin{theorem}
For $f\in C([a,b],\R)$ and every $\eps>0$, a real polynomial $p$ satisfies $\norm{f-p}_\infty<\eps$. The same assertion holds with complex coefficients for complex continuous $f$.
\end{theorem}
\begin{proof}
For $a<b$ apply the Bernstein construction to $g(t)=f(a+(b-a)t)$, and substitute $t=(x-a)/(b-a)$ into its approximating polynomials. This affine substitution preserves polynomials and the supremum error, proving the real result. For complex $f=u+iv$, approximate $u,v$ by real polynomials with errors $<\eps/2$ and combine them; the triangle inequality gives total error $<\eps$. If $a=b$, use the constant polynomial $f(a)$.
\end{proof}

\KEntry{K03.16.lattice}{Algebra closures and the absolute-value construction}
\begin{theorem}
The uniform closure $\overline A$ of a real unital subalgebra $A\subset C(K,\R)$ is a unital algebra and is closed under absolute values, finite maxima, and finite minima.
\end{theorem}
\begin{proof}
Addition and scalar multiplication pass to uniform limits. For multiplication, if $f_n\to f$ and $g_n\to g$ uniformly, then $f_n$ are uniformly bounded and
\[\norm{f_ng_n-fg}_\infty\leq\norm{f_n}_\infty\norm{g_n-g}_\infty+\norm{f_n-f}_\infty\norm g_\infty\longrightarrow0.\]
Thus $\overline A$ is an algebra. For $f\in\overline A$ choose $M\geq\norm f_\infty$. By the explicit Bernstein/affine construction in \KRef{K03.16.weierstrass} there are polynomials $p_j$ with $\sup_{[-M,M]}\abs{p_j(t)-\abs t}\to0$; if $M=0$ nothing is needed. Since constants and $f$ belong to $\overline A$, so does $p_j(f)$, and their uniform limit is $\abs f$. Finally
\[\max(f,g)=\tfrac12(f+g+\abs{f-g}),\qquad\min(f,g)=\tfrac12(f+g-\abs{f-g}),\]
proving finite lattice closure by iteration. More generally these formulas give lattice closure of any real algebra closure already known to contain absolute values.
\end{proof}

\KEntry{K03.16.stone_real}{Real Stone--Weierstrass on compact real sets}
\begin{theorem}
If a real subalgebra $A\subset C(K,\R)$ contains constants and separates points (for $x\ne y$ some $u\in A$ has $u(x)\ne u(y)$), then its uniform closure is $C(K,\R)$.
\end{theorem}
\begin{proof}
Fix $f\in C(K,\R)$ and $\eps>0$. For $x\ne y$, choose a separating $u$ and an affine expression $h_{xy}=\lambda u+\mu$ matching $f$ at $x$ and $y$; explicitly $\lambda=(f(x)-f(y))/(u(x)-u(y))$, $\mu=f(x)-\lambda u(x)$. For $y=x$ use the constant $f(x)$. Fixing $x$, the open sets $\{t:h_{xy}(t)>f(t)-\eps\}$, indexed by $y$, cover $K$ since each contains $y$. Choose finitely many and put $g_x=\max_i h_{xy_i}\in\overline A$. Then $g_x>f-\eps$ everywhere and $g_x(x)=f(x)$. The open sets $\{t:g_x(t)<f(t)+\eps\}$ cover $K$ as $x$ varies. Choose finitely many $x_j$ and set $h=\min_jg_{x_j}\in\overline A$. Every value of $h$ is $>f-\eps$, and at each point one selected $g_{x_j}$ is $<f+\eps$, so $\norm{h-f}_\infty\leq\eps$. Since $h\in\overline A$, approximate it within any further $\eps$ by an element of $A$. As $\eps$ is arbitrary this proves density.

A frequently used source variant replaces constants by ``no common zero.'' Its reduction needs proof: for each point choose $u\in A$ nonzero there; finitely many such neighborhoods cover $K$, so $t=\sum_{i=1}^ru_i^2\in A$ has range in $[m,M]$ for some $m>0$. Approximate $1/s$ on $[m,M]$ by polynomials $p_j(s)$ using Weierstrass. Then $t p_j(t)\in A$ (the polynomial $s p_j(s)$ has zero constant term) and $t p_j(t)\to1$ uniformly. Thus constants belong to $\overline A$, and the theorem just proved applied to this closed separating algebra gives $\overline A=C(K,\R)$. This establishes the variant rather than identifying its hypotheses by name.
\end{proof}

\KEntry{K03.16.stone_complex}{Complex Stone--Weierstrass with conjugation closure}
\begin{theorem}
If a complex subalgebra $A\subset C(K,\C)$ contains constants, separates points, and is closed under pointwise complex conjugation, then $\overline A=C(K,\C)$.
\end{theorem}
\begin{proof}
Its real-valued elements form a real unital algebra $B$. For $u\in A$, the functions $(u+\overline u)/2$ and $(u-\overline u)/(2i)$ lie in $B$. Whenever $u$ distinguishes two points, at least one of these two components also distinguishes them; hence $B$ separates points. The real theorem uniformly approximates the real and imaginary parts of any continuous complex function by members of $B$. Combining those two approximations gives a member of $A$ with arbitrarily small uniform error. Conjugation closure is exactly what made the component functions available and remains an explicit hypothesis.
\end{proof}

\KEntry{K03.16.smooth_density}{Smooth uniform approximation}
\begin{theorem}
Every continuous scalar function on a compact real interval is a uniform limit of smooth functions defined on a neighborhood of that interval.
\end{theorem}
\begin{proof}
Choose approximating polynomials from \KRef{K03.16.weierstrass} with errors $<1/n$. Polynomials have derivatives of all orders everywhere, so they supply the required smooth approximants. A discontinuous function on the interval cannot be a uniform limit of continuous, hence cannot be a uniform limit of smooth functions, by \KRef{K03.12.continuity}; no assertion about such functions follows from this density theorem.
\end{proof}

\KEntry{K03.16.approximate_identity}{Elementary convolution approximation}
\begin{theorem}
Let $f:\R\to\K$ be bounded and uniformly continuous. Let $k_j$ be nonnegative continuous functions with support contained in $[-\delta_j,\delta_j]$, $\delta_j\downarrow0$, and ordinary integral $1$. Then the proper Riemann convolutions
\[T_jf(x)=\int_{-\delta_j}^{\delta_j}f(x-t)k_j(t)\dd t\]
converge uniformly on $\R$ to $f$, with $\sup_x\abs{T_jf(x)-f(x)}\leq\omega_f(\delta_j)$.
\end{theorem}
\begin{proof}
Uniform continuity implies continuity of $f$, so each integrand is continuous on the compact integration interval and the integral exists. Since the kernel has integral one,
\[\abs{T_jf(x)-f(x)}\leq\int\abs{f(x-t)-f(x)}k_j(t)\dd t\leq\omega_f(\delta_j)\int k_j=\omega_f(\delta_j)\to0.\]
The bound is independent of $x$. Such kernels exist explicitly: $k_\delta(t)=\delta^{-1}\max(1-\abs t/\delta,0)$ has triangular area $1$, is continuous, nonnegative, and has the required support. General measurable or $L^p$ convolution theory is outside this assertion.
\end{proof}

\KEntry{K03.16.nowhere_differentiable}{A complete periodic-tent nowhere differentiable example}
\begin{construction}
Let $\phi$ be the $2$-periodic extension of $\abs x$ on $[-1,1]$. Then
\[F(x)=\sum_{n=0}^\infty(3/4)^n\phi(4^nx)\]
is continuous on $\R$ and has no finite derivative at any real point.
\end{construction}
\begin{proof}
The tent has $0\leq\phi\leq1$ and is affine with slope $1$ or $-1$ between consecutive integers. Subdivide any segment at the finitely many intervening integers and add the slope bounds to obtain $\abs{\phi(s)-\phi(t)}\leq\abs{s-t}$. Thus it is continuous. The geometric majorant $\sum(3/4)^n$ proves uniform convergence of the defining series and hence continuity of $F$.

Fix $x\in\R$. For every integer $m\geq1$, choose
\[h_m=\pm\tfrac12\,4^{-m}\]
so that the segment between $4^mx$ and $4^m(x+h_m)$ contains no integer in its interior. This is possible because within the unit interval containing $4^mx$, at least one side has available length at least $1/2$; at an integer choose either direction. On that segment the tent is affine, so
\[\gamma_{n,m}=\frac{\phi(4^n(x+h_m))-\phi(4^nx)}{h_m}\quad\text{satisfies}\quad
\abs{\gamma_{m,m}}=4^m,\qquad\abs{\gamma_{n,m}}\leq4^n\ (n<m).\]
If $n>m$, the displacement $4^nh_m=\pm2\cdot4^{n-m-1}$ is an even integer, so periodicity gives $\gamma_{n,m}=0$ exactly. Therefore the quotient of the uniformly convergent sum is here an actual finite sum:
\[\begin{split}
\abs{\frac{F(x+h_m)-F(x)}{h_m}}
&=\abs{\sum_{n=0}^m(3/4)^n\gamma_{n,m}}\\
&\geq3^m-\sum_{n=0}^{m-1}3^n=\tfrac12(3^m+1)\longrightarrow\infty.
\end{split}\]
Since $h_m\to0$, a finite derivative at this fixed $x$ is impossible. The argument holds at every $x$. No differentiation of an uncontrolled infinite series was used. The coefficients and signed half-cell increments are those of Rudin Theorem 7.18, verified against its displayed scan.
\end{proof}

% BEGIN REVISED SOURCE chapters/M03.18.tex SHA256=d27df2cd5543a89153a42131cc3fcbe9664f80c73bcccc35896eac968bbbf6b9
\chapter{Classical special integrals}
\label{module:M03.18}
\SourceRef{Zorich II, section 17.3; Rudin, section 8, Gamma-function part;
Olver, chapter 2 (real positive arguments only)}

\KEntry{K03.18.beta_definition}{The Beta integral and its domain}
\begin{definition}
For $p,q>0$ put
\[
 B(p,q)=\int_0^1 t^{p-1}(1-t)^{q-1}\dd t.
\]
This is an improper integral, with the two endpoints checked independently.
\end{definition}
\begin{proof}[Existence and symmetry]
On $(0,1/2]$ the second factor is bounded above and below by positive constants,
so convergence there follows from $\int_0^{1/2}t^{p-1}\dd t<\infty$.
On $[1/2,1)$ use the same argument with $(1-t)^{q-1}$. Both tests require the
indicated positive parameter; no linked truncation is being used. The integrand
is positive, so $B(p,q)>0$. Substitution $u=1-t$ first on regular compact
subintervals and then at each endpoint gives $B(p,q)=B(q,p)$.
\end{proof}

\KEntry{K03.18.beta_recurrence}{Beta recurrences}
\begin{proposition}
For $p,q>0$,
\[
 B(p+1,q)=\frac p{p+q}B(p,q),\qquad
 B(p,q+1)=\frac q{p+q}B(p,q),\qquad B(p,1)=\frac1p.
\]
Consequently, for integers $n\geq0$,
\begin{equation}\label{eq:beta-integer}
 B(s,n+1)=\frac{n!}{s(s+1)\cdots(s+n)}\quad(s>0).
\end{equation}
\end{proposition}
\begin{proof}
Adding the two integrands gives $B(p+1,q)+B(p,q+1)=B(p,q)$.
The derivative of $t^p(1-t)^q$ is
$pt^{p-1}(1-t)^q-qt^p(1-t)^{q-1}$. Integrating on $[\eta,1-\delta]$ and sending
$\eta,\delta$ independently to zero gives
$pB(p,q+1)=qB(p+1,q)$, since the boundary values tend to zero and both
derivative terms are integrable. Solving these two linear equations gives the
recurrences. Direct evaluation of $\int_0^1t^{p-1}\dd t$ gives $1/p$.
Iterate the recurrence in the second parameter to obtain
\eqref{eq:beta-integer}; this is finite induction, not a Gamma identity.
\end{proof}

\KEntry{K03.18.gamma_definition}{The real Gamma integral}
\begin{definition}
For $s>0$ define
\[
 \Gamma(s)=\int_0^\infty t^{s-1}e^{-t}\dd t.
\]
\end{definition}
\begin{proof}[Existence and positivity]
On $(0,1]$ the integrand is bounded by $t^{s-1}$, whose improper integral is
$1/s$. At infinity every fixed power is eventually dominated by an exponential:
choose an integer $m>s$ and use the exponential series, or
\KRef{K03.02.growth_scale}, to bound $t^{s-1}e^{-t}\leq C e^{-t/2}$ for
$t\geq1$. The latter function has a finite integral. The two positive improper
pieces therefore converge independently. The integrand is strictly positive on
every nonempty compact subinterval of $(0,\infty)$, so $\Gamma(s)>0$.
\end{proof}

\KEntry{K03.18.gamma_recurrence}{Recurrence and factorial values}
\begin{theorem}
For $s>0$, $\Gamma(s+1)=s\Gamma(s)$; moreover $\Gamma(1)=1$ and
$\Gamma(n+1)=n!$ for every integer $n\geq0$.
\end{theorem}
\begin{proof}
On $[\eta,R]$ integrate $t^se^{-t}$ by parts:
\[
 \int_\eta^R t^se^{-t}\dd t
 =[-t^se^{-t}]_\eta^R+s\int_\eta^R t^{s-1}e^{-t}\dd t.
\]
The upper boundary tends to zero by exponential domination; the lower boundary
tends to zero since $s>0$. Both integrals exist by the preceding entry, so the
recurrence follows from separate endpoint limits. The elementary integral of
$e^{-t}$ is $1$ on $(0,\infty)$. Induction now proves the factorial formula,
including $0!=1$.
\end{proof}

\KEntry{K03.18.gamma_euler_limit}{Euler's one-variable limit}
\begin{theorem}
For every $s>0$,
\begin{equation}\label{eq:gamma-euler}
 \Gamma(s)=\lim_{n\to\infty}n^sB(s,n+1)
 =\lim_{n\to\infty}\frac{n!n^s}{s(s+1)\cdots(s+n)}.
\end{equation}
\end{theorem}
\begin{proof}
Substitute $u=nt$ in the Beta integral:
\[
 n^sB(s,n+1)=\int_0^n u^{s-1}(1-u/n)^n\dd u.
\]
For $0\leq u\leq n$ the logarithmic inequality $\log(1-v)\leq-v$ gives
$(1-u/n)^n\leq e^{-u}$; the value at $u=n$ is interpreted by continuity.
On every fixed $[\eta,R]\subset(0,\infty)$,
$n\log(1-u/n)=-u+O(R^2/n)$ uniformly, so the integrands tend uniformly to
$u^{s-1}e^{-u}$. This uniform logarithmic estimate follows, for example, by
integrating $1/(1-v)-1$ on $0\leq v\leq R/n$; its absolute integral is at most
$2(R/n)^2$ for $n\geq2R$.

Here is the full improper-limit argument. On $(0,\eta)$ both nonnegative
integrands are bounded by $u^{s-1}$, so their integrals are at most $\eta^s/s$.
For $u\geq R$ their integrals are bounded by the tail of
$u^{s-1}e^{-u}$, uniformly in $n$, with the first integrand set to zero for
$u>n$. Given $\eps>0$, choose $\eta$ small and $R$ large to make these endpoint
contributions below $\eps$, and then use uniform convergence on $[\eta,R]$.
This proves the limit with ordinary proper integrals and tail estimates.
Equation \eqref{eq:beta-integer} supplies the product formula. No
multidimensional integration theorem is used.
\end{proof}

\KEntry{K03.18.gamma_regular}{Differentiation of Gamma}
\begin{theorem}
The function $\Gamma$ is $C^\infty$ on $(0,\infty)$ and, for $k\geq0$,
\[
 \Gamma^{(k)}(s)=\int_0^\infty t^{s-1}(\log t)^k e^{-t}\dd t.
\]
\end{theorem}
\begin{proof}
Fix $0<\alpha\leq s\leq\beta<\infty$. On $(0,1]$ the absolute value of the
$k$th parameter derivative is at most
$t^{\alpha-1}|\log t|^k$. Substitution $t=e^{-v}$ turns its integral into
$\int_0^\infty e^{-\alpha v}v^k\dd v$, which converges by the Gamma
existence estimate (or repeated integration by parts). On $[1,\infty)$ use
$t^{\beta-1}(\log t)^ke^{-t}\leq C_k e^{-t/2}$. These two integrable bounds
also apply to the next derivative. The tails are uniformly small in $s$.
On each compact $t$-interval the integrands and their parameter derivatives are
jointly continuous. Apply \KRef{K03.14.improper_derivative} separately at the
two ends, and iterate. The same uniform-tail argument gives continuity of each
derived integral. Since every positive parameter lies in such an interval,
the asserted smoothness follows.
\end{proof}

\KEntry{K03.18.gamma_logconvex}{Log-convexity}
\begin{theorem}
$\log\Gamma$ is convex on $(0,\infty)$.
\end{theorem}
\begin{proof}
For $0<\theta<1$ write
\[
 t^{(1-\theta)p+\theta q-1}e^{-t}
 =(t^{p-1}e^{-t})^{1-\theta}(t^{q-1}e^{-t})^\theta.
\]
Holder's inequality on $[\eta,R]$, with exponents $1/(1-\theta)$ and
$1/\theta$, bounds this integral by
$(\int_\eta^R t^{p-1}e^{-t})^{1-\theta}
 (\int_\eta^R t^{q-1}e^{-t})^\theta$.
All three positive integrals have independent finite endpoint limits, so passage
to those limits gives
$\Gamma((1-\theta)p+\theta q)\leq\Gamma(p)^{1-\theta}\Gamma(q)^\theta$.
Taking logarithms is legitimate by positivity. The cases $\theta=0,1$ are
equalities. This is exactly convexity, using only the properly proved
Riemann integral inequality \KRef{K03.09.holder}.
\end{proof}

\KEntry{K03.18.bohr_mollerup}{The Bohr--Mollerup characterization}
\begin{theorem}
If $F:(0,\infty)\to(0,\infty)$ satisfies $F(1)=1$,
$F(x+1)=xF(x)$ and convexity of $\log F$, then $F=\Gamma$.
\end{theorem}
\begin{proof}
Put $G=\log F$. Recurrence gives $F(n)=(n-1)!$.
For $0<x\leq1$ and integer $n\geq2$, the ordered-secant-slope inequalities
for the convex function $G$ yield
\[
 \log(n-1)\leq\frac{G(n+x)-G(n)}x\leq\log n.
\]
The upper inequality uses $[n,n+x]\subset[n,n+1]$; the lower one compares
with $[n-1,n]$. Since
$F(n+x)=F(x)\prod_{j=0}^{n-1}(x+j)$, exponentiating gives
\[
 \frac{(n-1)!(n-1)^x}{\prod_{j=0}^{n-1}(x+j)}
 \leq F(x)\leq
 \frac{(n-1)!n^x}{\prod_{j=0}^{n-1}(x+j)}.
\]
The left member tends to $\Gamma(x)$ by \eqref{eq:gamma-euler}; the quotient
of the upper and lower members is $(n/(n-1))^x\to1$. Squeeze proves
$F(x)=\Gamma(x)$ on $(0,1]$, and recurrence extends it to all positive $x$.
Gamma itself has all three properties, already proved above. Recurrence and
normalization alone are not unique: multiplication of Gamma by a positive
nonconstant period-one function equal to one at $1$ preserves them.
\end{proof}

\KEntry{K03.18.beta_gamma}{The Beta--Gamma identity without a hidden double integral}
\begin{theorem}
For $p,q>0$,
\[
 B(p,q)=\frac{\Gamma(p)\Gamma(q)}{\Gamma(p+q)}.
\]
\end{theorem}
\begin{proof}
The Beta recurrence gives
$B(p+n,q)=(p)_n B(p,q)/(p+q)_n$, where
$(p)_n=p(p+1)\cdots(p+n-1)$.
We first prove $n^qB(p+n,q)\to\Gamma(q)$.
With $u=n(1-t)$ the left member is
\[
 \int_0^n u^{q-1}(1-u/n)^{n+p-1}\dd u.
\]
For $n\geq2$, $n+p-1\geq n/2$, so this integrand is bounded by
$u^{q-1}e^{-u/2}$. On each $[\eta,R]$ it converges uniformly to
$u^{q-1}e^{-u}$. At zero use the bound $u^{q-1}$, and at infinity use
the displayed exponential bound. The three-interval argument in
\KRef{K03.18.gamma_euler_limit} therefore proves the claimed limit.

The quotient of the Euler product approximants for $\Gamma(p+q)$ and
$\Gamma(p)$, with index $n-1$, shows
\[
 (n-1)^q\frac{(p)_n}{(p+q)_n}\longrightarrow
 \frac{\Gamma(p+q)}{\Gamma(p)}.
\]
Every quantity is positive and the limit is nonzero, so inversion and
$(n/(n-1))^q\to1$ give
$\frac{(p+q)_n}{n^q(p)_n}\to\Gamma(p)/\Gamma(p+q)$.
Multiply this by the established limit of $n^qB(p+n,q)$ and use the exact
finite recurrence identity. This yields the formula without a polar,
Jacobian or full Fubini theorem.
\end{proof}

\KEntry{K03.18.wallis}{Wallis integrals and the product}
\begin{theorem}
Let $I_n=\int_0^{\pi/2}\sin^n x\dd x$, $n\geq0$. Then $I_0=\pi/2$,
$I_1=1$, $I_n=(n-1)I_{n-2}/n$ for $n\geq2$, and
\[
 \prod_{k=1}^n\frac{(2k)^2}{(2k-1)(2k+1)}\longrightarrow\frac\pi2.
\]
Writing $c_n=\binom{2n}{n}/4^n$, we also have
$\sqrt n\,c_n\to1/\sqrt\pi$.
\end{theorem}
\begin{proof}
The first two integrals follow from elementary primitives. Integrating
$\sin^{n-1}x\sin x$ by parts for $n\geq2$ gives
$I_n=(n-1)\int\sin^{n-2}x\cos^2x\dd x=(n-1)(I_{n-2}-I_n)$.
Induction gives
\[
 I_{2n}=\frac\pi2c_n,\qquad I_{2n+1}=\frac1{(2n+1)c_n}.
\]
Since $0\leq\sin x\leq1$,
$I_{2n+2}\leq I_{2n+1}\leq I_{2n}$.
The quotient $I_{2n+2}/I_{2n}=(2n+1)/(2n+2)\to1$, hence
$I_{2n+1}/I_{2n}\to1$. This quotient equals $2W_n/\pi$, where $W_n$
is the displayed finite product, proving Wallis' limit. Further,
$I_{2n}I_{2n+1}=\pi/[2(2n+1)]$. Combined with their quotient limit this
gives $I_{2n}\sim\sqrt\pi/(2\sqrt n)$ and the stated limit for $c_n$.
\end{proof}

\KEntry{K03.18.gaussian}{The Gaussian integral by a one-variable route}
\begin{theorem}
$\Gamma(1/2)=\sqrt\pi$ and $\int_{-\infty}^{\infty}e^{-x^2}\dd x=\sqrt\pi$.
\end{theorem}
\begin{proof}
The Euler approximant at $s=1/2$ simplifies algebraically to
\[
 \frac{n!\sqrt n}{(1/2)(3/2)\cdots(n+1/2)}
 =\frac{2\sqrt n}{(2n+1)c_n}\longrightarrow\sqrt\pi
\]
by Wallis' limit. It therefore equals $\Gamma(1/2)$ in the limit.
Substitution $t=x^2$ on positive regular intervals gives
$\Gamma(1/2)=2\int_0^\infty e^{-x^2}\dd x$, including the endpoint limits.
Evenness gives the two-sided formula. Both tails are absolutely convergent;
there is no principal-value interpretation and no use of a two-dimensional
coordinate substitution.
\end{proof}

\KEntry{K03.18.dirichlet_integral}{The Dirichlet sine integral}
\begin{theorem}
With $\sin x/x$ assigned its continuous value $1$ at zero,
\[
 \int_0^\infty\frac{\sin x}{x}\dd x=\frac\pi2.
\]
\end{theorem}
\begin{proof}
For $a>0$ define $J(a)=\int_0^\infty e^{-ax}\sin x/x\dd x$.
Near zero the integrand is bounded; at infinity it is absolutely integrable.
On a compact parameter interval $a\geq\alpha>0$, the parameter derivative
is $-e^{-ax}\sin x$, bounded by $e^{-\alpha x}$. Thus the proved improper
parameter differentiation rule gives
\[
 J'(a)=-\int_0^\infty e^{-ax}\sin x\dd x=-\frac1{1+a^2}.
\]
The last integral follows by evaluating the primitive
$e^{-ax}(-a\sin x-\cos x)/(1+a^2)$ at zero and infinity.
Also $|J(a)|\leq\int_0^\infty e^{-ax}\dd x=1/a$, since $|\sin x|\leq x$.
The function $\arctan(1/a)$ has the same derivative and tends to zero as
$a\to\infty$; constancy of their difference gives $J(a)=\arctan(1/a)$.

It remains to remove damping, rather than assume it legitimate. For $a\geq0$
and $R>0$, set $g_a(x)=e^{-ax}/x$. It is nonnegative decreasing, with
$g_a(R)\leq1/R$. The primitive $F_R(y)=\int_R^y\sin x\dd x$ has absolute
value at most $2$. Integration by parts gives, for $S>R$,
\[
 \left|\int_R^S g_a(x)\sin x\dd x\right|
 \leq2g_a(S)+2(g_a(R)-g_a(S))\leq\frac2R.
\]
The same bound holds on the infinite tail, uniformly even at $a=0$;
Dirichlet convergence at $a=0$ is included. On $[0,R]$ the damped integrand
converges uniformly to $\sin x/x$ as $a\downarrow0$. Split at $R$, make both
tails small, and then pass to the limit on the compact interval. Therefore
the undamped integral is $\lim_{a\downarrow0}J(a)=\pi/2$.
\end{proof}

\KEntry{K03.18.other_special_functions}{Boundary of the special-function supplement}
\begin{remark}[Routing record, not a new asserted theorem]
Incomplete Gamma, error and Fresnel integrals can be studied after their specific
convergence and error estimates are supplied. General Bessel, Airy, Legendre and
hypergeometric theories belong to D13/D24. Their names here do not claim that
their defining equations, existence theory or asymptotic formulas have been proved
in these notes. No proof above depends on such a theory.
\end{remark}

% BEGIN REVISED SOURCE chapters/M03.17.tex SHA256=6649e202f8e46925fc24be46a9d336e520bc559f1554e57c2ee31504e61131e2
\chapter{Asymptotic expansions and summation estimates}
\label{module:M03.17}
\SourceRef{Zorich I, Appendix D; Zorich II, chapter 19; selected real-variable
parts of Olver, chapters 1, 3 and 8}

\KEntry{K03.17.asymptotic_scale}{Scales and finite expansions}
\begin{definition}
Along a specified limit, an asymptotic scale is a family $\phi_0,\phi_1,\ldots$
of eventually nonzero scalar functions such that $\phi_{j+1}=o(\phi_j)$.
An expansion through index $m$ means the finite assertion
\[
 f=\sum_{j=0}^{m}a_j\phi_j+o(\phi_m).
\]
An all-orders expansion means this assertion for every finite $m$, with
consistent coefficients. It does not assert convergence of an infinite sum.
\end{definition}
For $j<k$, $\phi_k/\phi_j$ is a finite product of adjacent ratios tending to
zero and hence tends to zero. The domain and direction of the limit are part
of the definition; a two-sided real power scale often needs $|x|$ in estimates.

\KEntry{K03.17.coefficient_unique}{Uniqueness does not determine the function}
\begin{proposition}
The coefficients of an expansion in a fixed asymptotic scale are unique.
Functions can nevertheless have the same all-orders expansion.
\end{proposition}
\begin{proof}
Subtract two expansions of the same function and let $j$ be the least index
with unequal coefficients. Divide by $\phi_j$. All terms of larger index
tend to zero, as does the remainder divided by $\phi_j$; the nonzero constant
difference at index $j$ cannot tend to zero. This is a contradiction.
Uniqueness concerns coefficients of a given function, not uniqueness of that
function; \KRef{K03.17.flat_remainder} gives the latter distinction.
\end{proof}

\KEntry{K03.17.operations}{Arithmetic with controlled remainders}
\begin{proposition}
Finite asymptotic expansions can be added and multiplied by retaining the
resulting remainders. Division is valid after an eventually nonzero leading
term has been factored out. Closure of an arbitrary scale under multiplication
or composition is not automatic.
\end{proposition}
\begin{proof}
Write $f=P+r$, $g=Q+s$, where $P,Q$ are the stated finite sums. The exact
identities
\[
 f+g=P+Q+r+s,\qquad fg=PQ+Ps+Qr+rs
\]
reduce every truncation to the $O/o$ multiplication rules. For example, if
$f=\sum_{j=0}^m a_jt^j+o(t^m)$ and
$g=\sum_{j=0}^m b_jt^j+o(t^m)$ as $t\to0$, then $P,Q=O(1)$;
the coefficient at $t^k$ is $\sum_{j=0}^ka_jb_{k-j}$ for $k\leq m$, and
the discarded finite polynomial terms of degree $>m$ and the three remainder
terms are $o(t^m)$. For a nonzero leading coefficient, divide first by the
leading scale factor and write the remaining denominator as $a(1+u)$,
$a\ne0$, $u\to0$. The finite identity
\[
 \frac1{1+u}=\sum_{j=0}^{m}(-u)^j+
 \frac{(-u)^{m+1}}{1+u}
\]
both proves eventual invertibility and controls the quotient remainder.
Composition requires the inner variable to approach the expansion point and
its range to remain in the outer domain. For instance $u(t)=O(t)$ and an outer
remainder $o(u^m)$ imply a remainder $o(t^m)$ by the defining epsilon bound,
including points with $u=0$ after the remainder is assigned its correct value.
For a nonpower scale one must keep its product scales or explicitly justify
their order before dropping terms.
\end{proof}

\KEntry{K03.17.flat_remainder}{Beyond all powers}
\begin{example}
The function $e^{-1/x^2}$ for $x\ne0$, extended by zero at $0$, is
$o(|x|^m)$ for every nonnegative integer $m$, but is positive off $0$.
Thus its all-orders power expansion is the zero expansion without the function
being zero on a neighborhood.
\end{example}
\begin{proof}
Put $u=1/x^2$. For an integer $r>m/2$, the exponential series gives
$e^u\geq u^r/r!$ for $u>0$, so
$u^{m/2}e^{-u}\leq r!u^{m/2-r}\to0$. This proves the claimed estimates.
The full smooth-flat construction, including all derivatives at zero, is
\KRef{K03.06.smooth_flat}; no analyticity inference is made here.
\end{proof}

\KEntry{K03.17.divergent_asymptotic}{A divergent series that is asymptotically correct}
\begin{example}
For $x>0$ let $F(x)=\int_0^\infty e^{-t}/(1+xt)\dd t$. For every integer
$N\geq1$,
\[
 F(x)=\sum_{k=0}^{N-1}(-1)^kk!x^k+R_N(x),\qquad
 |R_N(x)|\leq N!x^N.
\]
The associated infinite power series has radius zero.
\end{example}
\begin{proof}
The integral is absolutely convergent by comparison with $e^{-t}$. Insert
the exact finite geometric identity into its integrand. Each moment is
$\Gamma(k+1)=k!$, and the remainder has absolute value at most
$x^N\int_0^\infty t^Ne^{-t}\dd t=N!x^N$, since $1+xt\geq1$.
For fixed $x\ne0$, consecutive absolute power-series terms have ratio
$(k+1)|x|\to\infty$, so they do not even tend to zero. The finite estimates
remain valid and, using one extra term, give an all-orders asymptotic expansion.
\end{proof}

\KEntry{K03.17.bernoulli}{Bernoulli polynomials and their convention}
\begin{definition}
Set $B_0(x)=1$. Recursively, $B_n$ is the unique polynomial satisfying
$B_n'=nB_{n-1}$ and $\int_0^1B_n(x)\dd x=0$ for $n\geq1$.
Write $B_n=B_n(0)$ for the Bernoulli number. In particular,
$B_1(x)=x-1/2$, so $B_1=-1/2$.
\end{definition}
\begin{proof}[Construction and identities]
Integrate the already constructed polynomial $nB_{n-1}$ and choose the unique
constant giving zero average. This proves existence and uniqueness. It gives
$B_2(x)=x^2-x+1/6$ and $B_4(x)=x^4-2x^3+x^2-1/30$ by coefficient comparison.
For $n\geq2$,
$B_n(1)-B_n(0)=n\int_0^1B_{n-1}=0$.
The polynomials $(-1)^nB_n(1-x)$ satisfy the same derivative and average
recursion, so uniqueness gives $B_n(1-x)=(-1)^nB_n(x)$.
Consequently $B_n(0)=0$ for odd $n\geq3$.

Induction on $n$, by differentiating and then evaluating at $x=0$, gives
$B_n(x+1)-B_n(x)=nx^{n-1}$; the case $n=1$ is direct and the endpoint
identity fixes the induction constant for $n\geq2$. Telescoping therefore
proves the finite power-sum formula
\[
 \sum_{k=0}^{N-1}k^p=
 \frac{B_{p+1}(N)-B_{p+1}(0)}{p+1}\quad(p\geq0).
\]
For $p=0$ interpret the summands as $1$, including the $k=0$ term.
\end{proof}

\KEntry{K03.17.euler_maclaurin}{Euler--Maclaurin with an actual remainder}
\begin{theorem}
Let integers $a<b$, $m\geq1$, and $f\in C^{2m}([a,b],\R\text{ or }\C)$.
Then
\begin{align*}
 \frac{f(a)}2+\sum_{a<k<b}f(k)+\frac{f(b)}2
 &=\int_a^b f(x)\dd x+
 \sum_{r=1}^m\frac{B_{2r}}{(2r)!}
 \bigl(f^{(2r-1)}(b)-f^{(2r-1)}(a)\bigr)\\
 &\quad-\frac1{(2m)!}\int_a^b
 B_{2m}(\{x\})f^{(2m)}(x)\dd x,
\end{align*}
where $\{x\}=x-\lfloor x\rfloor$.
\end{theorem}
\begin{proof}
Integrating $(x-k-1/2)f'(x)$ by parts on $[k,k+1]$ gives
\[
 \frac{f(k)+f(k+1)}2-\int_k^{k+1}f
 =\int_k^{k+1}B_1(x-k)f'(x)\dd x.
\]
Sum over the unit intervals. Put
$T_j=\int_a^b B_j(\{x\})f^{(j)}(x)\dd x/j!$.
For $1\leq j<2m$, integrate by parts on each unit interval and sum.
The equality $B_{j+1}(1)=B_{j+1}(0)$ cancels internal boundary terms, yielding
\[
 T_j=\frac{B_{j+1}}{(j+1)!}
 \bigl(f^{(j)}(b)-f^{(j)}(a)\bigr)-T_{j+1}.
\]
Odd Bernoulli numbers beyond $B_1$ vanish. Starting from $T_1$ and applying
this identity successively leaves exactly the even-index corrections and
$-T_{2m}$ in the theorem. The periodic polynomials have at most finitely
many relevant junctions on $[a,b]$; their endpoint values at those junctions
do not change their proper Riemann integrals. The stated regularity permits
every integration by parts used above.
\end{proof}

\KEntry{K03.17.em_bound}{A useful finite error bound}
\begin{corollary}
With $M_{2m}=\max_{0\leq t\leq1}|B_{2m}(t)|$, the absolute remainder is
at most
\[
 \frac{M_{2m}}{(2m)!}\int_a^b|f^{(2m)}(x)|\dd x.
\]
\end{corollary}
\begin{proof}
Take the absolute value inside the remainder integral and use the uniform
polynomial bound $M_{2m}$. This finite bound supplies no convergence
claim for an infinite Euler--Maclaurin series.
\end{proof}

\KEntry{K03.17.harmonic}{Euler's constant and harmonic corrections}
\begin{proposition}
There is a real constant $\gamma$ such that
\[
 H_n=\log n+\gamma+\frac1{2n}-\frac1{12n^2}
       +\frac1{120n^4}+O(n^{-6}),\qquad H_n=\sum_{k=1}^n\frac1k.
\]
More generally each fixed number of Bernoulli corrections has a controlled
inverse-power remainder.
\end{proposition}
\begin{proof}
Integral comparison gives $H_n\geq\log(n+1)>\log n$ and
$\log(1+1/n)>1/(n+1)$; hence $H_n-\log n$ is positive decreasing and has
a limit $\gamma$. Apply Euler--Maclaurin to $f(x)=1/x$ on $[n,M]$.
Its $(j)$th derivative is $(-1)^jj!x^{-j-1}$. The endpoint-half sum is
$H_M-H_n+1/(2n)-1/(2M)$. Sending $M\to\infty$ gives
\[
 H_n=\log n+\gamma+\frac1{2n}
 -\sum_{r=1}^m\frac{B_{2r}}{2r\,n^{2r}}+E_m(n),
 \qquad |E_m(n)|\leq\frac{M_{2m}}{2m\,n^{2m}}.
\]
The bound follows by integrating $x^{-2m-1}$ and justifies the improper
limit of the remainder. Taking the formula with $m=3$, and absorbing its
last correction and remainder into $O(n^{-6})$, proves the displayed
formula, since $B_2=1/6$ and $B_4=-1/30$. In general one more correction
than those displayed gives the desired order of error.
\end{proof}

\KEntry{K03.17.stirling}{Stirling's formula, including the constant}
\begin{theorem}
As $n\to\infty$ through positive integers,
\[
 n!=\sqrt{2\pi n}(n/e)^n\exp(O(1/n))
     \sim\sqrt{2\pi n}(n/e)^n.
\]
\end{theorem}
\begin{proof}
Put $s_n=\log(n!)-(n+1/2)\log n+n$. A direct subtraction gives
\[
 s_{n+1}-s_n=1-(n+1/2)\log(1+1/n)
            =-\frac1{12n^2}+O(n^{-3}).
\]
The last equality follows from the finite Taylor expansion of $\log(1+u)$
with a bounded fourth derivative on $[0,1/2]$; the finitely many smaller
indices do not matter. Thus $\sum|s_{n+1}-s_n|$ converges and its tail is
$O(1/n)$ by the $p$-series integral bound. Write $s_n=c+O(1/n)$ and
$A=e^c>0$. Exponentiation gives
$n!=A n^{n+1/2}e^{-n}\exp(O(1/n))$.

To determine $A$, form $c_n=\binom{2n}{n}/4^n$. The preceding formula yields
$\sqrt n\,c_n\to\sqrt2/A$. Independently, the Wallis calculation
\KRef{K03.18.wallis} gives $\sqrt n\,c_n\to1/\sqrt\pi$.
Uniqueness of a finite limit gives $A=\sqrt{2\pi}$. This identifies the
normalization rather than hiding it in an unspecified multiplicative constant.
The logarithmic $O(1/n)$ error also implies a relative $1+O(1/n)$ error,
by the mean-value bound for $e^u$ near zero.
\end{proof}

\KEntry{K03.17.laplace_local}{An interior Laplace minimum}
\begin{theorem}
Let $a<b$, let $\phi\in C^2([a,b],\R)$ have a unique global minimum at
$x_0\in(a,b)$, and suppose $d=\phi''(x_0)>0$. For continuous scalar $g$,
\[
 \sqrt\lambda\,e^{\lambda\phi(x_0)}
 \int_a^b e^{-\lambda\phi(x)}g(x)\dd x
 \longrightarrow g(x_0)\sqrt{\frac{2\pi}{d}}
 \quad(\lambda\to+\infty).
\]
The assertion includes $g(x_0)=0$; it is a normalized limit rather than an
invalid asymptotic equivalence to zero.
\end{theorem}
\begin{proof}
Fermat gives $\phi'(x_0)=0$. Taylor--Peano gives
$q(x_0+h):=\phi(x_0+h)-\phi(x_0)=dh^2/2+o(h^2)$.
Choose $r>0$ with $[x_0-r,x_0+r]\subset(a,b)$ and
$q(x_0+h)\geq c h^2$ there for some $c>0$, for instance $c=d/4$.
On the complementary compact set uniqueness of the minimum gives
$q\geq\eta>0$. If $M=\max|g|$, its normalized contribution is at most
$\sqrt\lambda M(b-a)e^{-\lambda\eta}\to0$.

In the local integral substitute $y=\sqrt\lambda(x-x_0)$. The normalized
integrand becomes
$g(x_0+y/\sqrt\lambda)e^{-\lambda q(x_0+y/\sqrt\lambda)}$ on
$|y|\leq r\sqrt\lambda$, bounded in absolute value by $Me^{-cy^2}$.
For fixed $R$, the defining $o(h^2)$ bound is uniform for
$|h|\leq R/\sqrt\lambda$ and proves uniform convergence on $[-R,R]$ to
$g(x_0)e^{-dy^2/2}$. The tails outside this fixed interval are bounded by
$M\int_{|y|>R}e^{-cy^2}\dd y$, uniformly in large $\lambda$.
They tend to zero as $R\to\infty$, by the Gaussian integral and a linear
substitution. First choose $R$ for the tails, then $\lambda$ for the proper
integrals. This proves convergence to
$g(x_0)\int_\R e^{-dy^2/2}\dd y=g(x_0)\sqrt{2\pi/d}$.
The argument uses explicit proper convergence and tails, not an unstated
dominated-convergence theorem.
\end{proof}
\begin{corollary}[A quantitative source variant]
If additionally $\phi$ is $C^3$ near $x_0$ and $g$ is locally Lipschitz there,
the normalized expression equals
$g(x_0)\sqrt{2\pi/d}+O(\lambda^{-1/2})$.
\end{corollary}
\begin{proof}
In the preceding variables,
$|\lambda q(x_0+y/\sqrt\lambda)-dy^2/2|
 \leq C|y|^3/\sqrt\lambda$ and
$|g(x_0+y/\sqrt\lambda)-g(x_0)|\leq L|y|/\sqrt\lambda$.
Both exponents are at least $c'y^2$ after reducing $r$.
The elementary inequality $|e^{-u}-e^{-v}|\leq|u-v|e^{-\min(u,v)}$
therefore bounds the local difference by
$C'\lambda^{-1/2}(|y|+|y|^3)e^{-c'y^2}$.
Its integral is finite, by substitution $t=c'y^2$ and Gamma moments.
The omitted Gaussian and exterior tails are exponentially small in
$\lambda$, proving the error estimate. These stronger hypotheses are not
attached to the weaker normalized-limit theorem.
\end{proof}

\KEntry{K03.17.watson}{A real Watson lemma with tail conditions}
\begin{theorem}
Let $N\geq1$, $\alpha>0$. Suppose $q$ is locally Riemann integrable on
$(0,\infty)$ and, for $0<t\leq\delta$,
\[
 q(t)=\sum_{k=0}^{N-1}a_kt^{\alpha+k-1}
             +O(t^{\alpha+N-1}).
\]
Suppose also $|q(t)|\leq C e^{ct}$ for all sufficiently large $t$.
Then, as $\lambda\to+\infty$,
\[
 \int_0^\infty e^{-\lambda t}q(t)\dd t
 =\sum_{k=0}^{N-1}\frac{a_k\Gamma(\alpha+k)}{\lambda^{\alpha+k}}
       +O(\lambda^{-\alpha-N}).
\]
\end{theorem}
\begin{proof}
The local expansion bounds $q$ by a constant times $t^{\alpha-1}$ near zero,
so the integral there converges absolutely. The growth bound controls the
infinite tail for $\lambda>c$. On $[\delta,A]$ local integrability gives a
finite absolute integral, bounded after damping by
$e^{-\lambda\delta}\int_\delta^A|q|$. On $[A,\infty)$ the bound is
$Ce^{-(\lambda-c)A}/(\lambda-c)$. Thus the whole tail beyond $\delta$ is
exponentially small for large $\lambda$; so are the tails of each of the
finitely many power monomials beyond $\delta$.

The local remainder contributes at most
\[
 C'\int_0^\delta e^{-\lambda t}t^{\alpha+N-1}\dd t
 \leq C'\Gamma(\alpha+N)\lambda^{-\alpha-N}.
\]
For each retained monomial, substitute $u=\lambda t$ in its full improper
integral to obtain $\Gamma(\alpha+k)\lambda^{-\alpha-k}$.
Adding the finite contributions and the controlled tails proves the formula.
The global growth condition is what prevents a formally correct local
expansion from being misused for an uncontrolled infinite integral.
\end{proof}

\KEntry{K03.17.oscillatory_bound}{An elementary oscillatory estimate}
\begin{proposition}
For $h\in C^1([a,b],\C)$ and nonzero real $\lambda$,
\[
 \left|\int_a^b h(x)e^{i\lambda x}\dd x\right|
 \leq\frac{|h(a)|+|h(b)|+\int_a^b|h'(x)|\dd x}{|\lambda|}.
\]
\end{proposition}
\begin{proof}
Real-parameter differentiation of the established complex exponential gives
$(e^{i\lambda x})'=i\lambda e^{i\lambda x}$. Integration by parts yields
the boundary term $[h(x)e^{i\lambda x}/(i\lambda)]_a^b$ and the remaining
integral $-(i\lambda)^{-1}\int h'e^{i\lambda x}$. Take absolute values and
use $|e^{i\lambda x}|=1$.
\end{proof}

\KEntry{K03.17.scope_boundary}{What this asymptotic supplement does not assert}
\begin{remark}[Routing record]
General stationary phase, steepest descent, Stokes transitions, WKB and uniform
Airy/Bessel approximations require additional D12/D13/D24 foundations and
estimates. This record routes those subjects and contains no claimed proof of
their general theories. All statements proved in this chapter are finite
real-variable estimates with explicit remainders or tail control.
\end{remark}

% BEGIN REVISED SOURCE chapters/M03.19.tex SHA256=5a56aa4148a66137f2b0f36f1badf3d7c98c67d9a3ebf1e601d0dca6a068d05f
\chapter{Classical Fourier series on an interval}
\label{module:M03.19}
\SourceRef{Zorich II, sections 18.1--18.2; Apostol, Fourier-series part of
chapter 11; Rudin, section 8 and the Fejer exercise}

\KEntry{K03.19.periodic_setup}{A fixed Fourier normalization}
\begin{definition}
Let $f:\R\to\C$ be $2\pi$-periodic and Riemann integrable on one period.
For $n\in\Z$ set
\[
 c_n(f)=\frac1{2\pi}\int_{-\pi}^{\pi}f(x)e^{-inx}\dd x,
 \qquad S_N(f)(x)=\sum_{n=-N}^Nc_n(f)e^{inx}.
\]
For real $f$, put
$a_n=\pi^{-1}\int_{-\pi}^{\pi}f(x)\cos nx\dd x$ and
$b_n=\pi^{-1}\int_{-\pi}^{\pi}f(x)\sin nx\dd x$ for $n\geq1$,
and $a_0=\pi^{-1}\int f$. Then
$c_0=a_0/2$, $c_n=(a_n-ib_n)/2$, $c_{-n}=(a_n+ib_n)/2$.
\end{definition}
These formulas follow by inserting Euler's identity into the integrals.
Finite values at a period seam may be assigned consistently with periodicity;
changing such isolated values does not change any coefficient. One-sided
limits at the seam still have to be computed from both adjacent periods.

\KEntry{K03.19.orthogonality}{Trigonometric orthogonality}
\begin{proposition}
With $e_n(x)=e^{inx}$ and
$\langle f,g\rangle=(2\pi)^{-1}\int_{-\pi}^{\pi}f\overline g$,
we have $\langle e_n,e_m\rangle=\delta_{nm}$.
The real sine/cosine orthogonality formulas follow from this identity.
\end{proposition}
\begin{proof}
If $n=m$ the normalized integral is $1$. Otherwise a primitive of
$e^{i(n-m)x}$ is $e^{i(n-m)x}/(i(n-m))$, whose endpoint values are equal
because $n-m$ is an integer. Expanding sine and cosine as
$(e_n-e_{-n})/(2i)$ and $(e_n+e_{-n})/2$ gives zero mixed integrals;
the integral mean of $\cos^2 nx$ or $\sin^2 nx$ is $1/2$ for $n\geq1$.
All products are Riemann integrable. This pairing is a seminorm pairing
on arbitrary integrable functions: zero integral norm need not mean equality
at every point. No unconstructed $L^2$ quotient is being assumed.
\end{proof}

\KEntry{K03.19.finite_projection}{Finite projection and best approximation}
\begin{theorem}
For a finite $E\subset\Z$ let $P_Ef=\sum_{n\in E}c_n(f)e_n$. For every
$p=\sum_{n\in E}d_ne_n$,
\[
 \|f-p\|_2^2=\|f-P_Ef\|_2^2+\sum_{n\in E}|c_n(f)-d_n|^2,
 \qquad \|h\|_2^2=\frac1{2\pi}\int_{-\pi}^{\pi}|h|^2.
\]
\end{theorem}
\begin{proof}
For $n\in E$, linearity and orthogonality give
$\langle f-P_Ef,e_n\rangle=c_n(f)-c_n(f)=0$.
Expand $f-p=(f-P_Ef)+(P_Ef-p)$ and expand its squared absolute value
inside the integral. Both cross terms vanish by those finite orthogonality
relations, while the second square has integral
$\sum_{n\in E}|c_n-d_n|^2$. This proves the identity and its minimizing
property without any claim that the infinite trigonometric system is complete.
\end{proof}

\KEntry{K03.19.bessel}{Bessel's inequality}
\begin{theorem}
For scalar Riemann-integrable $f$ and every finite $E\subset\Z$,
$\sum_{n\in E}|c_n(f)|^2\leq\|f\|_2^2$.
Consequently $\sum_{n\in\Z}|c_n(f)|^2$ exists as a nonnegative sum and
has that bound.
\end{theorem}
\begin{proof}
Take $p=0$ in the finite projection identity to get
$\|f\|_2^2=\|f-P_Ef\|_2^2+\sum_E|c_n|^2$.
The first term is nonnegative. Increasing finite sets then have bounded
nonnegative sums, whose supremum defines the asserted countable sum.
Integrability of $|f|^2$ follows from boundedness and continuous-composition
closure of proper Riemann integrability.
\end{proof}

\KEntry{K03.19.riemann_lebesgue}{A Riemann--Lebesgue lemma with its needed extension}
\begin{theorem}
For Riemann-integrable $f$ on $[a,b]$,
$\int_a^b f(x)e^{i\lambda x}\dd x\to0$ as real $|\lambda|\to\infty$.
In particular $c_n(f)\to0$ as $|n|\to\infty$.
The same conclusion holds for a locally Riemann-integrable function on a
bounded interval with an absolutely convergent improper integral.
\end{theorem}
\begin{proof}
First take real bounded $f$. The Darboux criterion gives a partition with
$\sum_j\operatorname{osc}_{I_j}(f)|I_j|<\eps$. Choosing any tag value in
each interval defines a step function $s$ with $\int|f-s|<\eps$; endpoint
values are irrelevant. For complex $f$ do this to real and imaginary parts,
splitting the error tolerance, to obtain the same integral error bound.
For each fixed step interval,
\[
 \left|\int_u^v e^{i\lambda x}\dd x\right|\leq2/|\lambda|,
\]
so $|\int s e^{i\lambda x}|\leq2\sum_j|s_j|/|\lambda|\to0$.
The difference is bounded by $\int|f-s|<\eps$, uniformly in $\lambda$.
First let $|\lambda|\to\infty$ and then $\eps\downarrow0$.

For an absolutely convergent improper integral choose regular endpoints
whose omitted absolute tails total less than $\eps$. The first result
applies between them, and the omitted oscillatory integrals are bounded by
those same absolute tails. This proves the extension by truncation, with
no measure-theoretic convergence theorem.
\end{proof}

\KEntry{K03.19.dirichlet_kernel}{The Dirichlet kernel and convolution identity}
\begin{construction}
Define $D_N(t)=\sum_{n=-N}^Ne^{int}$. The finite geometric identity gives
\[
 D_N(t)=\frac{\sin((N+1/2)t)}{\sin(t/2)}\quad(t\notin2\pi\Z),
 \qquad D_N(0)=2N+1.
\]
It is even and real, and its integral mean on a period is $1$.
\end{construction}
\begin{proof}[Convolution formula]
Insert the coefficient integrals into the finite sum and interchange only
that finite sum with the integral:
\[
 S_N(f)(x)=\frac1{2\pi}\int_{-\pi}^{\pi}f(y)D_N(x-y)\dd y
 =\frac1{2\pi}\int_{-\pi}^{\pi}f(x-t)D_N(t)\dd t.
\]
The second equality uses a translation followed by periodicity: a periodic
integrable function has the same integral over any interval of its period,
as follows by splitting an interval at one seam and translating the other
piece. The sine quotient follows by multiplying the finite sum by
$e^{it/2}-e^{-it/2}$ and canceling all middle terms. The excluded values
are assigned by the finite sum, so no singular integral is introduced.
Finally orthogonality of the nonconstant exponentials gives
$(2\pi)^{-1}\int D_N=1$.
\end{proof}

\KEntry{K03.19.local_convergence}{Dini's local condition}
\begin{theorem}
Let $f$ be periodic Riemann integrable, with finite limits $L_-,L_+$ at
$x$ from the two sides. If for some $0<\delta<\pi$,
\[
 \int_0^\delta
 \frac{|f(x+t)+f(x-t)-L_+-L_-|}{t}\dd t<\infty,
\]
then $S_N(f)(x)\to(L_-+L_+)/2$.
In particular the conclusion holds under piecewise $C^1$ regularity with
finite one-sided derivative bounds near the point.
\end{theorem}
\begin{proof}
By evenness and normalization of $D_N$,
\[
 S_N(f)(x)-\frac{L_-+L_+}{2}
 =\frac1{2\pi}\int_0^\pi
 \frac{f(x+t)+f(x-t)-L_+-L_-}{\sin(t/2)}
          \sin((N+1/2)t)\dd t.
\]
The quotient before the sine is locally Riemann integrable for $t>0$.
Near zero, $t/\sin(t/2)$ is bounded and tends to $2$, so the assumed
absolute improper integrability proves absolute integrability of that
quotient. Away from zero it is properly integrable. Apply the improper
extension of the preceding Riemann--Lebesgue lemma with frequency
$N+1/2$. If $f$ has one-sided Lipschitz estimates, the numerator is $O(t)$,
and its quotient by $t$ is bounded near zero. Piecewise $C^1$ functions
with finite one-sided derivative bounds have these estimates by the norm
mean-value inequality. Finite values assigned at the jump do not enter
the coefficient integrals or the limits.
\end{proof}

\KEntry{K03.19.jordan}{Dirichlet--Jordan convergence for BV functions}
\begin{theorem}
If a periodic scalar function has bounded variation on a period, then at
every $x$,
$S_N(f)(x)\to(f(x-)+f(x+))/2$.
\end{theorem}
\begin{proof}
We supply the oscillatory estimate needed when Dini's absolute condition
is unavailable. The integrals of $D_N$ over arbitrary subintervals of
$[0,\pi]$ have a common absolute bound $C$ independent of $N$.
Indeed write
\[
 \frac1{\sin(t/2)}=\frac2t+r(t),
\]
where $r$ extends continuously to $0$ by the Taylor estimate for sine.
The $r$ contribution is bounded by $\int_0^\pi|r|$. The other contribution
is $2\int_{\nu u}^{\nu v}\sin y/y\dd y$, $\nu=N+1/2$, whose absolute
value is uniformly bounded. To see the latter directly, the integral on
$[0,1]$ is bounded by $1$, and integration by parts on $[1,T]$ bounds
it by the sum of the endpoint terms and $\int_1^\infty y^{-2}\dd y$.
An interval integral is the difference of two such bounded primitives.

Consider first a real BV function $h$ on $[0,\delta]$ with $h(0)=0$ and
$h(t)\to0$ as $t\downarrow0$. Jordan decomposition writes it as a
difference of nondecreasing functions. Subtract their right limits at zero
(which have equal difference zero), to obtain
$h(t)=u(t)-v(t)$ for $t>0$, where $u,v$ are nonnegative nondecreasing,
vanish at zero and tend to zero there. Given $\eps>0$, take $\delta$ so
small that $u(\delta),v(\delta)<\eps$.
The second mean-value theorem \KRef{K03.09.second_mean_value}, reflected to a nondecreasing
weight, gives
\[
 \int_0^\delta u(t)D_N(t)\dd t
       =u(\delta)\int_\xi^\delta D_N(t)\dd t
\]
for some $\xi\in[0,\delta]$, and the same statement for $v$.
Its hypotheses hold because the monotone weights are integrable and
$D_N$ is continuous, including its removable value at zero. Thus
$|\int_0^\delta hD_N|\leq2C\eps$ uniformly in $N$.

Apply this to $f(x+t)-f(x+)$ and $f(x-t)-f(x-)$, assigning their limiting
value zero at $t=0$. They are BV and have the required limits. On
$[\delta,\pi]$ divide these bounded integrable functions by $\sin(t/2)$
and apply the ordinary Riemann--Lebesgue lemma; the resulting integrals
tend to zero. The even convolution identity, with the two constants
subtracted as in the previous proof, now bounds the limsup of the error
by a constant times $\eps$. Send $\eps$ to zero. For complex $f$ apply
the real proof to its two BV components; there is no complex order or
complex Jordan decomposition being assumed.
\end{proof}

\KEntry{K03.19.coefficient_decay}{Differentiation and Fourier coefficients}
\begin{theorem}
Suppose $f\in C^k([-\pi,\pi],\C)$, $k\geq1$, and
$f^{(j)}(-\pi)=f^{(j)}(\pi)$ for $0\leq j<k$. Then
\[
 c_n(f^{(k)})=(in)^k c_n(f),\qquad
 |c_n(f)|\leq\frac1{2\pi|n|^k}\int_{-\pi}^{\pi}|f^{(k)}(x)|\dd x
 \quad(n\ne0).
\]
In fact $c_n(f)=o(|n|^{-k})$.
\end{theorem}
\begin{proof}
For one derivative integrate $f'e^{-inx}$ by parts. The boundary term
$[f(x)e^{-inx}]_{-\pi}^{\pi}$ is zero, because the values of $f$ match
and $e^{in\pi}=e^{-in\pi}=(-1)^n$. Hence $c_n(f')=in c_n(f)$.
Iterate, using the stated lower-derivative endpoint equalities at each
step. The norm integral bound follows immediately; Riemann--Lebesgue
applied to the properly integrable $f^{(k)}$ sharpens the numerator to
$o(1)$. Endpoint matching is exactly what removes the boundary corrections.
\end{proof}

\KEntry{K03.19.c1_uniform}{Absolute convergence under first-derivative control}
\begin{theorem}
For $f\in C^1([-\pi,\pi],\C)$ with $f(-\pi)=f(\pi)$,
the Fourier series converges absolutely and uniformly to the continuous
periodic extension of $f$.
\end{theorem}
\begin{proof}
For $n\ne0$, $|c_n(f)|=|c_n(f')|/|n|$. Finite Cauchy--Schwarz gives
\[
 \sum_{1\leq|n|\leq N}|c_n(f)|
 \leq\left(\sum_{1\leq|n|\leq N}|c_n(f')|^2\right)^{1/2}
      \left(\sum_{n\ne0}\frac1{n^2}\right)^{1/2}.
\]
The first factor is bounded by $\|f'\|_2$ by Bessel and the second is
finite by the $p$-series test. Increasing partial sums therefore show
$\sum_{n\in\Z}|c_n(f)|<\infty$. The M-test gives a continuous periodic
uniform sum $g$. Termwise proper integration proves $c_n(g)=c_n(f)$
for every $n$. The coefficient uniqueness theorem proved below via Fejer
means therefore gives $g=f$. That theorem does not use the current
$C^1$ statement, so the forward reference is noncircular.

The same proof works for continuous periodic $f$ that is $C^1$ on a finite
partition with finite one-sided derivative limits: integrate by parts on
each piece, where internal boundary values of $f$ cancel. The resulting
piecewise derivative is properly Riemann integrable, and the Bessel
argument is unchanged.
\end{proof}

\KEntry{K03.19.fejer_kernel}{Fejer means and a positive kernel}
\begin{construction}
Put
\[
 \sigma_N(f)=\frac1{N+1}\sum_{k=0}^N S_k(f),\qquad
 F_N(t)=\frac1{N+1}\sum_{k=0}^ND_k(t).
\]
Then
\[
 F_N(t)=\frac1{N+1}\left|\sum_{j=0}^Ne^{ijt}\right|^2
       =\frac{\sin^2((N+1)t/2)}{(N+1)\sin^2(t/2)}
\]
away from $2\pi\Z$, with value $N+1$ there. It is nonnegative and has
integral mean $1$. Its normalized mass outside any fixed
$(-\delta,\delta)\subset(-\pi,\pi)$ tends to zero.
\end{construction}
\begin{proof}
Expand the squared finite sum. The coefficient of $e^{int}$ is
$(N+1-|n|)/(N+1)$ for $|n|\leq N$, exactly its coefficient in the average
of the $D_k$. The geometric-sum identity gives the sine quotient.
Orthogonality computes the integral mean as its constant coefficient $1$.
For $\delta\leq|t|\leq\pi$,
\[
 0\leq F_N(t)\leq\frac1{(N+1)\sin^2(\delta/2)},
\]
so even its normalized integral over that set has this upper bound.
Finite averaging of the convolution identities yields
$\sigma_N(f)(x)=(2\pi)^{-1}\int f(x-t)F_N(t)\dd t$.
\end{proof}

\KEntry{K03.19.fejer}{Uniform Fejer convergence}
\begin{theorem}
For continuous $2\pi$-periodic scalar $f$, $\sigma_N(f)\to f$ uniformly
on $\R$. In particular trigonometric polynomials are uniformly dense in
continuous periodic functions.
\end{theorem}
\begin{proof}
Such an $f$ is bounded, say $|f|\leq M$, and uniformly continuous on
$\R$: reduce one point by a period into $[-\pi,\pi]$, put the other
nearby point in $[-\pi-1,\pi+1]$, and use uniform continuity on this
compact interval. Choose $\delta$ so $|f(x-t)-f(x)|<\eps$ for
$|t|<\delta$, for all $x$. Positivity and normalization of the kernel give
\[
 |\sigma_N(f)(x)-f(x)|
 \leq\eps+\frac{2M}{(N+1)\sin^2(\delta/2)}.
\]
First let $N\to\infty$ and then $\eps\downarrow0$. Every $\sigma_N$ is a
trigonometric polynomial by its finite definition. This proves density
and makes no assertion about uniform convergence of the raw sums $S_N$.
\end{proof}

\KEntry{K03.19.uniqueness}{Uniqueness of coefficients for continuous functions}
\begin{corollary}
A continuous periodic function whose Fourier coefficients all vanish is zero.
Two continuous periodic functions with the same coefficients are equal.
\end{corollary}
\begin{proof}
All $S_N$, hence all $\sigma_N$, of the zero-coefficient function are
zero. Uniform Fejer convergence identifies their limit with the function.
Apply this to the difference of two functions for the second assertion.
The continuity assumption matters when equality is pointwise: changing
a single value of an integrable periodic function changes no coefficient.
\end{proof}

\KEntry{K03.19.parseval}{Parseval for continuous periodic functions}
\begin{theorem}
For continuous periodic $f$,
\[
 \sum_{n\in\Z}|c_n(f)|^2=\frac1{2\pi}\int_{-\pi}^{\pi}|f(x)|^2\dd x.
\]
\end{theorem}
\begin{proof}
The Fejer polynomial $\sigma_N$ lies in the span of $e_n$ for $|n|\leq N$.
Finite best approximation therefore gives
$\|f-S_N(f)\|_2^2\leq\|f-\sigma_N(f)\|_2^2$.
The right side is at most the square of its uniform norm and tends to zero.
The finite orthogonal identity gives
\[
 \|f\|_2^2=\sum_{|n|\leq N}|c_n(f)|^2+\|f-S_N(f)\|_2^2.
\]
Pass to the limit of the nonnegative coefficient sums. This proves the
formula in the stated Riemann/continuous setting, without assuming a
completed Hilbert space or Plancherel theorem.
\end{proof}

\KEntry{K03.19.examples}{Three Fourier examples and justified numerical sums}
\begin{example}
For the periodic sawtooth $f(x)=x$ on $(-\pi,\pi)$, assigned value $0$
at the seams, integration by parts gives
\[
 f(x)\sim2\sum_{n=1}^\infty\frac{(-1)^{n+1}}n\sin nx.
\]
For the square wave $g(x)=\operatorname{sgn}(\sin x)$, with value zero at
its jumps, integration over $(0,\pi)$ gives
\[
 g(x)\sim\frac4\pi\sum_{j=0}^\infty\frac{\sin((2j+1)x)}{2j+1}.
\]
Both are BV on a period. Dirichlet--Jordan proves equality with their
ordinary values at continuity points and with the midpoint value $0$ at
the jumps. At $x=\pi/2$ this gives the ordinary convergent sum
$\sum_{j\geq0}(-1)^j/(2j+1)=\pi/4$.
\end{example}
\begin{proof}[Coefficient computations and a polynomial example]
Oddness kills the cosine coefficients of both odd examples. For the
sawtooth,
$\int_0^\pi x\sin nx\dd x=-\pi(-1)^n/n$; multiplying by $2/\pi$ gives
the stated $b_n$. For the square wave the same factor multiplies
$\int_0^\pi\sin nx\dd x=(1-(-1)^n)/n$, giving zero for even $n$ and
$4/(\pi n)$ for odd $n$.

The periodic extension of $h(x)=x^2$ on $[-\pi,\pi]$ is continuous
and piecewise $C^1$. Direct integration gives $a_0/2=\pi^2/3$, $b_n=0$.
Twice integrating $x^2\cos nx$ by parts gives
$a_n=4(-1)^n/n^2$. Thus
\[
 x^2=\frac{\pi^2}3+4\sum_{n=1}^\infty\frac{(-1)^n}{n^2}\cos nx
 \quad(-\pi\leq x\leq\pi).
\]
The coefficients are absolutely summable, so the sum is uniform; its
identification with $h$ also follows from the already proved BV theorem.
At $x=\pi$ this yields $\sum_{n\geq1}n^{-2}=\pi^2/6$.
Parseval is applicable to this continuous $h$, and gives
\[
 \frac{\pi^4}{5}=\frac{\pi^4}{9}
                 +8\sum_{n=1}^\infty\frac1{n^4},
\]
hence $\sum n^{-4}=\pi^4/90$. We have not applied the
continuous-function Parseval theorem to the discontinuous square wave.
\end{proof}

\KEntry{K03.19.normalization_bridge}{Periods and real/complex representations}
\begin{proposition}
For period $T>0$, use
$\widehat f_T(n)=T^{-1}\int_{a}^{a+T}f(x)e^{-2\pi inx/T}\dd x$.
If $g(t)=f(Tt/(2\pi))$, then $c_n(g)=\widehat f_T(n)$ when both
origins are zero; translating the origin changes the expected phase.
\end{proposition}
\begin{proof}
The substitution $x=Tt/(2\pi)$ changes the normalized integral on
$[-\pi,\pi]$ into $T^{-1}\int_{-T/2}^{T/2}$ of the stated integrand.
It is $T$-periodic because $n$ is an integer, so the integral may be moved
to any interval $[a,a+T]$ without change. For $f_a(x)=f(x+a)$ a direct
substitution gives $\widehat{f_a}_T(n)=e^{2\pi ina/T}\widehat f_T(n)$.
Expanding $e^{-inx}=\cos nx-i\sin nx$ proves the real/complex
coefficient conversions stated at the start; combining the $n$ and $-n$
terms proves equality of the corresponding finite partial sums. Convergence
claims transfer only through those proved identities and their stated limits.
\end{proof}

\KEntry{K03.19.fourier_boundary}{The boundary with general Fourier analysis}
\begin{remark}[Routing record]
General Fourier transforms, Plancherel, arbitrary $L^p$ convergence,
distributions and quantitative Gibbs theory belong to D11/D24.
The present finite-interval results do not assert those theories.
In particular continuity alone is not offered as a theorem guaranteeing
pointwise convergence of raw Fourier partial sums. Fejer convergence,
absolute convergence under derivative control and the explicit Dini/BV
criteria remain distinct conclusions with distinct hypotheses.
\end{remark}
\appendix

% BEGIN REVISED SOURCE source-register.tex SHA256=9f3d629f3d7f7b961b3da208b2d6daf79c9ec8065099db05ca34d45860b2f045
\chapter{Source register and scope correspondence}

The following are comparison sources, not a license to replace missing proofs by book titles. Physical PDF pages are one-based. The source hashes, acquisition provenance and section-to-entry map are retained in the companion electronic register. Existing textbooks are unchanged.

\section*{Comparison volumes}

\paragraph{A01} Vladimir A. Zorich, \emph{Mathematical Analysis I}. Second English edition, copyright 2015.

\paragraph{A02} Vladimir A. Zorich, \emph{Mathematical Analysis II}. Second English edition, copyright 2016.

\paragraph{A03} Walter Rudin, \emph{Principles of Mathematical Analysis}. Third edition, 1976, international printing. Readable scan with OCR; check formulas against page images.

\paragraph{D03-S01} Tom M. Apostol, \emph{Mathematical Analysis}. Second edition. Readable scan with OCR; first PDF page is the cover, followed by the preface and contents. Separate title/copyright leaves are not present in this copy. Body starts at printed p. 1 (PDF p. 14), and final index reaches printed p. 492 (PDF p. 505).

\paragraph{D03-S02} Charles Chapman Pugh, \emph{Real Mathematical Analysis}. Second edition, copyright 2015; copyright page states corrected at second printing 2017. eBook ISBN 978-3-319-17771-7. Body begins at printed p. 1 (PDF p. 12); final index reaches printed p. 478 (PDF p. 486). Printed-to-PDF offsets change later; use verified per-location pages, not one global offset.

\paragraph{F03} Terence Tao, \emph{Analysis I}. Fourth edition, copyright 2022; Springer electronic edition 2023.

\paragraph{A37} Frank W. J. Olver, \emph{Asymptotics and Special Functions}. Original copyright 1974, later Academic Press printing. Original chapter archive: all 14 mathematical chapters (printed pages 1-544), answers, references and both indexes (through printed page 572). 23 PDFs total 578 PDF pages including selected frontmatter. No separate contents PDF; not a complete facsimile of every frontmatter page.

\section*{Located source sections}

\small

\begin{longtable}{p{22mm}p{99mm}p{20mm}}

\toprule Source key & Section or passage & PDF page \\\midrule\endhead

A01:3.1 & 3.1 The Limit of a Sequence & 100 \\

A01:3.2 & 3.2 The Limit of a Function & 127 \\

A01:4 & Chapter 4: Continuous Functions & 170 \\

A01:5.1 & 5.1 Differentiable Functions & 192 \\

A01:5.2 & 5.2 The Basic Rules of Differentiation & 211 \\

A01:5.3 & 5.3 The Basic Theorems of Differential Calculus & 232 \\

A01:5.4 & 5.4 The Study of Functions Using the Methods of Differential Calculus & 255 \\

A01:5.3.3 & 5.3.3 Taylor's Formula & 238 \\

A01:5.5 & 5.5 Complex Numbers and the Connections Among the Elementary Functions & 284 \\

A01:5.6.5 & 5.6.5 The Number e and the Function expx Revisited & 316 \\

A01:6.1 & 6.1 Definition of the Integral and Description of the Set of Integrable Functions & 351 \\

A01:6.3 & 6.3 The Integral and the Derivative & 380 \\

A01:6.3.3 & 6.3.3 Integration by Parts in the Definite Integral and Taylor's Formula & 384 \\

A01:6.5 & 6.5 Improper Integrals & 413 \\

A01:D & Appendix D: The Euler-MacLaurin Formula & 593 \\

A01:E & Appendix E: Riemann-Stieltjes Integral, Delta Function, and the Concept of Generalized Functions & 598 \\

A01:5.7.4 & 5.7.4 Primitives of the Form R(cosx, sinx)dx & 339 \\

A01:5.7.5 & 5.7.5 Primitives of the Form R(x,y(x))dx & 342 \\

A02:16 & Chapter 16: Uniform Convergence and the Basic Operations of Analysis on Series and Families of Functions & 380 \\

A02:16.4 & 16.4 *Compact and Dense Subsets of the Space of Continuous Functions & 412 \\

A02:17.1 & 17.1 Proper Integrals Depending on a Parameter & 421 \\

A02:17.2 & 17.2 Improper Integrals Depending on a Parameter & 429 \\

A02:17.3 & 17.3 The Eulerian Integrals & 449 \\

A02:18.1 & 18.1 Basic General Concepts Connected with Fourier Series & 509 \\

A02:18.2 & 18.2 Trigonometric Fourier Series & 536 \\

A02:19.1 & 19.1 Asymptotic Formulas and Asymptotic Series & 605 \\

A02:dini & Dini's theorem, Proposition 2 in section 16.3 & 401 \\

AP:7.3-exercise & Exercise 7.3, alternative mesh definition & 187 \\

A03:3 & Chapter 3 Numerical Sequences and Series & 58 \\

A03:4 & Chapter 4 Continuity & 94 \\

A03:6 & Chapter 6 The Riemann-Stieltjes Integral & 131 \\

A03:7 & Chapter 7 Sequences and Series of Functions & 154 \\

A03:8 & Chapter 8 Some Special Functions & 183 \\

A03:gamma & The Gamma Function & 203 \\

A03:fourier & Fourier Series & 196 \\

A03:mvt & Mean Value Theorems & 118 \\

A03:nowhere & Theorem 7.18: periodic-tent nowhere-differentiable construction & 165 \\

AP:5 & Chapter 5 Derivatives & 117 \\

AP:5.10 & 5.10 The Mean-Value Theorem for derivatives & 123 \\

AP:5.12 & 5.12 Taylor's formula with remainder & 126 \\

AP:7 & Chapter 7 The Riemann-Stieltjes Integral & 153 \\

AP:7.20 & 7.20 Second fundamental theorem of integral calculus & 175 \\

AP:7.23 & 7.23 Riemann-Stieltjes integrals depending on a parameter & 179 \\

AP:7.26 & 7.26 Lebesgue's criterion for existence of Riemann integrals & 182 \\

AP:8 & Chapter 8 Infinite Series and Infinite Products & 196 \\

AP:9 & Chapter 9 Sequences of Functions & 231 \\

AP:9.14 & 9.14 Power series & 247 \\

AP:9.19 & 9.19 The Taylor's series generated by a function & 254 \\

AP:10.13 & 10.13 Improper Riemann integrals & 289 \\

AP:11 & Chapter 11 Fourier Series and Fourier Integrals & 319 \\

AP:6 & Chapter 6 Functions of Bounded Variation and Rectifiable Curves & 140 \\

PUGH:2.2 & 2 Continuity & 71 \\

PUGH:3.1 & 1 Differentiation & 159 \\

PUGH:4.3 & 3 Compactness and Equicontinuity in C0 & 233 \\

PUGH:4.4 & 4 Uniform Approximation in C0 & 238 \\

F03:5.6 & 5.6 Real Exponentiation, Part I & 119 \\

F03:6 & 6 Limits of Sequences & 123 \\

F03:7 & 7 Series & 148 \\

F03:9 & 9 Continuous Functions on R & 196 \\

F03:11 & 11 The Riemann Integral & 243 \\

F03:11.9 & 11.9 The Two Fundamental Theorems of Calculus & 267 \\

A37:1 & Chapter 1 (chapter PDF within the Olver archive) & 1 \\

A37:2 & Chapter 2 (chapter PDF within the Olver archive) & 1 \\

A37:3 & Chapter 3 (chapter PDF within the Olver archive) & 1 \\

A37:8 & Chapter 8 (chapter PDF within the Olver archive) & 1 \\

\bottomrule\end{longtable}\normalsize

The accompanying chapter-routing matrix accounts for all 77 main chapters of the seven source units, including chapters shared with D01/D02 and chapters assigned to other directions. It does not claim to include every exercise or every theorem of all seven sources. The working D03 notes retain the fourteen core module identities and five bounded supplements.

The original sources sometimes use distinct definitions or compressed illustrative assertions. The notes use explicit scalar hypotheses for compactness, keep the refinement/mesh distinction for Stieltjes integration, and supply the one-variable special-integral dependencies. These are substantive source-comparison decisions, recorded in the mathematical audits rather than hidden in typography.

% BEGIN REVISED SOURCE verification-register.tex SHA256=422b45266c4a18faae27569c036602a4f07d5a7df1b98ce85d81b3607a71cdd3
\chapter{Coverage, verification and later reuse}
The companion coverage manifest links every stable knowledge identifier to its
chapter and distinguishes candidate mathematical content from the five routing
records. Definitions are checked for compatible conventions; theorem proofs,
construction arguments and counterexamples are reviewed as their respective
kinds of mathematical content. A routing paragraph is not counted as a proof
of the theory to which it points.

The source-level comparison, author self-checks and cross-author mathematical
reviews are saved separately. Their reports identify actual issues, repairs and
the versions reread. The technical build checks unique identifiers, resolved
references and a readable typeset document. These are different checks: a
successful LaTeX build is not evidence that a mathematical argument is valid,
and a mathematical review is not a Lean kernel check.

The existing organization snapshot is retained. Where a manuscript clarifies a
statement, proves a stronger useful variant or chooses a more explicit
construction, the coverage and review records explain the relation rather than
silently treating the old syllabus as an already verified proof. The number of
knowledge entries does not prescribe the number or boundaries of future Lean
goals.

Before formalization, compare exact assumptions and conclusions with the fixed
versions of Mathlib and MathGenome. Reuse matching definitions and results, and
prove required representation bridges. In particular, Riemann, oriented
interval, Stieltjes and measure-integral APIs need more than matching names;
real-variable derivatives of complex-valued functions are distinct from complex
derivatives. The existing libraries and their acceptance evidence are not
modified by this manuscript.

This document contains no new Lean axioms, no claimed accepted Lean declarations,
and no claim to have formalized every theorem in its comparison textbooks.
Formal acceptance and library release are later work. The present delivery is
an English mathematical manuscript with source correspondence, proofs,
counterexamples and documented mathematical review.
\end{document}
